\chapter{The construction}

\section{Necessary Conditions for Pattern Occurrence}

\textbf{Purpose:} Separate what an anchor proves from what it costs, letting a reader tell which claims here need no data and which ones are only as good as the measurements under them. \textbf{Scope:} \texttt{test/\allowbreak{}bench/\allowbreak{}bench\_\allowbreak{}ancorae\_\allowbreak{}sift.\allowbreak{}c}, \texttt{test/\allowbreak{}bench/\allowbreak{}bench\_\allowbreak{}ancorae\_\allowbreak{}lattice.\allowbreak{}c}, \texttt{test/\allowbreak{}bench/\allowbreak{}bench\_\allowbreak{}ancorae\_\allowbreak{}entropy.\allowbreak{}c}, \texttt{src/\allowbreak{}impensa\_\allowbreak{}ancorae\_\allowbreak{}acus/\allowbreak{}}, \texttt{test/\allowbreak{}support/\allowbreak{}mmgr\_\allowbreak{}sha256.\allowbreak{}h}, \texttt{test/\allowbreak{}support/\allowbreak{}mmgr\_\allowbreak{}sha256.\allowbreak{}c}

\subsection{Abstract}

An anchor is a condition copied out of a search pattern and tested at one position. This document separates two claims about that construction that are usually stated together. The first is that no choice of anchors can lose a true occurrence, which follows from the definition and holds for any selection rule, any number of anchors, any alphabet, and any index set. The second is how many false candidates survive, which is a property of a domain and is measured.

Three benches carry the measurements. Over byte strings, 9,396,207 true occurrences were examined under three anchor selection measures including one carrying no information, at one to six anchors, on five corpora, at needle lengths from 1 to 2048, with none refused. Over thirteen other geometries, including a two dimensional point set, that set rotated a quarter turn, a three dimensional volume, a domain whose symbols are complex numbers with irrational parts, and hypercubes of dimension one through eight, a further 213,840 occurrences were examined with none refused. Cost moved by up to a factor of 6.4 between measures. What a measure can be worth is given a closed form here, the uninformed rate over the expected minimum of \(m\) size biased draws, which equals one exactly when the distribution is uniform on its support. Entropy does not govern that quantity: one corpus here carries more entropy than another and has less than half its ceiling.

A refutation is shown to settle more than the position it was tested at. One read of one cell settles every alignment whose pattern symbol differs from what was found, which is \(m(1-2^{-H_2})\) of them in expectation, confirmed to within 0.7\% over 84 rows spanning three selection measures. The collision probability that sets the candidate rate therefore also sets the refutation distance, appearing once as a rate and once as a length.

The candidate count is shown to be an estimator of Renyi entropy of order two, equal in expectation to the plug in estimator when the probe's own match is counted and to the standard unbiased estimator when it is not. Scored against published estimators on sources with known parameters, it is worse in both accuracy and running time at every setting tested, converging toward the unbiased estimator as the probe budget approaches the corpus length. Its only advantage is the constraint it survives: it builds no histogram and never enumerates the alphabet.

The correction that separates the two forms is the exclusion of a coincident index, and applying it recursively gives the unbiased estimators at higher Renyi orders. Its magnitude at order two has the closed form \((1-C)/((N-1)C\ln 2)\) bits for a plug in collision probability \(C\), confirmed across nine sources spanning four orders of magnitude. The numerator vanishes exactly at a point mass. On a source carrying no entropy the correction is identically zero at every order and every length.

Put against the classical algorithms on one corpus and one counter, anchoring at the rarest symbol loses to a Horspool shift on every corpus and needle length by 1.2 to 25 times, because an anchor at offset \(a\) can advance by at most \(a+1\), and advance is what governs search cost. Candidate rate does not. Treating the two as factors of one product recovers the loss and wins 11\% on C source. Replacing the frequency model with a running advance per offset, which needs no alphabet, wins 13.2\% and beats a rule handed the corpus histogram outright.

Two negative results transfer beyond this construction. Sizing a cascade of anchors by multiplying their individual rates underestimates the survivors by up to four orders of magnitude on structured data, always in the direction that under-provisions. And selecting by the center of mass of a utility loses to selecting by its maximum on every row measured, by up to 42\%.

\subsection{1. Introduction}

A search for a pattern in a body of data can test every position, or it can find a cheap condition that most positions fail and test only the survivors. The second is the older idea and the cheap condition here is an anchor: one of the pattern's own symbols, checked at one offset.

Two questions about that arrangement get answered together in most treatments and come apart under examination. Whether the filter is safe is a question about the construction. How much the filter saves is a question about the data. This document keeps them apart, because the first needs no experiment and the second cannot be settled without one.

Section 2 states the construction and proves what follows from it. Section 3 describes what was measured and how. Section 4 reports the measurements. Section 5 discusses what the result is and what transfers. Section 6 states the limits as numbers. Sections 7 and 8 cover the prior art that was read and what remains open.

The digest implementation in \texttt{test/\allowbreak{}support} appears here as an instrument. It generates the uniform control corpus and serves as the test oracle. The strength of that control is the strength of its vectors. Appendix A covers it.

\subsection{2. Mathematical Framework}

\subsubsection{2.1 Ambient Algebra and Pattern Support}

Let \(\Sigma\) be the carrier of symbols. No structure on \(\Sigma\) is assumed: no order, no metric, no finite cardinality, and no way to enumerate its members. The only operation available is a decidable equality between two positions.

Let \((G,+)\) be an additive group of displacements and let a domain be a map \(x\) from an admissible subset of \(G\) to \(\Sigma\). A pattern is a finite displacement configuration with support \(D \subset G\), together with the symbols a chosen base position carries at those displacements. For a base \(t\) and a candidate base \(s\), define

\(\mathrm{Occ}(s) \iff \forall d \in D:\; x(s+d) = x(t+d)\)

An anchor choice is any subconfiguration \(A \subseteq D\), and

\(\mathrm{Anc}_A(s) \iff \forall a \in A:\; x(s+a) = x(t+a)\)

The null used here is the \(\Sigma\)-delta null, and its operational form is the permutation null already used throughout this work. Semantically, \(\Sigma\) projects its internal structure onto an induced topology \(\tau_\Sigma\); the instrument then examines that structured carrier against its own constrained reference. Let \(\Pi_\Sigma\) be the family of permutations that preserve the observed carrier counts and stated partition. The reference is sampled as \(x^\circ_\Sigma = \pi(x)\) for \(\pi \in \Pi_\Sigma\), and for a statistic \(M\) the measured residual is

\(\Delta_\Sigma[M] = M(x) - \mathbb{E}_{\pi \sim \Pi_\Sigma}\!\left[M\bigl(\pi(x)\bigr)\right]\)

This is not a new operator or code path: it is the permutation-null measure already implemented. The residual is the object of study. The measure does not claim that a hidden generator has been reconstructed, and it does not invoke Laplace's demon. A predictive or reconstructive engine that attempts to infer the complete state from its laws would be a separate measurement. Here \(\Sigma\) supplies both the topology and the reference, and the instrument examines the engine's residue in the object itself.

Nothing in either definition requires \(G\) to be the integers, \(D\) to be contiguous, or \(\Sigma\) to be known.

\subsubsection{2.2 Soundness of Subconfiguration Filtering}

For every \(A \subseteq D\), every domain, and every base,

\(\mathrm{Occ}(s) \;\Longrightarrow\; \mathrm{Anc}_A(s)\)

The proof is that \(A\) is a subconfiguration of \(D\), and a conjunction over \(D\) contains the conjunction over \(A\). The search runs the contrapositive:

\(\neg\,\mathrm{Anc}_A(s) \;\Longrightarrow\; \neg\,\mathrm{Occ}(s)\)

A position an anchor rejects is settled exactly and permanently, with no verification, and no later stage revisits it.

\subsubsection{2.3 Failure of Completeness}

For every proper subconfiguration \(A \subsetneq D\), the converse

\(\mathrm{Anc}_A(s) \;\Longrightarrow\; \mathrm{Occ}(s)\)

is not entailed, in any domain and in any state. A witness is any \(s\) agreeing with the base on \(A\) and differing on some \(d \in D \setminus A\), and the definitions permit such an \(s\) whenever \(|\Sigma| > 1\). The only \(A\) for which the converse holds is \(A = D\), where the filter has become the verification.

A matching anchor establishes nothing. The exact compare cannot be removed by any amount of knowledge about the domain, and adding anchors reduces the survivors without ever reaching certainty.

\subsubsection{2.4 Structural Corollaries}

\textbf{The selection rule is free.} Proposition 1 quantifies over every subconfiguration and never mentions what chose one. No cost table, no random table, and no absence of a table can produce a false negative. What a measure moves is how often the rejection fires, which is cost.

\textbf{The alphabet need not be known.} Both predicates are equalities between positions. Where \(\Sigma\) cannot be enumerated no frequency table exists, and the flat measure is what remains. Maximum entropy is therefore the base case for this construction, and a weighted table is the special case a known domain affords.

\textbf{The displacement support need not be ordered.} The definitions use \(+\) and never \(<\). A grid, a volume, and a rotated point configuration differ only in which members of \(G\) appear in \(D\). A rotation is a permutation of \(D\), which changes no term in either predicate.

\subsubsection{2.5 Probabilistic Cost Model}

Model occurrences of an anchor symbol as a Bernoulli process at rate \(q\). Over \(N\) candidate positions,

\(\mathbb{E}[C_1] = N q, \qquad \mathbb{E}[S] = \frac{1}{q}\)

\(\mathbb{E}[C_1]\) is what the verify step pays for and \(\mathbb{E}[S]\) is the run a skipping search never touches. Reducing \(N\) candidates to \(O(1)\) survivors requires \(\log_2 N\) bits of discrimination and one anchor supplies \(-\log_2 q\) bits.

The second expression is Kac's lemma. For a stationary ergodic source the expected return time to a set of measure \(q\) is exactly \(1/q\), which makes the skip a recurrence time and not an artifact of the filter. That names an assumption the construction does not supply. Proposition 1 guarantees that no occurrence is lost, and Section 2.5 prices the search, and neither says a pattern occurs at all. That last guarantee belongs to the source: a stationary ergodic process visits any event of positive probability almost surely, and a process that is not stationary or not ergodic carries no such promise. Three guarantees with three owners, and only the first is this document's.

For an uninformed anchor, the symbol is drawn from the domain by the domain's own frequencies, so

\[
q \;=\; \sum_{\sigma \in \Sigma} p(\sigma)^2 \;=\; 2^{-H_2}, \qquad
H_2 \;=\; -\log_2 \sum_{\sigma} p(\sigma)^2
\]

with \(H_2\) the Renyi entropy of order two. This predicts the uninformed candidate count from the domain alone, before any pattern is chosen, and Section 4.4 tests it.

\subsubsection{2.6 Filtering Pipeline}

\begin{lstlisting}[breaklines=true,breakatwhitespace=false,basicstyle=\ttfamily\small]
flowchart LR
    P["position s"] --> A{"anchor a in A<br/>x(s+a) == x(t+a) ?"}
    A -->|"fails"| R["REFUTED<br/>exact, final<br/>never revisited"]
    A -->|"holds"| C["candidate<br/>Proposition 2:<br/>proves nothing"]
    C --> V{"exact compare<br/>every d in D"}
    V -->|"agrees"| O["occurrence"]
    V -->|"differs"| F["false positive"]
    LIC["(c) 2026 Douglas Quigg (dstroy0)<br/>AGPL-3.0-or-later OR LicenseRef-Commercial"]
    style R fill:#1f6f43,color:#fff
    style V fill:#7a4a1f,color:#fff
    style LIC fill:#eeeeee,color:#333333,stroke-dasharray: 3 3
\end{lstlisting}

\textbf{Figure 1.} The anchor filter, as \texttt{bench\_\allowbreak{}ancorae\_\allowbreak{}lattice.\allowbreak{}c:193-\allowbreak{}239} implements it. The green box is the only place the method is certain. Everything to the right of the candidate is cost.

\subsection{3. Experimental Design}

\subsubsection{3.1 Benchmarks and Estimands}

\tiny
\setlength{\tabcolsep}{2pt}
\begin{longtable}{@{}>{\setlength{\rightskip}{0pt}\setlength{\parfillskip}{0pt}\arraybackslash}p{\dimexpr\textwidth/3-2\tabcolsep\relax}>{\setlength{\rightskip}{0pt}\setlength{\parfillskip}{0pt}\arraybackslash}p{\dimexpr\textwidth/3-2\tabcolsep\relax}>{\setlength{\rightskip}{0pt}\setlength{\parfillskip}{0pt}\arraybackslash}p{\dimexpr\textwidth/3-2\tabcolsep\relax}@{}}
\toprule
bench & question & reports \\
\midrule
\texttt{bench\_\allowbreak{}ancorae\_\allowbreak{}sift.\allowbreak{}c} & cost and independence over byte strings & candidates, skip, \(I(d)\), \(z\), refusals \\
\texttt{bench\_\allowbreak{}ancorae\_\allowbreak{}lattice.\allowbreak{}c} & whether Proposition 1 survives off the line & candidates, refusals \\
\texttt{bench\_\allowbreak{}ancorae\_\allowbreak{}entropy.\allowbreak{}c} & what the candidate count estimator is worth & bias, root mean square error \\
\texttt{bench\_\allowbreak{}ancorae\_\allowbreak{}ab.\allowbreak{}c} & how a whole search compares to one built another way & corpus symbol accesses \\
\bottomrule
\end{longtable}
\normalsize

All four report counts. Nothing is timed, no row is a performance claim, and every number is a property of the data and the geometry, identical on every part.

\subsubsection{3.2 Correctness and Cost Estimands}

Every row carries a correctness column and a cost column, graded differently. The correctness column counts true occurrences that an anchor rejected, and Proposition 1 says it has one acceptable value. A nonzero entry there would be a defect in the bench, never a property of a domain. The cost column is expected to move with the measure, the anchor count and the geometry.

A verdict of \texttt{none} is printed where a case had nothing to check (\texttt{test/\allowbreak{}bench/\allowbreak{}bench\_\allowbreak{}ancorae\_\allowbreak{}sift.\allowbreak{}c:715-\allowbreak{}757}).

\subsubsection{3.3 Byte-String Corpora}

Four corpora, 2048 bytes each except where stated, chosen to span the range of factor complexity:

\tiny
\setlength{\tabcolsep}{2pt}
\begin{longtable}{@{}>{\setlength{\rightskip}{0pt}\setlength{\parfillskip}{0pt}\arraybackslash}p{\dimexpr\textwidth/3-2\tabcolsep\relax}>{\setlength{\rightskip}{0pt}\setlength{\parfillskip}{0pt}\arraybackslash}p{\dimexpr\textwidth/3-2\tabcolsep\relax}>{\setlength{\rightskip}{0pt}\setlength{\parfillskip}{0pt}\arraybackslash}p{\dimexpr\textwidth/3-2\tabcolsep\relax}@{}}
\toprule
corpus & content & \(H_2\) (bits) \\
\midrule
\texttt{flat} & one byte repeated & 0.000 \\
\texttt{structured} & 1408 bytes of C source & 3.620 \\
\texttt{english} & non-repeating English prose & 3.692 \\
\texttt{periodic16} & 16 byte fixed width records, layout repeating and content not & 4.611 \\
\texttt{uniform} & SHA-256 in counter mode & 7.832 \\
\bottomrule
\end{longtable}
\normalsize

Needles are drawn from the corpus so every one genuinely occurs. A corpus built by repeating a unit is self similar at that scale and every needle then occurs once per repetition by construction, which ruined an earlier measurement recorded in Section 6.5. The fill truncates instead of repeating (\texttt{test/\allowbreak{}bench/\allowbreak{}bench\_\allowbreak{}ancorae\_\allowbreak{}sift.\allowbreak{}c:190-\allowbreak{}193}).

\subsubsection{3.4 Anchor-Selection Policies}

Three, differing as much as the interface allows (\texttt{test/\allowbreak{}bench/\allowbreak{}bench\_\allowbreak{}ancorae\_\allowbreak{}sift.\allowbreak{}c:342-\allowbreak{}347}, \texttt{361-\allowbreak{}420}):

\tiny
\setlength{\tabcolsep}{2pt}
\begin{longtable}{@{}>{\setlength{\rightskip}{0pt}\setlength{\parfillskip}{0pt}\arraybackslash}p{\dimexpr\textwidth/2-2\tabcolsep\relax}>{\setlength{\rightskip}{0pt}\setlength{\parfillskip}{0pt}\arraybackslash}p{\dimexpr\textwidth/2-2\tabcolsep\relax}@{}}
\toprule
policy & how a symbol is priced \\
\midrule
\texttt{table} & the linked cost table, ranking a byte by rarity \\
\texttt{random} & a permutation of 0 through 255 drawn from SHA-256, unrelated to frequency \\
\texttt{maxent} & one cost for every byte, so the table carries no information \\
\bottomrule
\end{longtable}
\normalsize

\texttt{random} is a permutation. Every cost stays distinct and the picker behaves exactly as under the real table. The only thing removed is the table being right. \texttt{maxent} makes the picker take the leftmost offsets and anchors land adjacent. That is the hardest arrangement for two rates to multiply.

\subsubsection{3.5 Geometric Configurations}

Six cases share one core that receives a list of valid base positions, a list of displacements, and a callback answering whether two positions carry the same symbol (\texttt{test/\allowbreak{}bench/\allowbreak{}bench\_\allowbreak{}ancorae\_\allowbreak{}lattice.\allowbreak{}c:80}, \texttt{193-\allowbreak{}239}). The core has no dimension parameter, because the geometry is entirely inside the base list, and no symbol type, because it only asks whether two agree.

\tiny
\setlength{\tabcolsep}{2pt}
\begin{longtable}{@{}>{\setlength{\rightskip}{0pt}\setlength{\parfillskip}{0pt}\arraybackslash}p{\dimexpr\textwidth/3-2\tabcolsep\relax}>{\setlength{\rightskip}{0pt}\setlength{\parfillskip}{0pt}\arraybackslash}p{\dimexpr\textwidth/3-2\tabcolsep\relax}>{\setlength{\rightskip}{0pt}\setlength{\parfillskip}{0pt}\arraybackslash}p{\dimexpr\textwidth/3-2\tabcolsep\relax}@{}}
\toprule
case & \(G\) & \(D\) \\
\midrule
\texttt{line1d} & \(\mathbb{Z}\) & 8 contiguous offsets \\
\texttt{grid2d\_\allowbreak{}box} & \(\mathbb{Z}^2\) & a 2 by 4 rectangle \\
\texttt{grid2d\_\allowbreak{}scatter} & \(\mathbb{Z}^2\) & 8 points in a 5 by 5 window, no row or column filled \\
\texttt{grid2d\_\allowbreak{}turned} & \(\mathbb{Z}^2\) & that scatter under \((r,c) \mapsto (c, 4-r)\) \\
\texttt{cube3d\_\allowbreak{}box} & \(\mathbb{Z}^3\) & a 2 by 2 by 2 block \\
\texttt{field1d\_\allowbreak{}complex} & \(\mathbb{Z}\) & 8 contiguous offsets over complex symbols \\
\texttt{cube1d} to \texttt{cube8d} & \(\mathbb{Z}^d\), \(d = 1 \dots 8\) & a star: the origin, then one step out along each axis in turn \\
\bottomrule
\end{longtable}
\normalsize

\texttt{field1d\_\allowbreak{}complex} holds \texttt{double \_\allowbreak{}Complex} values built from square roots of primes, compared over their storage (\texttt{test/\allowbreak{}bench/\allowbreak{}bench\_\allowbreak{}ancorae\_\allowbreak{}lattice.\allowbreak{}c:479-\allowbreak{}511}). No value is read, ordered, or interpreted anywhere in the file.

Anchor subconfigurations are chosen by three uninformed rules: the leading points, an even stride across them, and a permutation drawn from SHA-256 (\texttt{test/\allowbreak{}bench/\allowbreak{}bench\_\allowbreak{}ancorae\_\allowbreak{}lattice.\allowbreak{}c:291}). No cost table exists in that file and none can, since a domain with an unenumerable alphabet has no frequencies to weigh.

\subsubsection{3.6 Collision-Entropy Estimation}

Sources are written down as distributions before any data exists. The true collision probability is known in advance and no estimator is scored against a histogram of the corpus it was computed from (\texttt{test/\allowbreak{}bench/\allowbreak{}bench\_\allowbreak{}ancorae\_\allowbreak{}entropy.\allowbreak{}c:160}). Eight sources: uniform over 2, 16 and 256 symbols; Zipf at exponents 0.5, 1.0 and 1.5; and two point sources at 0.90 and 0.99. Corpora are drawn by inverse transform sampling at lengths 256, 1024, 4096 and 16384, with 64 independent trials per row.

Four estimators of the same quantity are scored by bias and root mean square error in bits:

\tiny
\setlength{\tabcolsep}{2pt}
\begin{longtable}{@{}>{\setlength{\rightskip}{0pt}\setlength{\parfillskip}{0pt}\arraybackslash}p{\dimexpr\textwidth/3-2\tabcolsep\relax}>{\setlength{\rightskip}{0pt}\setlength{\parfillskip}{0pt}\arraybackslash}p{\dimexpr\textwidth/3-2\tabcolsep\relax}>{\setlength{\rightskip}{0pt}\setlength{\parfillskip}{0pt}\arraybackslash}p{\dimexpr\textwidth/3-2\tabcolsep\relax}@{}}
\toprule
estimator & expression & needs \\
\midrule
plug in & \(\sum_\sigma \hat p(\sigma)^2\) & a histogram \\
unbiased & \(\sum_\sigma n_\sigma(n_\sigma-1) / N(N-1)\) & a histogram \\
probe, self counted & mean over probes of matching positions over \(N\) & neither \\
probe, self excluded & the same, less the probe's own match & neither \\
\bottomrule
\end{longtable}
\normalsize

The unbiased estimator and the first order bias correction of the plug in estimator are the same expression, which makes them one estimator and not two (\texttt{test/\allowbreak{}bench/\allowbreak{}bench\_\allowbreak{}ancorae\_\allowbreak{}entropy.\allowbreak{}c:256-\allowbreak{}344}).

A fifth quantity, the most common value estimate of min entropy from NIST SP 800-90B section 6.3.1, is computed for context and is not scored on the same scale, because min entropy is a different property (\texttt{test/\allowbreak{}bench/\allowbreak{}bench\_\allowbreak{}ancorae\_\allowbreak{}entropy.\allowbreak{}c:379}).

\subsection{4. Results}

\subsubsection{4.0 Overview of Identified Quantities}

The sections below read as a list of separate findings and are not one. They are a single quantity, the collision probability \(2^{-H_2}\), determining every cost in the construction, together with the three conditions under which it stops determining anything. This section says which is which so the rest can be read as one argument.

\textbf{Nothing about correctness is in this at all.} Sections 4.1 and 4.2 report zero refusals over 9,396,207 and 213,840 true occurrences, across three selection measures, thirteen geometries, dimensions one through eight, a rotated point set, and a complex alphabet. Proposition 1 mentions no distribution. No distributional parameter can reach it. Everything else on this page is cost.

\textbf{The cost is one number.} Each of these is a function of \(2^{-H_2}\) and the two sizes \(m\) and \(N\):

\tiny
\setlength{\tabcolsep}{2pt}
\begin{longtable}{@{}>{\setlength{\rightskip}{0pt}\setlength{\parfillskip}{0pt}\arraybackslash}p{\dimexpr\textwidth/3-2\tabcolsep\relax}>{\setlength{\rightskip}{0pt}\setlength{\parfillskip}{0pt}\arraybackslash}p{\dimexpr\textwidth/3-2\tabcolsep\relax}>{\setlength{\rightskip}{0pt}\setlength{\parfillskip}{0pt}\arraybackslash}p{\dimexpr\textwidth/3-2\tabcolsep\relax}@{}}
\toprule
quantity & form & where \\
\midrule
candidate rate per position & \(2^{-H_2}\) & 4.4 \\
alignments settled by one read & \(m(1-2^{-H_2})\) & 4.4.1 \\
estimator bias correction & \((1-C)/((N-1)C\ln 2)\) & 4.5.1 \\
fraction an in-order walk collects & \(1/((1+m2^{-H_2})(1-2^{-H_2}))\) & 4.9.4 \\
where the advance saturates & \(3 \cdot 2^{H_2}\) & 4.9.2 \\
anchors before the space is spent & \(\log_2 N / H_2\) & 4.6.1 \\
free order probe stride & \(m H_2 / \log_2(m H_2 \ln 2)\) & 4.9.5 \\
candidates in \(d\) dimensions & \(1 + (P-1)2^{-nH_2}\) & 4.2 \\
\bottomrule
\end{longtable}
\normalsize

The last row is the one that makes the others worth stating together. That form holds unchanged from a line to an eight dimensional hypercube, over an alphabet of complex numbers with irrational parts, which is where the parameter shows it does not depend on the geometry or on what a symbol is.

\textbf{It stops in three places, and every failure recorded in this document is one of them.}

\_The number is not unique, it is indexed by a slice.\_ There is no \(H_2\) of a corpus, only \(H_2(w)\) of a corpus read \(w\) bits at a time, and Section 4.10 measures it falling from 0.984 to 0.600 bits per bit between \(w = 1\) and \(w = 16\) on English. Every figure above uses \(w = 8\), which nothing here justifies.

\_A second moment is not a distribution.\_ Section 4.3.1 needs the expected minimum of \(m\) size biased draws, which is a property of the whole order statistics. Entropy orders the corpora wrongly there, and Section 4.8 prices a mismatched reference at 25 to 30 times, which no collision probability predicts.

\_Independence is assumed throughout and fails.\_ Section 4.6 measures the product rule wrong by four orders of magnitude, Section 4.7 finds correlation surviving to separation seven, Section 4.6.1 has 2.98 anchors predicted against 6 measured at \(m = 16\), and Section 4.9.3 has the distance field losing on periodic data alone. The parameter predicts none of these, because each is a statement about the joint distribution and it is a statement about a marginal.

\subsubsection{4.1 Soundness on Byte Strings}

Sweeping three policies, one to six anchors, five corpora, and needle lengths 1, 2, 3, 4, 16, 64, 256, 1024 and 2048:

\textbf{582 rows hold. 9,396,207 true occurrences examined. 0 refused.}

A further 210 rows report \texttt{none}, and each is a case where the needle carries fewer distinct offsets than the anchor count. That cause was checked and no row reports \texttt{none} for any other reason.

The degenerate lengths are in the sweep deliberately. A needle of one byte is the shortest thing that can carry an anchor, a needle as long as the corpus leaves one position, and the flat corpus makes every position an occurrence, which offers the largest number of true results available to lose.

\subsubsection{4.2 Soundness Across Geometries}

Sweeping three uninformed selection rules and one to six anchors over the six geometries above and a further eight hypercubes of dimension one through eight:

\textbf{252 rows hold. 213,840 true occurrences examined. 0 refused.}

Candidates admitted, rule \texttt{shuffled}, against the \(N \cdot 2^{-n}\) the two symbol alphabet predicts:

\tiny
\setlength{\tabcolsep}{2pt}
\begin{longtable}{@{}>{\setlength{\rightskip}{0pt}\setlength{\parfillskip}{0pt}\arraybackslash}p{\dimexpr\textwidth/7-2\tabcolsep\relax}>{\setlength{\rightskip}{0pt}\setlength{\parfillskip}{0pt}\arraybackslash}p{\dimexpr\textwidth/7-2\tabcolsep\relax}>{\setlength{\rightskip}{0pt}\setlength{\parfillskip}{0pt}\arraybackslash}p{\dimexpr\textwidth/7-2\tabcolsep\relax}>{\setlength{\rightskip}{0pt}\setlength{\parfillskip}{0pt}\arraybackslash}p{\dimexpr\textwidth/7-2\tabcolsep\relax}>{\setlength{\rightskip}{0pt}\setlength{\parfillskip}{0pt}\arraybackslash}p{\dimexpr\textwidth/7-2\tabcolsep\relax}>{\setlength{\rightskip}{0pt}\setlength{\parfillskip}{0pt}\arraybackslash}p{\dimexpr\textwidth/7-2\tabcolsep\relax}>{\setlength{\rightskip}{0pt}\setlength{\parfillskip}{0pt}\arraybackslash}p{\dimexpr\textwidth/7-2\tabcolsep\relax}@{}}
\toprule
domain & \(n{=}1\) & \(n{=}2\) & \(n{=}3\) & \(n{=}4\) & \(n{=}5\) & \(n{=}6\) \\
\midrule
predicted & 2044.5 & 1022.3 & 511.1 & 255.6 & 127.8 & 63.9 \\
\texttt{line1d} & 2044.7 & 1025.7 & 516.7 & 259.4 & 129.2 & 65.4 \\
\texttt{grid2d\_\allowbreak{}scatter} & 1797.8 & 901.5 & 449.5 & 223.1 & 112.6 & 56.9 \\
\texttt{grid2d\_\allowbreak{}turned} & 1801.7 & 902.7 & 448.6 & 223.3 & 112.1 & 56.7 \\
\texttt{cube3d\_\allowbreak{}box} & 1685.3 & 840.5 & 422.3 & 212.7 & 105.2 & 52.3 \\
\texttt{field1d\_\allowbreak{}complex} & 2046.4 & 1023.4 & 511.9 & 256.2 & 128.9 & 64.7 \\
\bottomrule
\end{longtable}
\normalsize

The grid and cube rows sit below the prediction because those geometries admit fewer base positions, 3600 and 3375 against 4089 on the line. Within each row the halving per anchor is what the model gives.

\textbf{A rotation is invisible.} \texttt{grid2d\_\allowbreak{}scatter} and \texttt{grid2d\_\allowbreak{}turned} are the same eight points a quarter turn apart and agree at every anchor count, which puts Section 2.4's third corollary on a measurement.

\textbf{The complex alphabet behaves as bytes do.} \texttt{field1d\_\allowbreak{}complex} matches \texttt{line1d} to within 0.1\% at every count, with the geometry held fixed and only the symbol type changed.

\textbf{The cost law holds at every dimension swept.} The hypercube cases place one pattern point at the origin and each of the rest one step out along the next axis. The pattern touches a new axis for every point it has and is never a lower dimensional figure sitting in a larger space (\texttt{test/\allowbreak{}bench/\allowbreak{}bench\_\allowbreak{}ancorae\_\allowbreak{}lattice.\allowbreak{}c:686}). Against \(1 + (P-1)2^{-n}\) for \(P\) positions and \(n\) anchors, which includes the one match the pattern is guaranteed against itself:

\tiny
\setlength{\tabcolsep}{2pt}
\begin{longtable}{@{}>{\setlength{\rightskip}{0pt}\setlength{\parfillskip}{0pt}\arraybackslash}p{\dimexpr\textwidth/9-2\tabcolsep\relax}>{\setlength{\rightskip}{0pt}\setlength{\parfillskip}{0pt}\arraybackslash}p{\dimexpr\textwidth/9-2\tabcolsep\relax}>{\setlength{\rightskip}{0pt}\setlength{\parfillskip}{0pt}\arraybackslash}p{\dimexpr\textwidth/9-2\tabcolsep\relax}>{\setlength{\rightskip}{0pt}\setlength{\parfillskip}{0pt}\arraybackslash}p{\dimexpr\textwidth/9-2\tabcolsep\relax}>{\setlength{\rightskip}{0pt}\setlength{\parfillskip}{0pt}\arraybackslash}p{\dimexpr\textwidth/9-2\tabcolsep\relax}>{\setlength{\rightskip}{0pt}\setlength{\parfillskip}{0pt}\arraybackslash}p{\dimexpr\textwidth/9-2\tabcolsep\relax}>{\setlength{\rightskip}{0pt}\setlength{\parfillskip}{0pt}\arraybackslash}p{\dimexpr\textwidth/9-2\tabcolsep\relax}>{\setlength{\rightskip}{0pt}\setlength{\parfillskip}{0pt}\arraybackslash}p{\dimexpr\textwidth/9-2\tabcolsep\relax}>{\setlength{\rightskip}{0pt}\setlength{\parfillskip}{0pt}\arraybackslash}p{\dimexpr\textwidth/9-2\tabcolsep\relax}@{}}
\toprule
dimension & 1 & 2 & 3 & 4 & 5 & 6 & 7 & 8 \\
\midrule
positions & 4089 & 3600 & 4913 & 4096 & 1024 & 729 & 2187 & 256 \\
ratio at \(n{=}1\) & 1.001 & 1.000 & 1.000 & 1.001 & 1.001 & 1.007 & 0.999 & 0.995 \\
ratio at \(n{=}3\) & 0.991 & 0.995 & 1.002 & 1.000 & 0.994 & 1.028 & 0.999 & 0.970 \\
ratio at \(n{=}6\) & 0.969 & 1.020 & 1.017 & 1.001 & 0.980 & 0.996 & 1.013 & 0.979 \\
\bottomrule
\end{longtable}
\normalsize

Across all 48 rows the ratio stays between 0.969 and 1.028. Seven points off the origin can touch at most seven axes, so the dimension eight pattern spans seven of its eight, which is a limit of the pattern size and not of the construction.

An earlier version of this table omitted the guaranteed self match from the prediction and drifted to a ratio of 1.22 at the smallest domain. That is the same error recorded twice already in Section 6.5, made a third time.

\subsubsection{4.3 Selection-Policy Efficiency}

Candidates admitted by one anchor at needle 16, english cost table, fingerprint \texttt{69c2e2df}:

\tiny
\setlength{\tabcolsep}{2pt}
\begin{longtable}{@{}>{\setlength{\rightskip}{0pt}\setlength{\parfillskip}{0pt}\arraybackslash}p{\dimexpr\textwidth/6-2\tabcolsep\relax}>{\setlength{\rightskip}{0pt}\setlength{\parfillskip}{0pt}\arraybackslash}p{\dimexpr\textwidth/6-2\tabcolsep\relax}>{\setlength{\rightskip}{0pt}\setlength{\parfillskip}{0pt}\arraybackslash}p{\dimexpr\textwidth/6-2\tabcolsep\relax}>{\setlength{\rightskip}{0pt}\setlength{\parfillskip}{0pt}\arraybackslash}p{\dimexpr\textwidth/6-2\tabcolsep\relax}>{\setlength{\rightskip}{0pt}\setlength{\parfillskip}{0pt}\arraybackslash}p{\dimexpr\textwidth/6-2\tabcolsep\relax}>{\setlength{\rightskip}{0pt}\setlength{\parfillskip}{0pt}\arraybackslash}p{\dimexpr\textwidth/6-2\tabcolsep\relax}@{}}
\toprule
corpus & \(H_2\) & \texttt{table} & \texttt{random} & \texttt{maxent} & table over maxent \\
\midrule
\texttt{flat} & 0.00 & 2032 & 2032 & 2032 & 1.00 \\
\texttt{structured} & 3.62 & 17.6 & 47.9 & 113.3 & 6.43 \\
\texttt{english} & 3.69 & 26.9 & 108.6 & 149.5 & 5.56 \\
\texttt{periodic16} & 4.61 & 44.4 & 96.8 & 82.3 & 1.85 \\
\texttt{uniform} & 7.83 & 7.53 & 7.84 & 7.61 & 1.01 \\
\bottomrule
\end{longtable}
\normalsize

On \texttt{flat} no anchor can ever fail. All eighteen combinations of policy and anchor count report 2032 candidates out of 2033 positions. On \texttt{uniform} the three policies agree to within 4\%. Between them the informed table is worth 6.4 times.

The column ordering rules out entropy as the governing quantity. \texttt{periodic16} carries more entropy than \texttt{english} and its ratio is less than half as large, and \texttt{structured} has the lowest entropy of the four non-degenerate corpora and the second largest ratio. Section 4.3.1 derives the quantity that does govern it.

\paragraph{4.3.1 Efficiency Bound for Size-Biased Order Statistics}

An informed anchor is the rarest symbol among the \(m\) that a pattern carries, and a pattern drawn from the corpus presents each symbol with probability equal to its own frequency. So the informed rate is the expectation of the minimum of \(m\) size biased draws. Writing

\(F(t) \;=\; \sum_{\sigma:\, p_\sigma \le t} p_\sigma\)

for the size biased distribution function, the survival form of the expectation gives

\[
\mathbb{E}[\min\nolimits_m] \;=\; \int_0^1 \bigl(1 - F(t)\bigr)^m \, dt, \qquad
\text{ceiling}(m) \;=\; \frac{\sum_\sigma p_\sigma^2}{\mathbb{E}[\min_m]}
\]

The numerator is the uninformed rate from Section 2.5, which makes the quotient what a table matched to the corpus would buy over no table at all. The ceiling is an upper bound on every measure, since no rule can pick a symbol rarer than the rarest one present.

\textbf{The ceiling is exactly 1 if and only if \(p\) is uniform on its support.} If every nonzero \(p_\sigma\) equals \(1/k\) then \(F\) is a single step, the integral collapses to \(1/k\), and \(\sum_\sigma p_\sigma^2\) is also \(1/k\). If two nonzero probabilities differ then the rarest of \(m\) draws is strictly below the size biased mean for \(m \ge 2\) and the quotient exceeds one. A point mass and a uniform distribution are both uniform on their support, so both ends of Section 4.3 report a ratio of one, and entropy plays no part in either case.

Computed at needle 16 from each corpus's own frequencies (\texttt{test/\allowbreak{}bench/\allowbreak{}bench\_\allowbreak{}ancorae\_\allowbreak{}sift.\allowbreak{}c:933}):

\tiny
\setlength{\tabcolsep}{2pt}
\begin{longtable}{@{}>{\setlength{\rightskip}{0pt}\setlength{\parfillskip}{0pt}\arraybackslash}p{\dimexpr\textwidth/7-2\tabcolsep\relax}>{\setlength{\rightskip}{0pt}\setlength{\parfillskip}{0pt}\arraybackslash}p{\dimexpr\textwidth/7-2\tabcolsep\relax}>{\setlength{\rightskip}{0pt}\setlength{\parfillskip}{0pt}\arraybackslash}p{\dimexpr\textwidth/7-2\tabcolsep\relax}>{\setlength{\rightskip}{0pt}\setlength{\parfillskip}{0pt}\arraybackslash}p{\dimexpr\textwidth/7-2\tabcolsep\relax}>{\setlength{\rightskip}{0pt}\setlength{\parfillskip}{0pt}\arraybackslash}p{\dimexpr\textwidth/7-2\tabcolsep\relax}>{\setlength{\rightskip}{0pt}\setlength{\parfillskip}{0pt}\arraybackslash}p{\dimexpr\textwidth/7-2\tabcolsep\relax}>{\setlength{\rightskip}{0pt}\setlength{\parfillskip}{0pt}\arraybackslash}p{\dimexpr\textwidth/7-2\tabcolsep\relax}@{}}
\toprule
corpus & \(H_2\) & uninformed rate & informed rate & ceiling & measured & captured \\
\midrule
\texttt{flat} & 0.000 & 1.000000 & 1.000000 & 1.000 & 1.00 & exact \\
\texttt{structured} & 3.620 & 0.081313 & 0.004483 & 18.14 & 6.43 & 35\% \\
\texttt{english} & 3.692 & 0.077371 & 0.011887 & 6.51 & 5.56 & 85\% \\
\texttt{periodic16} & 4.611 & 0.040927 & 0.013708 & 2.99 & 1.85 & 62\% \\
\texttt{uniform} & 7.832 & 0.004389 & 0.002172 & 2.02 & 1.01 & 50\% \\
\bottomrule
\end{longtable}
\normalsize

The ceiling orders the corpora where entropy does not, and it is exact at the point mass.

Two readings of the gap between ceiling and measured are separable. On \texttt{english} the linked table is the English table, and 85\% is the highest fraction captured anywhere in the study. That is the matched table effect appearing as a number. On \texttt{structured} the same table reaches 35\% of a much larger ceiling, which is a mismatch cost.

The \texttt{uniform} row needs its own reading. Its ceiling of 2.02 is not a property of the source, which is uniform and therefore has a ceiling of exactly 1 by the statement above. The ceiling is a property of the sample: 2033 draws over 256 symbols do not land uniformly, and the resulting fluctuation leaves something for an oracle to exploit. No fixed table can capture it, because the fluctuation is not a feature of the source and a table cannot know which symbols this particular corpus happened to under-represent. The measured 1.01 is that impossibility, quantified.

\subsubsection{4.4 Collision-Entropy Prediction}

Predicted \texttt{maxent} candidate count against measured, with one subtracted from each prediction because the cost rows exclude the needle's own occurrence:

\tiny
\setlength{\tabcolsep}{2pt}
\begin{longtable}{@{}>{\setlength{\rightskip}{0pt}\setlength{\parfillskip}{0pt}\arraybackslash}p{\dimexpr\textwidth/5-2\tabcolsep\relax}>{\setlength{\rightskip}{0pt}\setlength{\parfillskip}{0pt}\arraybackslash}p{\dimexpr\textwidth/5-2\tabcolsep\relax}>{\setlength{\rightskip}{0pt}\setlength{\parfillskip}{0pt}\arraybackslash}p{\dimexpr\textwidth/5-2\tabcolsep\relax}>{\setlength{\rightskip}{0pt}\setlength{\parfillskip}{0pt}\arraybackslash}p{\dimexpr\textwidth/5-2\tabcolsep\relax}>{\setlength{\rightskip}{0pt}\setlength{\parfillskip}{0pt}\arraybackslash}p{\dimexpr\textwidth/5-2\tabcolsep\relax}@{}}
\toprule
corpus & \(H_2\) & predicted & measured & error \\
\midrule
\texttt{flat} & 0.000 & 2033 & 2032 & 0.05\% \\
\texttt{structured} & 3.620 & 112.3 & 113.3 & 0.9\% \\
\texttt{english} & 3.692 & 156.3 & 149.5 & 4.3\% \\
\texttt{periodic16} & 4.611 & 82.2 & 82.3 & 0.1\% \\
\texttt{uniform} & 7.832 & 7.92 & 7.61 & 4.0\% \\
\bottomrule
\end{longtable}
\normalsize

The residual is sampling bias in the needle draw. Needles are taken at a fixed stride. On English the anchor byte is a biased letter draw, and the two rows with the largest error are the two whose stride interacts with their structure.

\paragraph{4.4.1 Refutation Radius}

Sections 2.2 and 4.1 treat a refutation as settling one position, and it settles many.

Testing the anchor at pattern offset \(a\) reads the cell at \(s+a\) and finds some symbol \(c\). For any shift \(\delta\), the alignment beginning at \(s+\delta\) places pattern offset \(a-\delta\) on that same cell. That alignment can hold only if \(\nu_{a-\delta} = c\). So one read settles every shift whose pattern symbol differs from what was found, and the number of them is

\(R \;=\; m \;-\; \bigl|\{\, i : \nu_i = c \,\}\bigr|\)

A symbol the pattern does not carry refutes every alignment touching the cell. Taking the expectation over a corpus, with the pattern drawn from that corpus so \(c\) arrives with its own frequency,

\(\mathbb{E}[R] \;=\; m\Bigl(1 - \sum_\sigma p_\sigma^2\Bigr) \;=\; m\bigl(1 - 2^{-H_2}\bigr)\)

The quantity that sets the candidate rate in Section 4.4 sets the refutation distance here, appearing once as a rate and once as a length. Measured at every needle length over every corpus (\texttt{test/\allowbreak{}bench/\allowbreak{}bench\_\allowbreak{}ancorae\_\allowbreak{}sift.\allowbreak{}c:1024}):

\tiny
\setlength{\tabcolsep}{2pt}
\begin{longtable}{@{}>{\setlength{\rightskip}{0pt}\setlength{\parfillskip}{0pt}\arraybackslash}p{\dimexpr\textwidth/6-2\tabcolsep\relax}>{\setlength{\rightskip}{0pt}\setlength{\parfillskip}{0pt}\arraybackslash}p{\dimexpr\textwidth/6-2\tabcolsep\relax}>{\setlength{\rightskip}{0pt}\setlength{\parfillskip}{0pt}\arraybackslash}p{\dimexpr\textwidth/6-2\tabcolsep\relax}>{\setlength{\rightskip}{0pt}\setlength{\parfillskip}{0pt}\arraybackslash}p{\dimexpr\textwidth/6-2\tabcolsep\relax}>{\setlength{\rightskip}{0pt}\setlength{\parfillskip}{0pt}\arraybackslash}p{\dimexpr\textwidth/6-2\tabcolsep\relax}>{\setlength{\rightskip}{0pt}\setlength{\parfillskip}{0pt}\arraybackslash}p{\dimexpr\textwidth/6-2\tabcolsep\relax}@{}}
\toprule
corpus & \(m{=}4\) & \(m{=}16\) & \(m{=}64\) & \(m{=}256\) & refuted over \(m\) \\
\midrule
\texttt{english} & 0.9962 & 0.9997 & 1.0000 & 1.0002 & 0.923 \\
\texttt{structured} & 0.9939 & 0.9991 & 0.9985 & 0.9934 & 0.917 \\
\texttt{periodic16} & 0.9999 & 0.9999 & 1.0001 & 1.0001 & 0.959 \\
\texttt{uniform} & 1.0001 & 1.0000 & 1.0000 & 1.0000 & 0.996 \\
\texttt{flat} & exact & exact & exact & exact & 0.000 \\
\bottomrule
\end{longtable}
\normalsize

Entries are measured over predicted. This policy produces 35 rows, of which 28 carry a ratio and it stays between 0.9934 and 1.0011. Over all three policies, 84 rows carry a ratio and the range is 0.9932 to 1.0011. On \texttt{flat} both sides are exactly zero at every length, since every cell holds the one byte the pattern is made of and no read can settle anything.

The last column is constant in \(m\) within each corpus, which the expectation requires: the fraction of alignments a read settles is \(1 - 2^{-H_2}\) and belongs to the corpus alone. One read settles 92\% of the alignments touching it on English prose and 99.6\% on a uniform corpus.

This is the bad character shift of classical string search, reached from Proposition 1 instead of from a shift table, and priced by the collision probability. Nothing in the derivation uses an order on positions, and the result carries to the geometries of Section 4.2, though it was measured only on byte strings.

\subsubsection{4.5 Candidate Count as an Entropy Estimator}

Bias and root mean square error in bits, 64 trials, probe budget 64:

\tiny
\setlength{\tabcolsep}{2pt}
\begin{longtable}{@{}>{\setlength{\rightskip}{0pt}\setlength{\parfillskip}{0pt}\arraybackslash}p{\dimexpr\textwidth/7-2\tabcolsep\relax}>{\setlength{\rightskip}{0pt}\setlength{\parfillskip}{0pt}\arraybackslash}p{\dimexpr\textwidth/7-2\tabcolsep\relax}>{\setlength{\rightskip}{0pt}\setlength{\parfillskip}{0pt}\arraybackslash}p{\dimexpr\textwidth/7-2\tabcolsep\relax}>{\setlength{\rightskip}{0pt}\setlength{\parfillskip}{0pt}\arraybackslash}p{\dimexpr\textwidth/7-2\tabcolsep\relax}>{\setlength{\rightskip}{0pt}\setlength{\parfillskip}{0pt}\arraybackslash}p{\dimexpr\textwidth/7-2\tabcolsep\relax}>{\setlength{\rightskip}{0pt}\setlength{\parfillskip}{0pt}\arraybackslash}p{\dimexpr\textwidth/7-2\tabcolsep\relax}>{\setlength{\rightskip}{0pt}\setlength{\parfillskip}{0pt}\arraybackslash}p{\dimexpr\textwidth/7-2\tabcolsep\relax}@{}}
\toprule
source & \(N\) & true \(H_2\) & plug in & unbiased & probe, self & probe, no self \\
\midrule
\texttt{uniform256} & 256 & 8.000 & -0.984 / 0.986 & +0.029 / 0.126 & -0.964 / 0.969 & +0.074 / 0.210 \\
\texttt{uniform256} & 1024 & 8.000 & -0.325 / 0.325 & -0.005 / 0.029 & -0.324 / 0.335 & -0.003 / 0.108 \\
\texttt{uniform256} & 4096 & 8.000 & -0.088 / 0.089 & -0.001 / 0.007 & -0.077 / 0.085 & +0.010 / 0.040 \\
\texttt{zipf1.\allowbreak{}0} & 4096 & 4.515 & -0.008 / 0.061 & -0.001 / 0.061 & -0.050 / 0.204 & -0.042 / 0.203 \\
\texttt{skew0.\allowbreak{}99} & 4096 & 0.029 & -0.000 / 0.005 & -0.000 / 0.005 & +0.000 / 0.019 & +0.000 / 0.020 \\
\bottomrule
\end{longtable}
\normalsize

**The probe counting its own match reproduces the plug in estimator, and excluding it reproduces the unbiased estimator.** At \texttt{uniform256} and \(N=256\) the plug in bias is \(-0.984\) bits and the probe with self is \(-0.964\); the unbiased bias is \(+0.029\) and the probe without self is \(+0.074\). The subtraction that Section 6.5 records as a measurement error is the standard bias correction, reached from the other direction.

\textbf{The probe is worse than both published estimators at every setting tested.} Its error is set by the probe budget and not by the corpus length. Root mean square error at \(N = 16384\), \texttt{zipf1.\allowbreak{}0}:

\tiny
\setlength{\tabcolsep}{2pt}
\begin{longtable}{@{}>{\setlength{\rightskip}{0pt}\setlength{\parfillskip}{0pt}\arraybackslash}p{\dimexpr\textwidth/6-2\tabcolsep\relax}>{\setlength{\rightskip}{0pt}\setlength{\parfillskip}{0pt}\arraybackslash}p{\dimexpr\textwidth/6-2\tabcolsep\relax}>{\setlength{\rightskip}{0pt}\setlength{\parfillskip}{0pt}\arraybackslash}p{\dimexpr\textwidth/6-2\tabcolsep\relax}>{\setlength{\rightskip}{0pt}\setlength{\parfillskip}{0pt}\arraybackslash}p{\dimexpr\textwidth/6-2\tabcolsep\relax}>{\setlength{\rightskip}{0pt}\setlength{\parfillskip}{0pt}\arraybackslash}p{\dimexpr\textwidth/6-2\tabcolsep\relax}>{\setlength{\rightskip}{0pt}\setlength{\parfillskip}{0pt}\arraybackslash}p{\dimexpr\textwidth/6-2\tabcolsep\relax}@{}}
\toprule
probe budget & 16 & 64 & 256 & 1024 & 4096 \\
\midrule
probe, no self & 0.602 & 0.252 & 0.122 & 0.062 & 0.039 \\
unbiased & 0.028 & 0.028 & 0.028 & 0.028 & 0.028 \\
ratio & 21.7 & 9.1 & 4.4 & 2.2 & 1.4 \\
\bottomrule
\end{longtable}
\normalsize

The probe's error falls as one over the square root of the budget and converges toward the unbiased estimator without reaching it. The probe is also more expensive, at \(O(kN)\) against \(O(N)\) for a histogram. Its one advantage is the constraint it survives: it builds no histogram and never enumerates the alphabet. The probe applies where the other two cannot be computed at all. In a search it is free, because the candidate count is a byproduct of work already being done.

\paragraph{4.5.1 Falling-Factorial Bias Correction}

The correction excludes a position from being paired with itself. Applied \(k\) deep it excludes every tuple in which any two indices coincide, which is a falling factorial in the numerator and the denominator, giving the unbiased estimator of \(\sum_\sigma p(\sigma)^k\) (\texttt{test/\allowbreak{}bench/\allowbreak{}bench\_\allowbreak{}ancorae\_\allowbreak{}entropy.\allowbreak{}c:463-\allowbreak{}500}). So the correction can be climbed, and each rung estimates Renyi entropy at the next order.

At order two the correction has a closed form. Writing \(C = \sum_\sigma \hat p(\sigma)^2\) for the plug in collision probability, the corrected form is \((CN-1)/(N-1)\), so the two differ by

\[
\Delta C \;=\; \frac{1 - C}{N-1}, \qquad
\Delta H_2 \;\approx\; \frac{1-C}{(N-1)\,C\,\ln 2}\ \text{bits}
\]

That numerator answers the question the ladder was built for. A point mass has \(C = 1\) exactly, so \(1 - C\) is exactly zero and the correction is exactly zero, for every length. The same holds at every order: with one symbol \(n_\sigma = N\), the plug in form gives \(\sum \hat p^m = 1\) and the corrected form gives \(N^{(m)}/N^{(m)} = 1\). Measured at orders 2 through 5 and lengths 256, 1024, 4096 and 16384, every entry is 0.0000 bits.

Correction magnitude in bits at order 2, \(N = 4096\), 64 trials, against the closed form:

\tiny
\setlength{\tabcolsep}{2pt}
\begin{longtable}{@{}>{\setlength{\rightskip}{0pt}\setlength{\parfillskip}{0pt}\arraybackslash}p{\dimexpr\textwidth/4-2\tabcolsep\relax}>{\setlength{\rightskip}{0pt}\setlength{\parfillskip}{0pt}\arraybackslash}p{\dimexpr\textwidth/4-2\tabcolsep\relax}>{\setlength{\rightskip}{0pt}\setlength{\parfillskip}{0pt}\arraybackslash}p{\dimexpr\textwidth/4-2\tabcolsep\relax}>{\setlength{\rightskip}{0pt}\setlength{\parfillskip}{0pt}\arraybackslash}p{\dimexpr\textwidth/4-2\tabcolsep\relax}@{}}
\toprule
source & \(H_2\) & predicted & measured \\
\midrule
\texttt{point1} & 0.000 & 0.0000 & 0.0000 \\
\texttt{skew0.\allowbreak{}99} & 0.029 & 0.0000 & 0.0000 \\
\texttt{skew0.\allowbreak{}90} & 0.286 & 0.0001 & 0.0001 \\
\texttt{uniform2} & 1.000 & 0.0004 & 0.0004 \\
\texttt{zipf1.\allowbreak{}5} & 2.364 & 0.0015 & 0.0015 \\
\texttt{uniform16} & 4.000 & 0.0053 & 0.0053 \\
\texttt{zipf1.\allowbreak{}0} & 4.515 & 0.0077 & 0.0077 \\
\texttt{zipf0.\allowbreak{}5} & 7.254 & 0.0534 & 0.0525 \\
\texttt{uniform256} & 8.000 & 0.0898 & 0.0871 \\
\bottomrule
\end{longtable}
\normalsize

The form holds across four orders of magnitude in the correction, with the largest disagreement at \texttt{uniform256} where the first order approximation of the logarithm is weakest.

The quantity governing the size is how many times a symbol was seen, since a falling factorial differs from a plain power in proportion to that count. A head symbol appearing 4055 times out of 4096 loses almost nothing by refusing to pair with itself. A symbol appearing 16 times loses a measurable fraction.

Climbing the ladder does not settle on one behavior, and which way a source moves depends on where its mass sits. On a spread source the correction grows with every rung, because thin counts stay thin at higher order: \texttt{uniform256} gives 0.0871, 0.1256, 0.1616 and 0.1955 at orders 2 through 5. On a skewed source it shrinks, because higher moments concentrate on the head symbol where the count is large: \texttt{zipf0.\allowbreak{}5} gives 0.0525, 0.0393, 0.0327 and 0.0331. On the point mass it is zero at every rung.

This correction and the ceiling in Section 4.3.1 answer to different properties of a distribution and neither one predicts the other. The ceiling is one when the probabilities are equal on the support, at any entropy and any corpus length. The correction is zero only at a point mass, and its size is set by how many times each symbol was seen. The correction falls with corpus length while the ceiling does not. A uniform source and a point mass share a ceiling of one and have corrections of 0.0871 and 0.0000 bits at \(N = 4096\). The two were measured on different source sets and no joint sweep was run.

\subsubsection{4.6 Dependence in Anchor Cascades}

Ratio of observed candidates to the independence prediction, needle 16 (\texttt{test/\allowbreak{}bench/\allowbreak{}bench\_\allowbreak{}ancorae\_\allowbreak{}sift.\allowbreak{}c:541-\allowbreak{}632}):

\tiny
\setlength{\tabcolsep}{2pt}
\begin{longtable}{@{}>{\setlength{\rightskip}{0pt}\setlength{\parfillskip}{0pt}\arraybackslash}p{\dimexpr\textwidth/7-2\tabcolsep\relax}>{\setlength{\rightskip}{0pt}\setlength{\parfillskip}{0pt}\arraybackslash}p{\dimexpr\textwidth/7-2\tabcolsep\relax}>{\setlength{\rightskip}{0pt}\setlength{\parfillskip}{0pt}\arraybackslash}p{\dimexpr\textwidth/7-2\tabcolsep\relax}>{\setlength{\rightskip}{0pt}\setlength{\parfillskip}{0pt}\arraybackslash}p{\dimexpr\textwidth/7-2\tabcolsep\relax}>{\setlength{\rightskip}{0pt}\setlength{\parfillskip}{0pt}\arraybackslash}p{\dimexpr\textwidth/7-2\tabcolsep\relax}>{\setlength{\rightskip}{0pt}\setlength{\parfillskip}{0pt}\arraybackslash}p{\dimexpr\textwidth/7-2\tabcolsep\relax}>{\setlength{\rightskip}{0pt}\setlength{\parfillskip}{0pt}\arraybackslash}p{\dimexpr\textwidth/7-2\tabcolsep\relax}@{}}
\toprule
corpus & policy & \(n{=}2\) & \(n{=}3\) & \(n{=}4\) & \(n{=}5\) & \(n{=}6\) \\
\midrule
\texttt{uniform} & \texttt{table} & 1.23 & - & - & - & - \\
\texttt{english} & \texttt{table} & 1.44 & 3.80 & 7.68 & 160.7 & - \\
\texttt{structured} & \texttt{table} & 8.46 & 163.2 & 663.4 & 2226.5 & 8186.1 \\
\texttt{english} & \texttt{maxent} & 2.20 & 12.1 & 77.4 & 438.4 & 2409.4 \\
\texttt{structured} & \texttt{maxent} & 2.34 & 8.53 & 28.9 & 118.0 & 389.7 \\
\bottomrule
\end{longtable}
\normalsize

A dash marks a cell where the survivors reached zero.

The product rule overpromises and the error compounds with every anchor added. On \texttt{structured} under the informed table it is wrong by three orders of magnitude at four anchors and by four at six. The direction is the dangerous one: more candidates survive than the arithmetic predicts. On \texttt{uniform} the ratio holds at 1.23 and the survivors reach zero by three anchors. That is the regime where the arithmetic is exact.

\paragraph{4.6.1 Cascade Depth Under Independence}

Section 2.5 budgets a search in bits: reducing \(N\) candidates to \(O(1)\) survivors needs \(\log_2 N\) bits and one anchor supplies \(-\log_2 q\) of them. For an uninformed anchor that is \(H_2\) bits, and a stack of \(n\) exhausts the space at

\(n \;=\; \frac{\log_2 N}{H_2}\)

Against the first anchor count whose measured excess candidate count reaches zero:

\tiny
\setlength{\tabcolsep}{2pt}
\begin{longtable}{@{}>{\setlength{\rightskip}{0pt}\setlength{\parfillskip}{0pt}\arraybackslash}p{\dimexpr\textwidth/6-2\tabcolsep\relax}>{\setlength{\rightskip}{0pt}\setlength{\parfillskip}{0pt}\arraybackslash}p{\dimexpr\textwidth/6-2\tabcolsep\relax}>{\setlength{\rightskip}{0pt}\setlength{\parfillskip}{0pt}\arraybackslash}p{\dimexpr\textwidth/6-2\tabcolsep\relax}>{\setlength{\rightskip}{0pt}\setlength{\parfillskip}{0pt}\arraybackslash}p{\dimexpr\textwidth/6-2\tabcolsep\relax}>{\setlength{\rightskip}{0pt}\setlength{\parfillskip}{0pt}\arraybackslash}p{\dimexpr\textwidth/6-2\tabcolsep\relax}>{\setlength{\rightskip}{0pt}\setlength{\parfillskip}{0pt}\arraybackslash}p{\dimexpr\textwidth/6-2\tabcolsep\relax}@{}}
\toprule
corpus & \(H_2\) & \(N\) & predicted & measured, \(m = 64\) and 256 & measured, \(m = 16\) \\
\midrule
\texttt{uniform} & 7.832 & 2048 & 1.40 & 3 & 3 \\
\texttt{periodic16} & 4.611 & 2048 & 2.39 & 3 to 4 & 4 \\
\texttt{structured} & 3.620 & 1408 & 2.89 & beyond 6 & beyond 6 \\
\texttt{english} & 3.692 & 2048 & 2.98 & 3 & 6 \\
\texttt{flat} & 0.000 & 2048 & unbounded & never & never \\
\bottomrule
\end{longtable}
\normalsize

The prediction assumes each anchor contributes its full \(H_2\) bits, which requires them to be independent. Section 4.7 puts that at separation 7 and above. At \(m = 64\) and beyond there is room for that and the bound is exact on English, predicting 2.98 against a measured 3. At \(m = 16\) there is not, the anchors sit close enough to be correlated, each supplies less than \(H_2\) bits, and English needs 6.

**The gap between predicted and measured exhaustion is the independence failure of Section 4.6, read a second way.** \texttt{structured} has the largest cascade ratio and the largest gap. \texttt{flat} carries no bits at all. No stack of any size removes anything, and the row that never exhausts is the same corpus every other instrument reports zero on.

A stack terminates at three or four here, and it terminates because the space is spent and not because the arithmetic fails.

\subsubsection{4.7 Dependence by Separation and Periodicity}

Independence \(I(d)\) against anchor separation, needle 64, policy \texttt{table}:

\tiny
\setlength{\tabcolsep}{2pt}
\begin{longtable}{@{}>{\setlength{\rightskip}{0pt}\setlength{\parfillskip}{0pt}\arraybackslash}p{\dimexpr\textwidth/9-2\tabcolsep\relax}>{\setlength{\rightskip}{0pt}\setlength{\parfillskip}{0pt}\arraybackslash}p{\dimexpr\textwidth/9-2\tabcolsep\relax}>{\setlength{\rightskip}{0pt}\setlength{\parfillskip}{0pt}\arraybackslash}p{\dimexpr\textwidth/9-2\tabcolsep\relax}>{\setlength{\rightskip}{0pt}\setlength{\parfillskip}{0pt}\arraybackslash}p{\dimexpr\textwidth/9-2\tabcolsep\relax}>{\setlength{\rightskip}{0pt}\setlength{\parfillskip}{0pt}\arraybackslash}p{\dimexpr\textwidth/9-2\tabcolsep\relax}>{\setlength{\rightskip}{0pt}\setlength{\parfillskip}{0pt}\arraybackslash}p{\dimexpr\textwidth/9-2\tabcolsep\relax}>{\setlength{\rightskip}{0pt}\setlength{\parfillskip}{0pt}\arraybackslash}p{\dimexpr\textwidth/9-2\tabcolsep\relax}>{\setlength{\rightskip}{0pt}\setlength{\parfillskip}{0pt}\arraybackslash}p{\dimexpr\textwidth/9-2\tabcolsep\relax}>{\setlength{\rightskip}{0pt}\setlength{\parfillskip}{0pt}\arraybackslash}p{\dimexpr\textwidth/9-2\tabcolsep\relax}@{}}
\toprule
stride & 1 & 2 & 3 & 5 & 7 & 8 & 13 & 17 \\
\midrule
\texttt{english} & 6.05 & 2.59 & 1.50 & 1.21 & 0.99 & 1.02 & 1.07 & 1.21 \\
\texttt{structured} & 6.20 & 3.70 & 3.79 & 2.88 & 1.46 & 1.54 & 1.07 & 1.17 \\
\bottomrule
\end{longtable}
\normalsize

Correlation is a function of distance and not of needle length, dying by stride 7 in both corpora. The \texttt{uniform} control is reported as \(z\) and not as \(I\): its joint excess runs from 0.00 to 0.09 against a prediction of 0.03, so the ratio swings between 0.00 and 3.03 on counts too small to carry one. Every \(|z|\) in that row is at most 1.7, consistent with independence at every stride measured.

On \texttt{periodic16}, reported as \(z\):

\tiny
\setlength{\tabcolsep}{2pt}
\begin{longtable}{@{}>{\setlength{\rightskip}{0pt}\setlength{\parfillskip}{0pt}\arraybackslash}p{\dimexpr\textwidth/6-2\tabcolsep\relax}>{\setlength{\rightskip}{0pt}\setlength{\parfillskip}{0pt}\arraybackslash}p{\dimexpr\textwidth/6-2\tabcolsep\relax}>{\setlength{\rightskip}{0pt}\setlength{\parfillskip}{0pt}\arraybackslash}p{\dimexpr\textwidth/6-2\tabcolsep\relax}>{\setlength{\rightskip}{0pt}\setlength{\parfillskip}{0pt}\arraybackslash}p{\dimexpr\textwidth/6-2\tabcolsep\relax}>{\setlength{\rightskip}{0pt}\setlength{\parfillskip}{0pt}\arraybackslash}p{\dimexpr\textwidth/6-2\tabcolsep\relax}>{\setlength{\rightskip}{0pt}\setlength{\parfillskip}{0pt}\arraybackslash}p{\dimexpr\textwidth/6-2\tabcolsep\relax}@{}}
\toprule
stride & 3 & 7 & 11 & 16 & 17 \\
\midrule
needle 32 & 0.9 & \textbf{4.0} & -0.7 & \textbf{5.3} & 0.7 \\
needle 64 & -0.3 & \textbf{6.6} & -0.1 & \textbf{12.3} & 0.4 \\
needle 128 & -0.1 & \textbf{6.5} & 0.7 & \textbf{13.1} & 0.5 \\
needle 256 & -0.4 & \textbf{5.5} & 0.7 & \textbf{12.6} & 1.4 \\
\bottomrule
\end{longtable}
\normalsize

A stride equal to the record period is the strongest correlation measured anywhere here, \(I = 5.82\) at \(z = 12.3\), reproducing at every needle length. A coprime stride of 17 is clean.

\(I(d)\) is an exponentiated pointwise mutual information, so \(\mathrm{PMI}(d) = \log_2 I(d)\) and the second anchor's contribution is \(-\log_2 p_\beta - \mathrm{PMI}(d)\). At \(I = 5.82\) the loss is 2.54 bits, so an anchor nominally worth 8 delivers 5.5.

For fixed width records of period \(T\) with field offsets \(F\), a layout bound predicts

\(D_x \;\subseteq\; \{T\}\ \cup\ \{\,|f_i-f_j| \;:\; f_i,f_j\in F\,\}\)

Here \(T = 16\) and the punctuation sits at \(F = \{4,11,15\}\), so \(\hat{D} = \{4,7,11,16\}\). Correlated is \(|z| > 2\), and lags 1 and 2 are excluded because short range correlation is present in every corpus here, English included at \(I(1) = 6.05\).

\tiny
\setlength{\tabcolsep}{2pt}
\begin{longtable}{@{}>{\setlength{\rightskip}{0pt}\setlength{\parfillskip}{0pt}\arraybackslash}p{\dimexpr\textwidth/6-2\tabcolsep\relax}>{\setlength{\rightskip}{0pt}\setlength{\parfillskip}{0pt}\arraybackslash}p{\dimexpr\textwidth/6-2\tabcolsep\relax}>{\setlength{\rightskip}{0pt}\setlength{\parfillskip}{0pt}\arraybackslash}p{\dimexpr\textwidth/6-2\tabcolsep\relax}>{\setlength{\rightskip}{0pt}\setlength{\parfillskip}{0pt}\arraybackslash}p{\dimexpr\textwidth/6-2\tabcolsep\relax}>{\setlength{\rightskip}{0pt}\setlength{\parfillskip}{0pt}\arraybackslash}p{\dimexpr\textwidth/6-2\tabcolsep\relax}>{\setlength{\rightskip}{0pt}\setlength{\parfillskip}{0pt}\arraybackslash}p{\dimexpr\textwidth/6-2\tabcolsep\relax}@{}}
\toprule
\(d\) & in \(\hat{D}\) & \(I(d)\) & \(z\) & measured & verdict \\
\midrule
3 & no & 0.95 & -0.3 & independent & agrees \\
4 & \textbf{yes} & 0.96 & -0.3 & independent & \textbf{necessity fails} \\
5 & no & 1.31 & 2.1 & correlated & \textbf{sufficiency fails} \\
7 & \textbf{yes} & 2.46 & 6.6 & correlated & agrees \\
8 & no & 1.64 & 4.1 & correlated & \textbf{sufficiency fails} \\
11 & \textbf{yes} & 0.99 & -0.1 & independent & \textbf{necessity fails} \\
13 & no & 1.13 & 0.9 & independent & agrees \\
16 & \textbf{yes} & 5.82 & 12.3 & correlated & agrees \\
17 & no & 1.08 & 0.4 & independent & agrees \\
\bottomrule
\end{longtable}
\normalsize

Five of nine agree, and the four failures are the result. The \(d = 5\) row sits at \(z = 2.1\) against a threshold of 2.0.

Necessity fails at \(d \in \{4,11\}\) because a predicted distance only bites when the picker selects both of its endpoints, and selection is driven by the cost table. Sufficiency fails at \(d = 8\) because \(8 = T/2\) and a periodic sequence carries autocorrelation at divisors of its period, which the bound as written omits. Adding harmonics would capture that and would widen \(\hat{D}\), worsening the necessity side. So \(\hat{D}\) bounds the risky set in neither direction, and the rule that survives is the weak one: a stride is safe when it has been measured to be.

One confound is stated here instead of left for a reader to find. The anchor is picked by the English cost table applied to non-English data. These rows measure the profile mismatch as well as the periodicity.

\subsubsection{4.8 Reference-Measure Mismatch}

\(I(d)\) and the skip are both divergences from a reference, and the cost table is that reference. A build links exactly one of five. Skip distance at needle 64, stride 0:

\tiny
\setlength{\tabcolsep}{2pt}
\begin{longtable}{@{}>{\setlength{\rightskip}{0pt}\setlength{\parfillskip}{0pt}\arraybackslash}p{\dimexpr\textwidth/5-2\tabcolsep\relax}>{\setlength{\rightskip}{0pt}\setlength{\parfillskip}{0pt}\arraybackslash}p{\dimexpr\textwidth/5-2\tabcolsep\relax}>{\setlength{\rightskip}{0pt}\setlength{\parfillskip}{0pt}\arraybackslash}p{\dimexpr\textwidth/5-2\tabcolsep\relax}>{\setlength{\rightskip}{0pt}\setlength{\parfillskip}{0pt}\arraybackslash}p{\dimexpr\textwidth/5-2\tabcolsep\relax}>{\setlength{\rightskip}{0pt}\setlength{\parfillskip}{0pt}\arraybackslash}p{\dimexpr\textwidth/5-2\tabcolsep\relax}@{}}
\toprule
corpus & english & generic & uri & route \\
\midrule
english prose & 159.1 & \textbf{168.0} & 5.8 & 6.0 \\
C source & 147.7 & \textbf{206.7} & 5.9 & 6.8 \\
periodic records & 73.2 & 73.2 & 16.6 & 9.5 \\
uniform (control) & 277.1 & 256.0 & 277.7 & 277.1 \\
\bottomrule
\end{longtable}
\normalsize

This table was measured against an earlier control generator and its uniform row is on that generator. The claim it carries is that the four profiles agree there, and that claim does not depend on which uniform source was used.

The wrong reference costs 25 to 30 times. \texttt{uri} and \texttt{route} take the skip on real text from about 160 to about 6. That is the largest effect measured anywhere in this document.

The prediction that a corpus matched table wins is refuted. \texttt{generic} beats \texttt{english} on English prose and beats it decisively on C source. So the case for instantiating the table is not that it should match the data, which is unsupported. The case is that the wrong table costs 25 times and a build holds one.

On the uniform corpus every reference performs the same, at 277.1, 256.0, 277.7 and 277.1. At maximum entropy the choice of \(Q\) stops mattering, because there is no gradient for a reference to be a reference to, and \(D(P\|P) = 0\). The policy sweep in Section 4.3 reaches the same conclusion from the other side, removing the reference entirely instead of replacing it, and the uniform corpus again refuses to separate.

\subsubsection{4.9 Ordered Search Versus Anchor Filtering}

Sections 4.3 through 4.6 measure candidate rates, which is not the same objective as the work a whole search performs. \texttt{bench\_\allowbreak{}ancorae\_\allowbreak{}ab.\allowbreak{}c} puts every arm on one corpus, one needle set, and one counter, which is corpus symbol accesses, and checks every arm's occurrence set against a brute force scan. A row whose arms disagree prints BROKEN and is not read.

\paragraph{4.9.1 Rare-Anchor Search Cost}

An anchor at offset \(a\) can advance the search by at most \(a+1\), since past that the pattern has left the cell behind and the read constrains nothing. Horspool anchors at \(m-1\), the largest offset a pattern has. That takes the largest advance available and the worst candidate rate. Anchoring at the rarest byte takes the best candidate rate and gives up advance.

Corpus reads per search, mean over 64 needles, english cost table:

\tiny
\setlength{\tabcolsep}{2pt}
\begin{longtable}{@{}>{\setlength{\rightskip}{0pt}\setlength{\parfillskip}{0pt}\arraybackslash}p{\dimexpr\textwidth/6-2\tabcolsep\relax}>{\setlength{\rightskip}{0pt}\setlength{\parfillskip}{0pt}\arraybackslash}p{\dimexpr\textwidth/6-2\tabcolsep\relax}>{\setlength{\rightskip}{0pt}\setlength{\parfillskip}{0pt}\arraybackslash}p{\dimexpr\textwidth/6-2\tabcolsep\relax}>{\setlength{\rightskip}{0pt}\setlength{\parfillskip}{0pt}\arraybackslash}p{\dimexpr\textwidth/6-2\tabcolsep\relax}>{\setlength{\rightskip}{0pt}\setlength{\parfillskip}{0pt}\arraybackslash}p{\dimexpr\textwidth/6-2\tabcolsep\relax}>{\setlength{\rightskip}{0pt}\setlength{\parfillskip}{0pt}\arraybackslash}p{\dimexpr\textwidth/6-2\tabcolsep\relax}@{}}
\toprule
corpus & \(m\) & Horspool & rare anchor & salted anchor & rare over Horspool \\
\midrule
\texttt{english} & 16 & 361.0 & 671.8 & 694.9 & 1.86 \\
\texttt{structured} & 16 & 170.3 & 363.0 & 332.8 & 2.13 \\
\texttt{periodic16} & 16 & 360.1 & 1551.2 & 859.1 & 4.31 \\
\texttt{uniform} & 16 & 281.6 & 3517.0 & 720.5 & 12.49 \\
\bottomrule
\end{longtable}
\normalsize

The rare anchor loses on every corpus and every needle length tested, by 1.2 to 25 times, on the mean and on the worst case over the 64 needles. Salting the offset does not recover it. Adding the rare byte as a filter on top of Horspool buys between 0.0\% and 2.6\%.

The loss is largest on \texttt{uniform}, where no rare byte exists and the cost table selects an offset unrelated to the data, and the advance ceiling collapses with it. Search cost is governed by advance distance and candidate rate is not the binding constraint.

This comparison is only available on a line. Horspool needs a total order to have a last position and translation along it to shift, and neither exists for the index sets of Section 4.2. What it measures is the cost of generality where the extra structure is present to be exploited.

\paragraph{4.9.2 Joint Objective: Rejection and Advance}

Rejection rate and advance distance are two factors of one quantity. A read at offset \(a\) returning symbol \(c\) advances by \(\mathrm{adv}_a(c)\), so an anchor is worth

\(\mathbb{E}[\text{advance}] \;=\; \sum_c p(c)\,\mathrm{adv}_a(c), \qquad \mathrm{adv}_a(c) \le a+1\)

Horspool maximizes the ceiling and ignores the frequencies. The rare anchor maximizes the rejection rate and ignores the ceiling. Choosing the offset that maximizes the product, using the corpus's own frequencies:

\tiny
\setlength{\tabcolsep}{2pt}
\begin{longtable}{@{}>{\setlength{\rightskip}{0pt}\setlength{\parfillskip}{0pt}\arraybackslash}p{\dimexpr\textwidth/4-2\tabcolsep\relax}>{\setlength{\rightskip}{0pt}\setlength{\parfillskip}{0pt}\arraybackslash}p{\dimexpr\textwidth/4-2\tabcolsep\relax}>{\setlength{\rightskip}{0pt}\setlength{\parfillskip}{0pt}\arraybackslash}p{\dimexpr\textwidth/4-2\tabcolsep\relax}>{\setlength{\rightskip}{0pt}\setlength{\parfillskip}{0pt}\arraybackslash}p{\dimexpr\textwidth/4-2\tabcolsep\relax}@{}}
\toprule
corpus & \(m\) & chosen offset of \(m-1 = 31\) & vs Horspool \\
\midrule
\texttt{uniform} & 32 & 31.00 & 1.0000 \\
\texttt{english} & 32 & 28.97 & 0.9692 \\
\texttt{structured} & 32 & 29.73 & 0.8902 \\
\texttt{periodic16} & 32 & 30.66 & 1.0273 \\
\bottomrule
\end{longtable}
\normalsize

On \texttt{uniform} it selects \(m-1\) at every needle length and ties Horspool to four decimal places. On C source it gives up two units of ceiling to reach a rarer byte and wins 11\%. This arm reads the whole corpus's frequencies before searching and that cost is not counted, which makes it a ceiling on anchor choice and not a usable rule.

Which of those two happens is decided in advance by one number. Modeling the needle as carrying the observed symbol at rate \(q = 2^{-H_2}\) per position, the distance to the nearest earlier occurrence has \(\Pr[d > j] = (1-q)^j\), so

\(\mathbb{E}[\text{advance}(a)] \;=\; \sum_{j=0}^{a}(1-q)^j \;=\; \frac{1 - (1-q)^{a+1}}{q}\)

This climbs and saturates at \(1/q = 2^{H_2}\), passing 95\% of that ceiling near \(a = 3 \cdot 2^{H_2}\). Below that radius the advance is still growing and the largest offset wins, and a rarer symbol cannot pay for the ceiling it costs. Above it the advance is flat and the rarity term decides alone.

The chosen offset over \(m-1\), against the radius, with the bar at the predicted crossing:

\tiny
\setlength{\tabcolsep}{2pt}
\begin{longtable}{@{}>{\setlength{\rightskip}{0pt}\setlength{\parfillskip}{0pt}\arraybackslash}p{\dimexpr\textwidth/8-2\tabcolsep\relax}>{\setlength{\rightskip}{0pt}\setlength{\parfillskip}{0pt}\arraybackslash}p{\dimexpr\textwidth/8-2\tabcolsep\relax}>{\setlength{\rightskip}{0pt}\setlength{\parfillskip}{0pt}\arraybackslash}p{\dimexpr\textwidth/8-2\tabcolsep\relax}>{\setlength{\rightskip}{0pt}\setlength{\parfillskip}{0pt}\arraybackslash}p{\dimexpr\textwidth/8-2\tabcolsep\relax}>{\setlength{\rightskip}{0pt}\setlength{\parfillskip}{0pt}\arraybackslash}p{\dimexpr\textwidth/8-2\tabcolsep\relax}>{\setlength{\rightskip}{0pt}\setlength{\parfillskip}{0pt}\arraybackslash}p{\dimexpr\textwidth/8-2\tabcolsep\relax}>{\setlength{\rightskip}{0pt}\setlength{\parfillskip}{0pt}\arraybackslash}p{\dimexpr\textwidth/8-2\tabcolsep\relax}>{\setlength{\rightskip}{0pt}\setlength{\parfillskip}{0pt}\arraybackslash}p{\dimexpr\textwidth/8-2\tabcolsep\relax}@{}}
\toprule
corpus & \(3\cdot 2^{H_2}\) & 8 & 16 & 32 & 64 & 128 & 256 \\
\midrule
\texttt{uniform} & 683.6 & 1.000 & 1.000 & 1.000 & 1.000 & 1.000 & 0.998 \\
\texttt{periodic16} & 73.3 & 1.000 & 0.996 & 0.989 & 0.974 & 0.888 & 0.770 \\
\texttt{english} & 38.8 & 0.993 & 0.975 & 0.935 & 0.911 & 0.869 & 0.774 \\
\texttt{structured} & 36.9 & 0.993 & 0.968 & 0.959 & 0.963 & 0.913 & 0.893 \\
\bottomrule
\end{longtable}
\normalsize

\texttt{uniform} has a radius larger than any needle tested. No needle reaches the flat zone and the optimum stays at \(m-1\) throughout, which it does to three decimals at five lengths of six. \texttt{periodic16} crosses between 64 and 128, and its chosen offset falls from 0.974 to 0.888 across exactly that gap. The two corpora with radii near 37 have already begun drifting while nominally below it, so \(3 \cdot 2^{H_2}\) is the right scale and not a sharp edge.

\paragraph{4.9.3 Online Advance Estimation}

A frequency table needs the alphabet to be finite and enumerable, which Section 2.1 does not assume. What stays finite when the alphabet does not is the advance, because an advance is a count of positions. So the field below holds one running advance per offset, \(m\) numbers, and the symbol side of the problem never appears.

Two properties come from the construction and not from tuning. Each offset starts at \(a+1\), the advance it can be proved to reach. The field opens at the last offset, which is Horspool, and can only be revised downward as answers arrive. And one answer updates every offset, because the advance another offset would have taken from the same cell is fixed by where the needle carries the observed symbol, which one backward pass over the needle computes without any table over symbols.

\tiny
\setlength{\tabcolsep}{2pt}
\begin{longtable}{@{}>{\setlength{\rightskip}{0pt}\setlength{\parfillskip}{0pt}\arraybackslash}p{\dimexpr\textwidth/5-2\tabcolsep\relax}>{\setlength{\rightskip}{0pt}\setlength{\parfillskip}{0pt}\arraybackslash}p{\dimexpr\textwidth/5-2\tabcolsep\relax}>{\setlength{\rightskip}{0pt}\setlength{\parfillskip}{0pt}\arraybackslash}p{\dimexpr\textwidth/5-2\tabcolsep\relax}>{\setlength{\rightskip}{0pt}\setlength{\parfillskip}{0pt}\arraybackslash}p{\dimexpr\textwidth/5-2\tabcolsep\relax}>{\setlength{\rightskip}{0pt}\setlength{\parfillskip}{0pt}\arraybackslash}p{\dimexpr\textwidth/5-2\tabcolsep\relax}@{}}
\toprule
corpus & \(m\) & Horspool & product rule, given the histogram & field, given nothing \\
\midrule
\texttt{structured} & 32 & 136.6 & 121.6 & \textbf{118.5} \\
\texttt{structured} & 64 & 129.5 & 128.9 & \textbf{121.7} \\
\texttt{structured} & 4 & 445.2 & 445.2 & \textbf{426.0} \\
\texttt{structured} & 128 & 177.6 & 170.2 & \textbf{168.3} \\
\texttt{english} & 128 & 274.4 & 255.7 & 266.4 \\
\texttt{uniform} & all &  &  & 1.000 to 1.008 of Horspool \\
\bottomrule
\end{longtable}
\normalsize

The field beats Horspool on 15 of 28 rows, by 13.2\% at best, and beats the product rule on every \texttt{structured} row despite the product rule being given the corpus histogram it was never shown. The product rule derives advance through a global frequency model, and the field measures advance directly and forgets at a rate of 0.1. The field follows local structure a global histogram averages away.

Crediting only the offset that asked, instead of all of them, leaves the field at parity with Horspool (0.95 to 1.05). The gain is in the counterfactual credit and not in the adaptation.

That credit carries an assumption, and the one corpus where the field loses is the one that violates it. A read at offset \(a\) returns the symbol in cell \(s+a\), and the credit given to offset \(a'\) is the advance that symbol would have bought there. The cell offset \(a'\) would actually have read is \(s+a'\). The credit is sound only where the symbol distribution does not depend on position. Fixed width records make it depend on the column, and the two cells are different columns.

The field over Horspool, by needle length:

\tiny
\setlength{\tabcolsep}{2pt}
\begin{longtable}{@{}>{\setlength{\rightskip}{0pt}\setlength{\parfillskip}{0pt}\arraybackslash}p{\dimexpr\textwidth/8-2\tabcolsep\relax}>{\setlength{\rightskip}{0pt}\setlength{\parfillskip}{0pt}\arraybackslash}p{\dimexpr\textwidth/8-2\tabcolsep\relax}>{\setlength{\rightskip}{0pt}\setlength{\parfillskip}{0pt}\arraybackslash}p{\dimexpr\textwidth/8-2\tabcolsep\relax}>{\setlength{\rightskip}{0pt}\setlength{\parfillskip}{0pt}\arraybackslash}p{\dimexpr\textwidth/8-2\tabcolsep\relax}>{\setlength{\rightskip}{0pt}\setlength{\parfillskip}{0pt}\arraybackslash}p{\dimexpr\textwidth/8-2\tabcolsep\relax}>{\setlength{\rightskip}{0pt}\setlength{\parfillskip}{0pt}\arraybackslash}p{\dimexpr\textwidth/8-2\tabcolsep\relax}>{\setlength{\rightskip}{0pt}\setlength{\parfillskip}{0pt}\arraybackslash}p{\dimexpr\textwidth/8-2\tabcolsep\relax}>{\setlength{\rightskip}{0pt}\setlength{\parfillskip}{0pt}\arraybackslash}p{\dimexpr\textwidth/8-2\tabcolsep\relax}@{}}
\toprule
corpus & 4 & 8 & 16 & 32 & 64 & 128 & 256 \\
\midrule
\texttt{periodic16} & 1.000 & 1.003 & \textbf{1.050} & \textbf{1.056} & 1.024 & 1.006 & 0.953 \\
\texttt{structured} & 0.957 & 0.946 & 0.908 & 0.868 & 0.940 & 0.947 & 0.992 \\
\texttt{english} & 1.002 & 1.003 & 0.998 & 0.994 & 0.985 & 0.971 & 0.981 \\
\texttt{uniform} & 1.000 & 1.000 & 1.004 & 1.005 & 1.008 & 1.004 & 1.000 \\
\bottomrule
\end{longtable}
\normalsize

\texttt{periodic16} is the only corpus above one across a run of lengths, and the loss peaks at 16 and 32, the record period and twice it, falling away on either side. The three corpora with no column structure show nothing. So the mechanism that buys 13.2\% and the mechanism that costs 5.6\% are the same one, and what separates the two cases is whether position carries information the credit ignores.

Reading the second cell instead of inferring it would remove the assumption at the price of a read. That read is the sample efficiency the counterfactual exists to buy. That trade is not measured here.

\paragraph{4.9.4 Ordered Harvest of Refutations}

Section 4.4.1 counts the alignments one read makes it possible to refute, and Section 4.9.1 counts the reads a greedy in-order shift performs. Those two are enough to price the ordering constraint without measuring anything further.

A read returning symbol \(c\) refutes every alignment whose needle position holds something else, which is \(m(1-2^{-H_2})\) of them in expectation. A shift that must proceed in order can only take the contiguous run up to the nearest alignment that survives, which stops it at the nearest occurrence of \(c\) in the needle. That symbol occurs about \(m\,2^{-H_2}\) times, placing the nearest one about \(m/(1+m\,2^{-H_2})\) away, so the fraction of the available refutation a greedy shift collects is

\(\text{harvest} \;\approx\; \frac{1}{\bigl(1 + m\,2^{-H_2}\bigr)\bigl(1 - 2^{-H_2}\bigr)}\)

Against the reads already reported, with no parameter fitted. The measured column is a lower bound, since a read count includes verification and therefore understates the mean advance:

\tiny
\setlength{\tabcolsep}{2pt}
\begin{longtable}{@{}>{\setlength{\rightskip}{0pt}\setlength{\parfillskip}{0pt}\arraybackslash}p{\dimexpr\textwidth/5-2\tabcolsep\relax}>{\setlength{\rightskip}{0pt}\setlength{\parfillskip}{0pt}\arraybackslash}p{\dimexpr\textwidth/5-2\tabcolsep\relax}>{\setlength{\rightskip}{0pt}\setlength{\parfillskip}{0pt}\arraybackslash}p{\dimexpr\textwidth/5-2\tabcolsep\relax}>{\setlength{\rightskip}{0pt}\setlength{\parfillskip}{0pt}\arraybackslash}p{\dimexpr\textwidth/5-2\tabcolsep\relax}>{\setlength{\rightskip}{0pt}\setlength{\parfillskip}{0pt}\arraybackslash}p{\dimexpr\textwidth/5-2\tabcolsep\relax}@{}}
\toprule
corpus & \(m\) & \(H_2\) & predicted & measured at least \\
\midrule
\texttt{structured} & 16 & 3.620 & 47.3\% & 47\% \\
\texttt{english} & 16 & 3.692 & 48.4\% & 52\% \\
\texttt{periodic16} & 16 & 4.611 & 63.0\% & 74\% \\
\texttt{uniform} & 16 & 7.832 & 93.9\% & 91\% \\
\texttt{english} & 64 & 3.692 & 18.2\% & 19\% \\
\texttt{english} & 256 & 3.692 & 5.2\% & 3\% \\
\bottomrule
\end{longtable}
\normalsize

The harvest is monotone in \(H_2\) at fixed \(m\), reading 47\%, 52\%, 74\% and 91\% across the four corpora in entropy order, and it collapses as \(m\) grows because a longer needle carries more copies of whatever was read.

\textbf{The consequence for a search that gives up the order.} For large \(m\) the harvest tends to \(2^{H_2}/m\), so the ceiling on what dropping the constraint can buy is \(m\,2^{-H_2}\), the expected number of times the observed symbol occurs in the needle. That is roughly 20 on English prose at \(m = 256\) and 1.07 on the uniform corpus at \(m = 16\). The regime where order costs something is low entropy and long needles, and there is nothing to collect at high entropy and short ones.

Every refutation a read makes available is exact and permanent by Proposition 1. None of this information is destroyed by being ignored. A greedy in-order shift declines to collect between 3\% and 99\% of it, and the table below gives the figure per row.

\paragraph{4.9.5 Unordered Probe Allocation}

Section 4.9.4 says how much refutation an in-order walk declines to collect. This one spends it.

If order is free then a read is not attached to an anchor at all. One read of one cell refutes every alignment whose needle position holds a different symbol, so there is nothing to choose except which cells to read, and those can be fixed in advance. Reading at stride \(k\) covers each alignment \(m/k\) times, so with \(x = m/k\) the survivors are \(N 2^{-H_2 x}\), the probe reads are \(Nx/m\), and the total is minimized at

\[
x \;=\; \frac{\log_2\!\left(m H_2 \ln 2\right)}{H_2},
\qquad \text{reads} \;\approx\; \frac{N}{m}\left[x + \frac{1}{H_2 \ln 2}\right]
\]

That model predicted savings of 11 to 26 times at \(m = 256\). It was wrong by an order of magnitude, and the way it was wrong is the useful part.

Measured, as in-order reads over free-order reads, so above one is a saving:

\tiny
\setlength{\tabcolsep}{2pt}
\begin{longtable}{@{}>{\setlength{\rightskip}{0pt}\setlength{\parfillskip}{0pt}\arraybackslash}p{\dimexpr\textwidth/6-2\tabcolsep\relax}>{\setlength{\rightskip}{0pt}\setlength{\parfillskip}{0pt}\arraybackslash}p{\dimexpr\textwidth/6-2\tabcolsep\relax}>{\setlength{\rightskip}{0pt}\setlength{\parfillskip}{0pt}\arraybackslash}p{\dimexpr\textwidth/6-2\tabcolsep\relax}>{\setlength{\rightskip}{0pt}\setlength{\parfillskip}{0pt}\arraybackslash}p{\dimexpr\textwidth/6-2\tabcolsep\relax}>{\setlength{\rightskip}{0pt}\setlength{\parfillskip}{0pt}\arraybackslash}p{\dimexpr\textwidth/6-2\tabcolsep\relax}>{\setlength{\rightskip}{0pt}\setlength{\parfillskip}{0pt}\arraybackslash}p{\dimexpr\textwidth/6-2\tabcolsep\relax}@{}}
\toprule
corpus & \(m = 4\) & \(m = 16\) & \(m = 64\) & \(m = 128\) & \(m = 256\) \\
\midrule
\texttt{english} & 0.96 & 0.85 & 1.38 & \textbf{1.47} & 1.33 \\
\texttt{structured} & 1.07 & 0.81 & 1.13 & 1.16 & 1.08 \\
\texttt{periodic16} & 0.63 & 0.76 & 0.99 & 1.20 & 1.36 \\
\texttt{uniform} & 0.44 & 0.27 & 0.93 & 0.94 & 0.98 \\
\texttt{mixed} & 0.66 & 0.82 & 0.80 & 0.93 & 1.05 \\
\bottomrule
\end{longtable}
\normalsize

The dependency depth is 2 on every row, as the construction requires.

\textbf{The probe phase behaved exactly as modeled and the model still failed.} Survivors after probing at \(m = 256\) are 1, 7, 8, 11 and 2 across the five corpora, against a prediction of 1.25 to 3.84. The reads went somewhere the model never counted.

Every needle here is drawn from the corpus it is searched in. Every search must confirm one genuine occurrence, and confirming a match of length \(m\) costs \(m\) reads. That term is absent from the derivation above. Adding it accounts for the measurements to within 3\%:

\tiny
\setlength{\tabcolsep}{2pt}
\begin{longtable}{@{}>{\setlength{\rightskip}{0pt}\setlength{\parfillskip}{0pt}\arraybackslash}p{\dimexpr\textwidth/6-2\tabcolsep\relax}>{\setlength{\rightskip}{0pt}\setlength{\parfillskip}{0pt}\arraybackslash}p{\dimexpr\textwidth/6-2\tabcolsep\relax}>{\setlength{\rightskip}{0pt}\setlength{\parfillskip}{0pt}\arraybackslash}p{\dimexpr\textwidth/6-2\tabcolsep\relax}>{\setlength{\rightskip}{0pt}\setlength{\parfillskip}{0pt}\arraybackslash}p{\dimexpr\textwidth/6-2\tabcolsep\relax}>{\setlength{\rightskip}{0pt}\setlength{\parfillskip}{0pt}\arraybackslash}p{\dimexpr\textwidth/6-2\tabcolsep\relax}>{\setlength{\rightskip}{0pt}\setlength{\parfillskip}{0pt}\arraybackslash}p{\dimexpr\textwidth/6-2\tabcolsep\relax}@{}}
\toprule
corpus & probes & \(+\,m\) & \(+\) false survivors & predicted & measured \\
\midrule
\texttt{english} & 25.1 & 256 & 9.0 & 290.1 & 289.7 \\
\texttt{structured} & 9.7 & 256 & 0.0 & 265.7 & 267.0 \\
\texttt{periodic16} & 31.5 & 256 & 10.5 & 298.0 & 296.0 \\
\texttt{uniform} & 19.8 & 256 & 15.0 & 290.8 & 287.6 \\
\bottomrule
\end{longtable}
\normalsize

\textbf{A long search is verification bound and not filter bound.} The floor is \(m\) reads for any algorithm that must confirm the match, and at \(m = 256\) Horspool runs at 1.10 to 1.57 times that floor while free order runs at 1.04 to 1.16. Both are already near a bound that no filtering strategy can move, so an order of magnitude was never available. The saving that does exist is the gap between those two ratios, and 1.47 at \(m = 128\) on English is the largest of it measured.

At short needles free order loses, badly on the uniform corpus at 0.27, because probing at a stride below the needle length costs more reads than a shift does and there is no verification floor to hide behind.

\textbf{The accumulation is bitwise, and that makes the count reachable.} A refutation set is a set of alignments, sets union commutatively by Proposition 1, and a set of alignments is a bit vector. So accumulating a read costs \(\lceil m/W \rceil\) word operations for a machine word of \(W\) bits, the surviving set is the complement, and no other structure is needed. At \(m = 256\) with 64 bit words the whole probe phase is 25 reads and 100 word operations, at a dependency depth of two, against 385 strictly sequential reads.

The mask table is indexed by pattern position, which makes it \(m\) bits wide whatever the alphabet is and whatever the dimension is. Section 2.1 can therefore decline to bound either one at no cost: neither appears in the state.

\paragraph{4.9.6 Verification Necessity}

Section 4.9.5 finds a long search verification bound. The obvious question is whether the verification can be dropped. Proposition 1 says every occurrence survives every anchor set, and a survivor set the same size as the occurrence set contains exactly the occurrences and confirming distinguishes nothing. The arm below probes at the stride that leaves one survivor in expectation, declares the survivors, and never confirms.

The arm is cheap and it is wrong.

\tiny
\setlength{\tabcolsep}{2pt}
\begin{longtable}{@{}>{\setlength{\rightskip}{0pt}\setlength{\parfillskip}{0pt}\arraybackslash}p{\dimexpr\textwidth/6-2\tabcolsep\relax}>{\setlength{\rightskip}{0pt}\setlength{\parfillskip}{0pt}\arraybackslash}p{\dimexpr\textwidth/6-2\tabcolsep\relax}>{\setlength{\rightskip}{0pt}\setlength{\parfillskip}{0pt}\arraybackslash}p{\dimexpr\textwidth/6-2\tabcolsep\relax}>{\setlength{\rightskip}{0pt}\setlength{\parfillskip}{0pt}\arraybackslash}p{\dimexpr\textwidth/6-2\tabcolsep\relax}>{\setlength{\rightskip}{0pt}\setlength{\parfillskip}{0pt}\arraybackslash}p{\dimexpr\textwidth/6-2\tabcolsep\relax}>{\setlength{\rightskip}{0pt}\setlength{\parfillskip}{0pt}\arraybackslash}p{\dimexpr\textwidth/6-2\tabcolsep\relax}@{}}
\toprule
corpus & \(m\) & in-order reads & unconfirmed reads & saving & searches wrong, of 64 \\
\midrule
\texttt{structured} & 256 & 287.5 & 14.0 & 20.5 & 51 \\
\texttt{english} & 256 & 384.9 & 34.0 & 11.3 & 32 \\
\texttt{uniform} & 256 & 281.6 & 25.0 & 11.3 & 64 \\
\texttt{structured} & 128 & 177.6 & 27.0 & 6.6 & 49 \\
\texttt{english} & 4 & 916.5 & 2789.0 & 0.33 & 0 \\
\bottomrule
\end{longtable}
\normalsize

The last row is the only shape that is ever right: at stride one every cell is read, nothing false can survive, and the cost is worse than reading the corpus straight through. Everywhere else the saving is real and the answer is wrong on half to all of the searches.

\textbf{The two propositions appear here as measurements.} Across all 35 rows the count of searches that over-report equals the count that are wrong. Not one search in any corpus, at any needle length, at any stride, ever reported fewer occurrences than exist. Proposition 1 has been an argument in this document until now, and this is the first place it is visible in data. Proposition 2 is visible in the same column: survivors are a strict superset and probing does not close the gap short of reading everything.

\textbf{The error is one directional, which makes it a signal.} Since a miss is impossible, any discrepancy is an over-count. An over-count is detectable without knowing the answer and means exactly one thing: the stride is too coarse. The transition is sharp. On English at \(m = 4\) the stride moves from 1 to 2 and the searches in error move from 0 to 61 of 64.

\textbf{Calibrating the stride on one known pattern does not make it safe for others.} A signal that is one directional and monotone in the stride can be bisected. One needle whose occurrence count is known can find the largest stride reproducing that count, in \(\log_2 m\) passes. Applying that stride to the other 63 needles:

\tiny
\setlength{\tabcolsep}{2pt}
\begin{longtable}{@{}>{\setlength{\rightskip}{0pt}\setlength{\parfillskip}{0pt}\arraybackslash}p{\dimexpr\textwidth/6-2\tabcolsep\relax}>{\setlength{\rightskip}{0pt}\setlength{\parfillskip}{0pt}\arraybackslash}p{\dimexpr\textwidth/6-2\tabcolsep\relax}>{\setlength{\rightskip}{0pt}\setlength{\parfillskip}{0pt}\arraybackslash}p{\dimexpr\textwidth/6-2\tabcolsep\relax}>{\setlength{\rightskip}{0pt}\setlength{\parfillskip}{0pt}\arraybackslash}p{\dimexpr\textwidth/6-2\tabcolsep\relax}>{\setlength{\rightskip}{0pt}\setlength{\parfillskip}{0pt}\arraybackslash}p{\dimexpr\textwidth/6-2\tabcolsep\relax}>{\setlength{\rightskip}{0pt}\setlength{\parfillskip}{0pt}\arraybackslash}p{\dimexpr\textwidth/6-2\tabcolsep\relax}@{}}
\toprule
corpus & \(m\) & calibrated & derived & wrong of 63 & over in-order reads \\
\midrule
\texttt{structured} & 256 & 98 & 91 & 23 & 22.1 \\
\texttt{english} & 256 & 75 & 84 & 28 & 10.1 \\
\texttt{uniform} & 256 & 145 & 170 & 59 & 9.7 \\
\texttt{periodic16} & 256 & 66 & 99 & 1 & 6.4 \\
\texttt{uniform} & 16 & 9 & 11 & 62 & 0.6 \\
\texttt{english} & 8 & 1 & 3 & 0 & 0.19 \\
\bottomrule
\end{longtable}
\normalsize

Every row with no errors is a row where calibration drove the stride to one or two, which reads nearly every cell and costs more than an in-order search. Every row that is fast is also wrong on a third to almost all of the unseen needles.

The calibrated stride and the derived stride agree closely, 98 against 91 and 75 against 84 and 44 against 45, so the failure is not a mis-set constant. The safe stride is a property of the pattern and the corpus together: a pattern carrying rare symbols tolerates a long stride and one carrying common symbols does not, and a measurement on one pattern bounds nothing about another. One known constant calibrates one pattern.

\textbf{What the failure was, since it was predicted to be something else.} This section originally reported the model's 11 to 26 times as its result, flagged as resting on the independence assumption that Section 4.6 measures failing worst on \texttt{structured}. That flag was aimed at the wrong thing. The independence assumption held: predicted survivors 1.25 to 3.84 against measured 1 to 11. What the model omitted was the \(m\) reads needed to confirm the occurrence every needle is guaranteed to have, a term with nothing probabilistic in it.

A model can be wrong in the place it was not being watched. The caveat attached to this section named the weakest assumption and the error was in an assumption not stated at all.

\paragraph{4.9.7 Adaptive Probe Selection}

Section 7.1 separates two arrangements. A fixed probe set can be made safe by deriving it from the pattern in advance. Deterministic sampling does that. A fixed probe set without that derivation is what Sections 4.9.5 and 4.9.6 measure failing. The third arrangement holds no probe set at all.

The rule needs no pattern analysis and no distribution. The symbol about to be read is unknown. Every candidate position kills the same expected fraction of the alignments it touches, and what differs between positions is how many living alignments touch them at all. So the next read goes where the survivors overlap most, which is a function of the answers so far and of nothing else. The rule stops when a read separates nothing and confirms what remains (\texttt{test/\allowbreak{}bench/\allowbreak{}bench\_\allowbreak{}ancorae\_\allowbreak{}ab.\allowbreak{}c:1122}).

\tiny
\setlength{\tabcolsep}{2pt}
\begin{longtable}{@{}>{\setlength{\rightskip}{0pt}\setlength{\parfillskip}{0pt}\arraybackslash}p{\dimexpr\textwidth/5-2\tabcolsep\relax}>{\setlength{\rightskip}{0pt}\setlength{\parfillskip}{0pt}\arraybackslash}p{\dimexpr\textwidth/5-2\tabcolsep\relax}>{\setlength{\rightskip}{0pt}\setlength{\parfillskip}{0pt}\arraybackslash}p{\dimexpr\textwidth/5-2\tabcolsep\relax}>{\setlength{\rightskip}{0pt}\setlength{\parfillskip}{0pt}\arraybackslash}p{\dimexpr\textwidth/5-2\tabcolsep\relax}>{\setlength{\rightskip}{0pt}\setlength{\parfillskip}{0pt}\arraybackslash}p{\dimexpr\textwidth/5-2\tabcolsep\relax}@{}}
\toprule
corpus & reads at \(m = 256\) & probes among them & reads over the floor & in-order over the floor \\
\midrule
\texttt{structured} & 262.1 & 9.5 & 1.02 & 1.12 \\
\texttt{uniform} & 277.7 & 23.5 & 1.08 & 1.10 \\
\texttt{english} & 279.4 & 26.6 & 1.09 & 1.50 \\
\texttt{mixed} & 281.8 & 28.7 & 1.10 & 1.21 \\
\texttt{periodic16} & 307.4 & 45.4 & 1.20 & 1.57 \\
\bottomrule
\end{longtable}
\normalsize

The floor is the \(m\) reads Section 4.9.5 identifies as irreducible. This runs within 2\% to 20\% of it where an in-order search runs within 10\% to 57\%, and it is exact on all 35 rows with a mirror residual of zero. The rule reads fewer cells than an in-order search on 22 of the 35, by up to 1.57 times at \(m = 64\) and 128.

\textbf{The selecting phase is small without being made small.} Nine probes settle \texttt{structured} at \(m = 256\), against \(\log_2 256 = 8\), the size deterministic sampling reaches by analyzing the pattern first. English and uniform take 23 to 27. Nothing here derives a sample, carries state between searches, or looks at the pattern before the first read.

\textbf{Two limits.} The logarithmic size in the published method is a worst case guarantee and this is a mean over 64 needles on five corpora with no guarantee of any kind. And it loses on \texttt{periodic16} at needle lengths 8 through 32, between 0.64 and 0.86. That is the positional structure Sections 4.7 and 4.9.3 fail on for the same reason.

\paragraph{4.9.8 Negative Control Mechanisms}

Recorded because each was built, measured on the case it was designed for, and kept. Two of them are control laws over the field of Section 4.9.3, where the process variable is an advance, which is a small integer with a variance close to its mean.

\textbf{Discarding the field on divergence.} A short window mean of the advance against a long window mean detects a corpus whose character has changed, and every offset returns to its provable ceiling, which restores the anchor to \(m-1\). Over the four homogeneous corpora the detector never fires on 15 of 28 rows and moves the result by at most 1\% where it does. A fifth corpus of four regions laid end to end, English then C source then records then uniform bytes, was built for it, and every arm on that corpus converges to within 1\% of every other. The mechanism has no measured benefit and no measured cost.

\textbf{Pulling the anchor to the field's center of mass.} Selecting the offset at \(\sum_i v_i i / \sum_i v_i\) instead of the offset with the largest \(v_i\) loses on every row, by 2\% to 42\% against the largest and by up to 48\% against Horspool. The field is roughly increasing in offset because each entry starts at \(a+1\). Its balance point sits near two thirds of the way along and carries a mediocre ceiling. The centroid is a canonical location and the mean of a utility is not the argument that maximizes it, so the construction that gives Section 2.1 its reference point is the wrong one for choosing among values.

\textbf{Adding accumulated and trend terms to the update.} The field's update is a first order lag, which is a proportional response with a leak. Adding an accumulated term at 0.002 and a trend term at 0.05 gives a mean of 1.005 against the lag alone and is worse on 24 of 35 rows, with the largest losses on the periodic corpus where the trend term follows the oscillation. An advance is a small integer whose variance is close to its mean, and a difference between consecutive errors is mostly noise and a term proportional to it mostly injects noise.

\textbf{Rejecting answers at the ambient level and amplifying the rest.} Replacing the lag with a deadband against the running mean advance, full weight applied outside it and settling back to the provable ceiling inside it, gives a mean of 1.014 against the lag and is better on 2 of 35 rows.

\textbf{What the five have in common.} Each changes how the search responds to information it already has. The one change in Section 4.9.3 that did pay, crediting every offset from one answer instead of only the offset that asked, changes how much information one observation yields. Over the seven arrangements measured here, extracting more per observation is worth 13.2\% and responding to it more cleverly is worth nothing.

\subsubsection{4.10 Symbol Width and Entropy Rate}

Section 2.1 declines to assume the alphabet is enumerable. It does not follow that a symbol is a byte, and nothing in this document justifies eight bits. A corpus is a bit pattern, and where one symbol ends and the next begins is a slice imposed on it. Every \(H_2\) reported before this section is a property of the corpus and that slice together.

Proposition 1 does not mention the width. Soundness holds at every width at once. What the width moves is cost. A read of \(w\) bits costs \(w\) bits of reading and yields \(H_2(w)\) bits of discrimination, so what a search spends is \(w \log_2 N / H_2(w)\) and the quantity to compare across widths is \(H_2(w)/w\).

Measured over windows taken at every bit offset, since a read can begin anywhere (\texttt{test/\allowbreak{}bench/\allowbreak{}bench\_\allowbreak{}ancorae\_\allowbreak{}sift.\allowbreak{}c:1122}):

\tiny
\setlength{\tabcolsep}{2pt}
\begin{longtable}{@{}>{\setlength{\rightskip}{0pt}\setlength{\parfillskip}{0pt}\arraybackslash}p{\dimexpr\textwidth/7-2\tabcolsep\relax}>{\setlength{\rightskip}{0pt}\setlength{\parfillskip}{0pt}\arraybackslash}p{\dimexpr\textwidth/7-2\tabcolsep\relax}>{\setlength{\rightskip}{0pt}\setlength{\parfillskip}{0pt}\arraybackslash}p{\dimexpr\textwidth/7-2\tabcolsep\relax}>{\setlength{\rightskip}{0pt}\setlength{\parfillskip}{0pt}\arraybackslash}p{\dimexpr\textwidth/7-2\tabcolsep\relax}>{\setlength{\rightskip}{0pt}\setlength{\parfillskip}{0pt}\arraybackslash}p{\dimexpr\textwidth/7-2\tabcolsep\relax}>{\setlength{\rightskip}{0pt}\setlength{\parfillskip}{0pt}\arraybackslash}p{\dimexpr\textwidth/7-2\tabcolsep\relax}>{\setlength{\rightskip}{0pt}\setlength{\parfillskip}{0pt}\arraybackslash}p{\dimexpr\textwidth/7-2\tabcolsep\relax}@{}}
\toprule
corpus & \(w{=}1\) & \(w{=}2\) & \(w{=}4\) & \(w{=}8\) & \(w{=}12\) & \(w{=}16\) \\
\midrule
\texttt{english} & 0.984 & 0.984 & 0.969 & 0.856 & 0.696 & 0.600 \\
\texttt{structured} & 0.967 & 0.966 & 0.954 & 0.850 & 0.670 & 0.541 \\
\texttt{periodic16} & 0.989 & 0.988 & 0.944 & 0.888 & 0.792 & 0.726 \\
\texttt{uniform} & 1.000 & 1.000 & 1.000 & 0.997 & 0.973 & 0.855 \\
\texttt{flat} & 0.678 & 0.708 & 0.670 & 0.375 & 0.250 & 0.188 \\
\bottomrule
\end{longtable}
\normalsize

The ratio falls with width on every corpus. Reading English a byte at a time yields 0.856 bits per bit read against 0.984 at a single bit. The byte slice gives up 15\% of the discrimination available in the same number of bits, and a sixteen bit symbol gives up 40\%.

\textbf{The flat corpus is not flat.} It is \texttt{0x41} repeated, \texttt{01000001} has period eight, and a window of six bits or more takes exactly eight distinct values whatever the corpus length. Its collision entropy is therefore \(\log_2 8 = 3\) bits exactly. The sweep reports that at widths 6, 8, 12 and 16, and 0.678 bits at width one. The object every other section of this document reports zero on carries three bits under a slice one bit narrower.

\textbf{What this does not settle.} \(H_2(w)/w\) is the right comparison only where reading costs by the bit. Where a read costs the same whatever its width, as it does for a machine word, the quantity to maximize is \(H_2(w)\) itself, and that grows monotonically with \(w\) on every row above. The three cost models disagree about which width is best, this document counts reads under the third of them, and none of the earlier sections states which one it is assuming.

\subsubsection{4.11 Period Detection Without a Pattern}

Every measurement before this one is given a needle. The construction does not require one. A shift \(d\) is the hypothesis that the corpus agrees with itself \(d\) apart, a pair of read cells whose positions differ by \(d\) is a test of it, and one disagreement refutes it permanently. So the pattern is not one string, it is every shift at once, and the reads are shared across all of them: \(k\) cells carry \(k(k-1)/2\) tests. Here 512 reads test 130,816 pairs (\texttt{test/\allowbreak{}bench/\allowbreak{}bench\_\allowbreak{}ancorae\_\allowbreak{}ab.\allowbreak{}c:1315}).

\textbf{The threshold cannot be trusted and is not.} Tests of one shift share read positions and are not independent, and a standard error computed as though they were is too small. The first run of this reported twelve shifts on the uniform corpus, which is SHA-256 output and has nothing to find.

A shuffle of the same bytes settles it. It preserves the byte histogram exactly, so the collision probability is unchanged, and it destroys every relation between positions, and that is all a shift can be. Whatever the detector reports there is what it reports on nothing.

\tiny
\setlength{\tabcolsep}{2pt}
\begin{longtable}{@{}>{\setlength{\rightskip}{0pt}\setlength{\parfillskip}{0pt}\arraybackslash}p{\dimexpr\textwidth/5-2\tabcolsep\relax}>{\setlength{\rightskip}{0pt}\setlength{\parfillskip}{0pt}\arraybackslash}p{\dimexpr\textwidth/5-2\tabcolsep\relax}>{\setlength{\rightskip}{0pt}\setlength{\parfillskip}{0pt}\arraybackslash}p{\dimexpr\textwidth/5-2\tabcolsep\relax}>{\setlength{\rightskip}{0pt}\setlength{\parfillskip}{0pt}\arraybackslash}p{\dimexpr\textwidth/5-2\tabcolsep\relax}>{\setlength{\rightskip}{0pt}\setlength{\parfillskip}{0pt}\arraybackslash}p{\dimexpr\textwidth/5-2\tabcolsep\relax}@{}}
\toprule
corpus & shifts reported & on the shuffle & difference & strongest shift \\
\midrule
\texttt{periodic16} & 92 & 0 & \textbf{92} & 560 \\
\texttt{uniform} & 12 & 16 & -4 & 919 \\
\texttt{mixed} & 4 & 6 & -2 & 416 \\
\texttt{structured} & 1 & 1 & 0 & 984 \\
\texttt{english} & 0 & 0 & 0 & 1719 \\
\bottomrule
\end{longtable}
\normalsize

One corpus has structure of this kind and it is the one built with a period. Its strongest shift is 560, and \(560 = 35 \times 16\). The detector recovered a multiple of the record length without being given the record length or a pattern to look for. The \texttt{mixed} corpus peaks at 416, also a multiple of 16, which is its periodic quarter showing through below the threshold.

\textbf{What it does not find is as informative.} English prose and C source report nothing above their shuffles, and English peaks lower than its own shuffle does. Both plainly have structure. What they do not have is agreement with themselves at a fixed offset. That is the only hypothesis this formulation tests. Repetition at irregular positions, which is what a word or an identifier is, needs a different query. The detector finds periodicity and is silent on everything else.

\subsubsection{4.12 Boundary Detection}

Section 4.11 asks whether a corpus repeats at a fixed offset, and prose does not. A different question reaches it. Whatever else a language is, it composes units, so something in it marks where one unit ends, by a separator or by a rule. If nothing does then nothing is a unit. That has a signature the same reads can measure.

If a symbol's positions carry no structure, the gaps between its occurrences are geometric, so the ratio of the variance of the gaps to their squared mean sits near one. A separator has bounded unit length. Its gaps are far tighter and the ratio falls. The statistic has no units and no scale, so one threshold reads the same on any alphabet and any corpus (\texttt{test/\allowbreak{}bench/\allowbreak{}bench\_\allowbreak{}ancorae\_\allowbreak{}ab.\allowbreak{}c:1482}). The shuffle calibrates it as before: it holds every symbol's count and destroys where they sit, which makes every symbol's gaps geometric by construction.

\tiny
\setlength{\tabcolsep}{2pt}
\begin{longtable}{@{}>{\setlength{\rightskip}{0pt}\setlength{\parfillskip}{0pt}\arraybackslash}p{\dimexpr\textwidth/6-2\tabcolsep\relax}>{\setlength{\rightskip}{0pt}\setlength{\parfillskip}{0pt}\arraybackslash}p{\dimexpr\textwidth/6-2\tabcolsep\relax}>{\setlength{\rightskip}{0pt}\setlength{\parfillskip}{0pt}\arraybackslash}p{\dimexpr\textwidth/6-2\tabcolsep\relax}>{\setlength{\rightskip}{0pt}\setlength{\parfillskip}{0pt}\arraybackslash}p{\dimexpr\textwidth/6-2\tabcolsep\relax}>{\setlength{\rightskip}{0pt}\setlength{\parfillskip}{0pt}\arraybackslash}p{\dimexpr\textwidth/6-2\tabcolsep\relax}>{\setlength{\rightskip}{0pt}\setlength{\parfillskip}{0pt}\arraybackslash}p{\dimexpr\textwidth/6-2\tabcolsep\relax}@{}}
\toprule
corpus & symbol found & mean gap & dispersion & on the shuffle & ratio \\
\midrule
\texttt{english} & \texttt{0x20}, space & 5.58 & 0.162 & 0.599 & 3.70 \\
\texttt{periodic16} & \texttt{0x0A}, newline & 16.00 & 0.0000 & 0.653 & unbounded \\
\texttt{structured} & \texttt{0x3B}, semicolon & 59.9 & 0.353 & 0.413 & 1.17 \\
\texttt{uniform} & \texttt{0x45} & 177.6 & 0.324 & 0.339 & 1.05 \\
\texttt{mixed} & \texttt{0x6B} & 134.0 & 0.814 & 0.400 & 0.49 \\
\bottomrule
\end{longtable}
\normalsize

Each corpus that has a boundary marker returns it. English returns the space, at a mean gap of 5.58, which is a word and its separator. The record corpus returns the newline at exactly 16.00 with a dispersion of exactly zero. C source returns the semicolon, which is its statement terminator, at a weak 1.17, because statement lengths vary far more than word lengths do.

The uniform corpus returns 1.05, which is nothing, and the mixed corpus returns 0.49, which is less regular than its own shuffle. Four languages laid end to end have no single marker that is regular across all of them, and the statistic says so instead of reporting the marker of whichever one is longest.

Nothing here is given a pattern, a dictionary, or a notion of what a unit is. What makes the comparison sound is that both objects are finite and hold the same symbols, so the null is constructed and not assumed.

\paragraph{4.12.1 Boundary Spacing as a Filter}

A known boundary symbol turns a corpus into a sequence of gaps, and a needle carrying \(k\) boundaries carries \(k-1\) gaps between them that any true occurrence has to reproduce. That run is a filter on its own, before any other symbol is looked at (\texttt{test/\allowbreak{}bench/\allowbreak{}bench\_\allowbreak{}ancorae\_\allowbreak{}ab.\allowbreak{}c:1605}).

\tiny
\setlength{\tabcolsep}{2pt}
\begin{longtable}{@{}>{\setlength{\rightskip}{0pt}\setlength{\parfillskip}{0pt}\arraybackslash}p{\dimexpr\textwidth/5-2\tabcolsep\relax}>{\setlength{\rightskip}{0pt}\setlength{\parfillskip}{0pt}\arraybackslash}p{\dimexpr\textwidth/5-2\tabcolsep\relax}>{\setlength{\rightskip}{0pt}\setlength{\parfillskip}{0pt}\arraybackslash}p{\dimexpr\textwidth/5-2\tabcolsep\relax}>{\setlength{\rightskip}{0pt}\setlength{\parfillskip}{0pt}\arraybackslash}p{\dimexpr\textwidth/5-2\tabcolsep\relax}>{\setlength{\rightskip}{0pt}\setlength{\parfillskip}{0pt}\arraybackslash}p{\dimexpr\textwidth/5-2\tabcolsep\relax}@{}}
\toprule
corpus & \(m\) & boundaries in the needle & survivors & reduction \\
\midrule
\texttt{english} & 16 & 2.88 & 37.3 & 74 \\
\texttt{english} & 32 & 5.72 & 1.75 & 1572 \\
\texttt{english} & 64 & 11.45 & 1.00 & 2726 \\
\texttt{english} & 256 & 46.0 & 1.00 & 2534 \\
\texttt{structured} & 256 & 4.27 & 1.01 & 923 \\
\texttt{periodic16} & 64 & 4.00 & 252.1 & 16.0 \\
\texttt{periodic16} & 256 & 16.00 & 240.3 & 16.0 \\
\bottomrule
\end{longtable}
\normalsize

On English the spacing alone locates the needle uniquely from a needle length of 64 upward, using a symbol that is 17.9\% of the corpus and is therefore the worst anchor in it by rarity. The filter strengthens as the needle lengthens because a longer needle carries more boundaries, and the boundary count is the property being tested.

\textbf{The record corpus is the counterexample and it does not move.} Its reduction is exactly 16.0 at every needle length and at every boundary count, because its newlines are perfectly periodic: the spacing fixes the alignment modulo 16 and says nothing further. A boundary of period \(T\) carries \(\log_2 T\) bits and no more, however many of them a needle holds.

That inverts Section 4.12. \texttt{periodic16} reported the tightest spacing found anywhere, a dispersion of exactly zero, and its boundaries are the least useful for locating anything. English reported 0.162 and its boundaries locate uniquely. A boundary has to be regular enough to be found and irregular enough to carry a position, and the variation in the spacing is the part that does the second job.

The uniform corpus also reduces to a single survivor, at a marker rate of 0.0049. That is the rare anchor of Section 4.3 working normally and not a boundary. A corpus with no boundaries should produce that.

\subsubsection{4.13 Cross-Linguistic Regularities}

Section 4.12 finds a boundary and Section 4.12.1 uses it. Neither says whether what it found is a property of language or of the one English sample it was run on, which was written by the author of this document. Two regularities are claimed to hold of every natural language and of nothing else: the frequency of a unit falls as roughly the reciprocal of its rank, and the frequent units are the short ones. Both are consequences of communicating under a fixed cognitive budget. Neither should care what the language is or when it was written.

Six public domain texts, fetched by \texttt{maint/\allowbreak{}data/\allowbreak{}fetch/\allowbreak{}fetch\_\allowbreak{}language\_\allowbreak{}corpora.\allowbreak{}py} and read by \texttt{bench\_\allowbreak{}ancorae\_\allowbreak{}ab.\allowbreak{}c} given a path. Line endings are folded into the space first, because a plain text file wrapped at a fixed width carries a return every fifty odd bytes and the detector finds the publisher's formatting before it finds the language.

\tiny
\setlength{\tabcolsep}{2pt}
\begin{longtable}{@{}>{\setlength{\rightskip}{0pt}\setlength{\parfillskip}{0pt}\arraybackslash}p{\dimexpr\textwidth/7-2\tabcolsep\relax}>{\setlength{\rightskip}{0pt}\setlength{\parfillskip}{0pt}\arraybackslash}p{\dimexpr\textwidth/7-2\tabcolsep\relax}>{\setlength{\rightskip}{0pt}\setlength{\parfillskip}{0pt}\arraybackslash}p{\dimexpr\textwidth/7-2\tabcolsep\relax}>{\setlength{\rightskip}{0pt}\setlength{\parfillskip}{0pt}\arraybackslash}p{\dimexpr\textwidth/7-2\tabcolsep\relax}>{\setlength{\rightskip}{0pt}\setlength{\parfillskip}{0pt}\arraybackslash}p{\dimexpr\textwidth/7-2\tabcolsep\relax}>{\setlength{\rightskip}{0pt}\setlength{\parfillskip}{0pt}\arraybackslash}p{\dimexpr\textwidth/7-2\tabcolsep\relax}>{\setlength{\rightskip}{0pt}\setlength{\parfillskip}{0pt}\arraybackslash}p{\dimexpr\textwidth/7-2\tabcolsep\relax}@{}}
\toprule
text & family & year & boundary & gap & Zipf slope & brevity \\
\midrule
Cervantes, \_Don Quixote\_ & Romance & 1605 & \texttt{0x20} & 5.13 & -1.047 & -0.305 \\
King James Bible & Germanic & 1611 & \texttt{0x20} & 4.70 & -0.925 & -0.244 \\
Goethe, \_Faust\_ & Germanic & 1808 & \texttt{0x20} & 5.09 & -0.723 & -0.364 \\
Austen, \_Pride and Prejudice\_ & Germanic & 1813 & \texttt{0x20} & 4.94 & -0.875 & -0.301 \\
Hugo, \_Les Misérables\_ & Romance & 1862 & \texttt{0x20} & 5.40 & -0.950 & -0.321 \\
\_Kalevala\_ & Uralic & 1849 & \texttt{0x2C} & 43.7 & -0.563 & -0.153 \\
\bottomrule
\end{longtable}
\normalsize

Five of the six locate the same boundary symbol without being told there is one, at a mean gap between 4.70 and 5.40. The Zipf slopes sit between -0.72 and -1.05 against a stated -1, and the brevity correlation is negative in every row. Cervantes and Hugo are 257 years and two countries apart and differ by 0.10 in slope and 0.016 in brevity.

\textbf{The Finnish row is a genre confound and not a property of Finnish.} Every other text here is running prose and the \_Kalevala\_ is verse in strict trochaic tetrameter, compiled from oral songs. Its lines are more regular than its words and the detector finds the line. Replacing it with Finnish prose settles it:

\tiny
\setlength{\tabcolsep}{2pt}
\begin{longtable}{@{}>{\setlength{\rightskip}{0pt}\setlength{\parfillskip}{0pt}\arraybackslash}p{\dimexpr\textwidth/6-2\tabcolsep\relax}>{\setlength{\rightskip}{0pt}\setlength{\parfillskip}{0pt}\arraybackslash}p{\dimexpr\textwidth/6-2\tabcolsep\relax}>{\setlength{\rightskip}{0pt}\setlength{\parfillskip}{0pt}\arraybackslash}p{\dimexpr\textwidth/6-2\tabcolsep\relax}>{\setlength{\rightskip}{0pt}\setlength{\parfillskip}{0pt}\arraybackslash}p{\dimexpr\textwidth/6-2\tabcolsep\relax}>{\setlength{\rightskip}{0pt}\setlength{\parfillskip}{0pt}\arraybackslash}p{\dimexpr\textwidth/6-2\tabcolsep\relax}>{\setlength{\rightskip}{0pt}\setlength{\parfillskip}{0pt}\arraybackslash}p{\dimexpr\textwidth/6-2\tabcolsep\relax}@{}}
\toprule
text & boundary & gap & Zipf & brevity & distinct per token \\
\midrule
Kivi, \_Seitsemän veljestä\_, prose, 1870 & \texttt{0x20} & 6.82 & -0.828 & -0.298 & 0.37 \\
\_Kalevala\_, verse, 1849 & \texttt{0x2C} & 43.7 & -0.563 & -0.153 & 0.91 \\
Austen, prose, 1813 & \texttt{0x20} & 4.94 & -0.875 & -0.301 & 0.11 \\
\bottomrule
\end{longtable}
\normalsize

Finnish prose finds the space and returns a brevity of -0.298 against English prose at -0.301. The meter produced the failure, not the morphology, and an earlier draft of this section said otherwise.

\textbf{Agglutination is visible, in two places that are not the failure.} The mean word is 6.82 bytes against 4.94, and the distinct forms per token are 0.37 against 0.11, because a case heavy language spells one root many ways. Neither impairs the detector.

\textbf{What this does and does not establish.} Six prose texts, four languages, three families including one that is not Indo-European, from 1605 to 1870, every one in an alphabetic script. It is evidence that the regularities survive changes of language, family, morphology and century within that setting. It is not evidence about logographic scripts or spoken language, and the one verse text in the set behaves differently enough that genre is a variable this has not controlled for beyond replacing it.

\paragraph{4.13.0 Format Artifacts}

Four times in this section the tightest spacing in a file belonged to the file and not to what was written in it.

\tiny
\setlength{\tabcolsep}{2pt}
\begin{longtable}{@{}>{\setlength{\rightskip}{0pt}\setlength{\parfillskip}{0pt}\arraybackslash}p{\dimexpr\textwidth/2-2\tabcolsep\relax}>{\setlength{\rightskip}{0pt}\setlength{\parfillskip}{0pt}\arraybackslash}p{\dimexpr\textwidth/2-2\tabcolsep\relax}@{}}
\toprule
what was found & where \\
\midrule
\texttt{0x0D}, the line ending of a text wrapped at a fixed width & 5 of the 6 texts, before line endings were folded \\
\texttt{Heading}, Project Gutenberg's own markup & Austen's narrowest company in Section 4.14 \\
\texttt{0x2A} at a gap of 2.02 and a dispersion of 0.005, a ruled separator & a Greek text \\
\texttt{0x30} at a gap of 816, verse numbering & the same Greek text \\
\bottomrule
\end{longtable}
\normalsize

The most regular thing in a text file is almost never the language, because a language is produced by somebody composing and a file format is produced by a rule. Every result in Sections 4.12 to 4.14 required a layer of the format to be removed first, and the ones above were caught by reading the output and not by any check inside the method.

Two conditions now sit on a candidate boundary and both came out of these failures. A boundary has to have a dispersion above 0.05, because Section 4.12.1 shows that a perfectly regular boundary fixes the phase of its own period and carries nothing else. Decoration and padding do that. And it has to occur at least once in 64 bytes. A symbol can be regular and still be a page number, which marks the page and not the unit. Both were added while trying to read a text that was failing for a third reason covered in Section 4.13.05, and both leave every corpus that already worked unchanged.

\paragraph{4.13.05 Unicode Symbolization}

Every measurement above counted bytes, and no part of this document justified that. Section 4.10 already establishes the slice as a choice the method has to make, and the byte was chosen here by habit.

A Greek text exposes it. The detector returned \texttt{0x68} at a mean gap of 21.24, which is not the word boundary, and the universals collapsed to a Zipf slope of \texttt{-\allowbreak{}0.\allowbreak{}049}. The cause is visible in the file without knowing what language it holds. UTF-8 announces its own framing, a lead byte from \texttt{0x\allowbreak{}C2} to \texttt{0x\allowbreak{}DF} opening a two byte sequence and every following byte lying between \texttt{0x80} and \texttt{0x\allowbreak{}BF}, and reading that framing out of the Greek text gives 71.3\% two byte sequences against 99.4\% one byte sequences for the English one (\texttt{data/\allowbreak{}normalize\_\allowbreak{}symbols.\allowbreak{}py:29-\allowbreak{}58}). A byte level detector on the Greek text was measuring code units, and half a letter is not a symbol.

The alphabet is finite, so this is recoverable. Greek has been written with a bounded set of letters in every period it was written in, and the text carries 144 distinct symbols, so each one fits in a byte with room left. Re-seating them in order of first appearance, which is arbitrary with respect to every statistic measured, and folding line endings as the byte path already does, the detector returns \texttt{U+0020} as the boundary at a mean gap of 4.42 and the text reads:

\tiny
\setlength{\tabcolsep}{2pt}
\begin{longtable}{@{}>{\setlength{\rightskip}{0pt}\setlength{\parfillskip}{0pt}\arraybackslash}p{\dimexpr\textwidth/5-2\tabcolsep\relax}>{\setlength{\rightskip}{0pt}\setlength{\parfillskip}{0pt}\arraybackslash}p{\dimexpr\textwidth/5-2\tabcolsep\relax}>{\setlength{\rightskip}{0pt}\setlength{\parfillskip}{0pt}\arraybackslash}p{\dimexpr\textwidth/5-2\tabcolsep\relax}>{\setlength{\rightskip}{0pt}\setlength{\parfillskip}{0pt}\arraybackslash}p{\dimexpr\textwidth/5-2\tabcolsep\relax}>{\setlength{\rightskip}{0pt}\setlength{\parfillskip}{0pt}\arraybackslash}p{\dimexpr\textwidth/5-2\tabcolsep\relax}@{}}
\toprule
corpus & zipf slope & brevity & types at 25k & bits per word \\
\midrule
\texttt{greek\_\allowbreak{}iliad} & -0.878 & -0.352 & 7001 & 10.318 \\
\texttt{english\_\allowbreak{}1813\_\allowbreak{}austen} & -0.875 & -0.285 & 5690 & 9.792 \\
\texttt{french\_\allowbreak{}1862\_\allowbreak{}hugo} & -0.950 & -0.311 & 7432 & 10.324 \\
\texttt{german\_\allowbreak{}1808\_\allowbreak{}goethe} & -0.723 & -0.357 & 8064 & 10.891 \\
\texttt{finnish\_\allowbreak{}1870\_\allowbreak{}kivi\_\allowbreak{}prose} & -0.828 & -0.313 & 11841 & 11.900 \\
\bottomrule
\end{longtable}
\normalsize

Greek sits inside the group on every column, carrying the second strongest brevity correlation of any natural text measured here. Nothing in the detector knows Greek, and the only thing it needed was the width.

Two checks say the re-slice is not doing the work itself. Applied to the English text, where the byte was already the right width, it moves the Zipf slope from -0.875 to -0.875 and the type count from 5691 to 5690. And the boundary it finds is not planted: seats are assigned by first appearance, the detector ranks candidates on gap regularity, and it recovers \texttt{U+0020} in all four texts re-sliced this way.

\paragraph{4.13.06 Translation and Composition}

The 1611 King James and the 1623 Shakespeare are twelve years apart in one language, one translated from Hebrew and Aramaic and one composed. Capped at the same 1048576 symbols, Shakespeare uses fewer tokens, 188474 against 199878, and 92\% more distinct types, 25031 against 13058. So the King James carries the smallest vocabulary of any prose measured here, and the century is not what put it there.

The brevity correlation separates the same pair without using vocabulary at all. It runs -0.247 for the King James against -0.315 for Shakespeare, and every text composed in its own language sits between -0.285 and -0.357. Brevity is a property of production, since the speaker shortening a frequent word is the one paying for the length. A translator is not that speaker and chooses for sense, which decouples the frequency from the length.

Two measures with nothing in common therefore separate a translation from a composition of the same decade, and a source language whose words carry several senses at once is the mechanism both of them would report. Neither measurement identifies that mechanism, and neither reads a word.

\paragraph{4.13.07 Genre Controls}

Epic was tested for a distinct register and twice returned nothing, both times for a reason in the instrument.

Formula reuse was the first attempt, since epic is composed from repeated phrases. The Greek Iliad returned 108 formulas per 25000 tokens against 142 for Shakespeare and 128 for Austen, which is no signature.

Inflection was the first explanation offered for that, on the reasoning that a Homeric epithet declines with its noun so no exact match survives, and it is wrong. English barely inflects. The test was repeated on two English epics that are each a single work. Milton returns 120 and Pope's Iliad returns 98, both under Austen's 128 on the same prose novel the Greek text was already under. Pope carries the same story as the Greek text in a language where the stated mechanism cannot operate, and the count goes down instead of up.

So the surface form argument does not survive, and neither does any claim that epic carries elevated formula reuse. Either the register has no such signature or this counter does not measure what an oral formula is.

Burstiness was the second, since a large named cast should cluster. The Iliad returned 4.03, second lowest of every corpus, while Shakespeare returned 40.06 and the King James 20.77. Those two are anthologies, 37 plays and 66 books, and a name confined to one member of a collection bursts for that reason. The measure is reporting how many separate works were concatenated.

So this document makes no claim about epic register, and it withdraws the explanation it first gave for not finding one. Three measurements over five texts in three languages, two of them chosen to remove the confound blamed the first time, return nothing that separates epic from prose.

The burst figures cannot be read back on a re-sliced corpus at all, because the reporting path prints the matched unit as bytes and those bytes are seat indices after Section 4.13.05 re-seats them.

There is a limit behind all three failures that no better instrument reaches. Every measure here reads how a text was produced, and a story someone invented is produced by the same species using the same machinery as a report of something witnessed. The two come out alike because they are alike in the only respect being measured. This is Proposition 2 arriving from the direction of the corpus: the signature is a necessary condition of human production and never a sufficient one for what happened. A text can carry every regularity in Section 4.13 and describe nothing that occurred.

\paragraph{4.13.08 Diachronic Drift}

Drift is a distance between two texts, which needs a comparison the per corpus measures do not make. Two frequency bands are compared as a Jaccard overlap (\texttt{examples/\allowbreak{}language/\allowbreak{}4\_\allowbreak{}measure/\allowbreak{}measure\_\allowbreak{}drift.\allowbreak{}py:38-\allowbreak{}52}). The top 100 words of an English text are almost all function words and move slowly. Ranks 500 to 2000 are mostly content and follow the subject.

\tiny
\setlength{\tabcolsep}{2pt}
\begin{longtable}{@{}>{\setlength{\rightskip}{0pt}\setlength{\parfillskip}{0pt}\arraybackslash}p{\dimexpr\textwidth/5-2\tabcolsep\relax}>{\setlength{\rightskip}{0pt}\setlength{\parfillskip}{0pt}\arraybackslash}p{\dimexpr\textwidth/5-2\tabcolsep\relax}>{\setlength{\rightskip}{0pt}\setlength{\parfillskip}{0pt}\arraybackslash}p{\dimexpr\textwidth/5-2\tabcolsep\relax}>{\setlength{\rightskip}{0pt}\setlength{\parfillskip}{0pt}\arraybackslash}p{\dimexpr\textwidth/5-2\tabcolsep\relax}>{\setlength{\rightskip}{0pt}\setlength{\parfillskip}{0pt}\arraybackslash}p{\dimexpr\textwidth/5-2\tabcolsep\relax}@{}}
\toprule
pair & years apart & form & top 100 & ranks 500 to 2000 \\
\midrule
Milton 1667, Pope 1720 & 53 & both epic verse & 0.5625 & 0.2310 \\
King James 1611, Shakespeare 1623 & 12 & scripture, plays & 0.5038 & 0.1596 \\
Pope 1720, Austen 1813 & 93 & epic verse, novel & 0.3889 & 0.1207 \\
King James 1611, Austen 1813 & 202 & scripture, novel & 0.3699 & 0.1054 \\
Kalevala 1849, Kivi 1870 & 21 & Finnish verse, Finnish novel & 0.2121 & 0.0684 \\
Austen 1813, Hugo 1862 & control & different languages & 0.0204 & 0.0132 \\
\bottomrule
\end{longtable}
\normalsize

The last row shows the measure behaves, since two languages share almost nothing and the figure goes to the floor. The first two rows show it cannot answer the question. Fifty three years with the form held constant preserves more of the common vocabulary than twelve years across a change of form. The genre term is larger than the time term and one genre matched pair cannot separate them.

A second objection is stronger than the statistical one. A drift floor is a claim about a population living out normal lifespans in one place with nothing interrupting the chain of transmission, and the only genre matched pair available spans the English Civil War, the regicide, the plague of 1665 and the fire of 1666. That interval is the opposite of the condition, so 0.5625 is not a floor. Read as an upper bound on how fast the common vocabulary moves, it says roughly 72 of the 100 most frequent words survived that half century unchanged.

Measuring the floor needs several pairs matched for form at different separations, drawn from a period with no disaster and no displacement. This corpus contains no such series.

Held against one 1611 reference the English texts give 0.5038 at 1623, 0.4286 at 1667, 0.4184 at 1720 and 0.3699 at 1813 in the top 100 band, which decreases at every step. The content band reverses once, rising from 0.1596 at 1623 to 0.1632 at 1667, and Milton is retelling the subject of the book being compared against. The topic produced that reversal, not the century.

Converting the top 100 band to a loss per year gives 0.00171 across 1623 to 1667, 0.00019 across 1667 to 1720, and 0.00052 across 1720 to 1813. The first interval covers the civil war, the regicide, the plague and the fire, and it moves 8.9 times faster than the second, which covers the settled period after the Restoration. That ordering is what a drift rate responding to disruption would produce.

That ordering is not evidence of one. The 1667 to 1720 interval is the only pair in the series matched for form. The slowest interval is also the one interval where the largest confound was removed. The disruption term and the genre term take the same value on the same rows, and nothing measured here separates them. Whether a perturbation produces an oscillation is a further question this corpus cannot reach at all: five samples spread unevenly over 202 years are fewer than two per period for any period worth proposing, so no periodicity is detectable and none is excluded.

Holding the subject fixed removes the genre term by construction. Three translations of one book were measured against each other. Douay Rheims 1609 and the King James 1611 are two years apart and were made by rival translators from different sources, which prices translator choice with the era held still. The World English Bible is 389 years after the King James.

\tiny
\setlength{\tabcolsep}{2pt}
\begin{longtable}{@{}>{\setlength{\rightskip}{0pt}\setlength{\parfillskip}{0pt}\arraybackslash}p{\dimexpr\textwidth/4-2\tabcolsep\relax}>{\setlength{\rightskip}{0pt}\setlength{\parfillskip}{0pt}\arraybackslash}p{\dimexpr\textwidth/4-2\tabcolsep\relax}>{\setlength{\rightskip}{0pt}\setlength{\parfillskip}{0pt}\arraybackslash}p{\dimexpr\textwidth/4-2\tabcolsep\relax}>{\setlength{\rightskip}{0pt}\setlength{\parfillskip}{0pt}\arraybackslash}p{\dimexpr\textwidth/4-2\tabcolsep\relax}@{}}
\toprule
pair & years apart & top 100 & ranks 500 to 2000 \\
\midrule
Douay 1609, King James 1611 & 2 & 0.8519 & 0.5440 \\
King James 1611, World English 2000 & 389 & 0.7857 & 0.6243 \\
Douay 1609, World English 2000 & 391 & 0.7699 & 0.4684 \\
\bottomrule
\end{longtable}
\normalsize

Every figure here is far above the genre varied pairs above, which confirms the subject was carrying most of that difference. Within these three the era is the smaller term. Two years across a change of translator costs 0.1481 of the top 100 band and 389 years costs 0.0662. The language moving for four centuries does less to the common vocabulary than the choice of who is translating.

The content band settles it. The King James and the World English agree better across 389 years, 0.6243, than the two translations two years apart do, 0.5440. Douay Rheims came through the Latin Vulgate and the World English descends from the American Standard in the King James line. That column is reporting which textual tradition a translation stands in.

Three designs have now been tried and each one lost to a different confound: pairs matched for form lost to genre, a fixed reference series lost to genre, and a fixed subject lost to translator lineage. Every series reachable here varies something that moves vocabulary harder than time does. No drift rate is reported and the bunching of changes against periods of disruption is not tested. What the last table does support is the stability itself. A top 100 overlap of 0.7857 across 389 years is a common vocabulary that barely moves, and any bunching would have to be visible against it.

There is a further limit in the measure that no additional corpus removes. Across 389 years of one language it reads 0.7857, and across two languages it reads 0.0204, and nothing measured here falls between those. The instrument has a slow regime and a total regime with no resolved middle.

That gap is where a civilization ending sits. A port silting up until the city is abandoned does not speed up the drift of the vocabulary spoken there, it removes the speakers, and the ground then carries a different language. Ephesus runs through Anatolian, Ionic Greek, Koine, Byzantine Greek and Turkish by that route. A change at the edge of a civilization substitutes one language for another instead of altering a rate, and substitution is discrete, which is a mechanism for changes arriving in groups that does not require the drift rate to vary at all. Measuring it needs a replacement event with text on both sides of it. The drift rate is the only quantity these measurements compute.

For the oldest replacement events there is no text on either side and there will not be one. Göbekli Tepe was backfilled by the people who built it around 8000 BC and carries sliced pillars, repeated animal figures and an evident organizing intent, with no recoverable language attached to any of it. The community ending events in the archaeological record are the same shape, legible in the deposit and silent in every respect this document measures. So the limit there is not the size of the corpus available. There is no reading that recovers what those marks meant.

Writing is also one channel among several, and the others carry information this document cannot read at all. Grave assemblages running to tens of thousands of worked beads, in materials sourced across a continent, are evidence of an organized system carrying status and obligation, and the labor in one of them is measured in years. Bodies assembled from several individuals, kept above ground for generations and then buried under an occupied floor, are a deliberate link to the past maintained in bone. Both are societies holding continuity hard, and neither leaves a term that can be read. A gap in the written record across an event is not evidence that continuity failed there, and this document measures the one channel that happens to survive as text.

There is a selection effect in this that no corpus work removes. Britain and Ireland carry no indigenous writing until the Roman period. The community ending deposits found there predate any local written record by two thousand years at the recent end and by far more at the older end. The events most likely to show changes arriving in groups are therefore the events that systematically leave nothing to read, and the intervals that can be measured are the ones mild enough that a writing tradition survived them. Those are the intervals where drift is what happened, and the measurements above report that. A corpus of written language is a sample of the surviving cases, and it is biased against the phenomenon being asked about.

That is the strongest form of the asymmetry this document is built on. A pattern can be certainly present and its meaning permanently unrecoverable, which is Proposition 1 holding while Proposition 2 refuses, stated by the record instead of by an argument.

None of this makes the regularities fragile, and the reason is the one result in this section that came from a controlled destruction instead of a comparison. Section 4.13.1 permutes the Morse code table, the part a person designed and the part actually transmitted, and the brevity correlation returns at -0.550 against -0.652 for the intact table while the Zipf slope does not move, -1.183 against -1.185. The assignment of codes to letters was destroyed and the measurements reassembled out of the wreckage, because they were never held in that assignment. They belong to the distribution of the words being encoded, and they reappear through any re-slice of the symbols.

So these regularities are not transmitted and do not need a surviving channel between two instances of them. They are properties of whatever produces a message, and they recur wherever that kind of producer recurs. Three language families, two scripts, 265 years and one percussive medium give the same measures for that reason and not because anything passed between them.

The mechanism is that a shorter code for a more frequent symbol is what minimizes an expected transmission cost, which is a result in coding theory and not a fact about people. Any system paying per unit sent, over enough repetitions, arrives at it. So does a telegraph operator assigning codes by hand. Section 4.13.1 finds most of Morse's brevity surviving the destruction of the table he designed. A language is solved against one set of constraints, a vocal tract of fixed bandwidth, a working memory holding a few items, a listener decoding in real time and a finite childhood to learn in, and the optimum belongs to that constraint set instead of to any lineage. Nothing inherits the answer. Each language derives it again, because the problem is still there. Losing the content therefore does not remove the regularity: the regularity was never stored, and the conditions that regenerate it outlast every artifact.

One caution belongs with this and it is now measured. Distributions of this shape also arise from processes doing no optimizing at all, and Section 7.4 runs that counterexample against these same measures. Independent character draws carrying a delimiter return a Zipf slope of -0.988, inside the range every natural corpus here occupies. The Zipf slope in this section carries no evidence about production and is withdrawn as such. The brevity correlation survives weakened, at -0.222 for the control against -0.247 to -0.364 for natural text. The boundary regularity against a permutation null survives intact, at 1.00 for the control against 1.79 and above for natural text. That counterexample does not weaken the method in this document, and strengthens it. A structure produced by nearly any generator is a better thing to key on than one belonging to a special class of them, and the requirement here was never that the source be a language. What the method needs is that something regular produced the bytes, so it needs no alphabet, no dimension and no interpretation.

The measures in Section 4.13 are also close to a question that has been asked before, since deciding whether an undeciphered corpus is a language at all is what a Zipf slope and a brevity correlation get used for. That literature exists and it is contested. It is not cited here because it has not been read in full, and Section 7 states what was read and what it cost.

\paragraph{4.13.09 Frequency-Stratified Structure}

Section 7.4.1 finds the two halves of the frequency distribution behaving differently, the frequent one carrying regular boundaries and the rare one carrying clustering. Neither measurement reads a word. Printing the twelve most frequent words of each corpus (\texttt{examples/\allowbreak{}language/\allowbreak{}4\_\allowbreak{}measure/\allowbreak{}top\_\allowbreak{}words.\allowbreak{}py}) shows what fills the half that all of the statistics describe without naming:

\tiny
\setlength{\tabcolsep}{2pt}
\begin{longtable}{@{}>{\setlength{\rightskip}{0pt}\setlength{\parfillskip}{0pt}\arraybackslash}p{\dimexpr\textwidth/2-2\tabcolsep\relax}>{\setlength{\rightskip}{0pt}\setlength{\parfillskip}{0pt}\arraybackslash}p{\dimexpr\textwidth/2-2\tabcolsep\relax}@{}}
\toprule
corpus & twelve most frequent words \\
\midrule
\texttt{english\_\allowbreak{}1813\_\allowbreak{}austen} & the to of and her i a in was she that it \\
\texttt{english\_\allowbreak{}1611\_\allowbreak{}kjv} & the and of to that in he shall unto for i his \\
\texttt{french\_\allowbreak{}1862\_\allowbreak{}hugo} & de il la et le l à un les que une d \\
\texttt{german\_\allowbreak{}1808\_\allowbreak{}goethe} & und ich die der nicht das ein ist zu du in sie \\
\texttt{spanish\_\allowbreak{}1605\_\allowbreak{}cervantes} & que de y la a en el no los se con por \\
\texttt{finnish\_\allowbreak{}1870\_\allowbreak{}kivi\_\allowbreak{}prose} & ja mutta hän juhani on niin kuin nyt oli ei he hänen \\
\texttt{greek\_\allowbreak{}iliad} & και κι του ο το να τον τα τους με τ η \\
\bottomrule
\end{longtable}
\normalsize

Every head carries reference to persons. English has \texttt{her}, \texttt{i}, \texttt{she} and \texttt{it}, German has \texttt{ich}, \texttt{du} and \texttt{sie}, Finnish has \texttt{hän}, \texttt{he} and \texttt{hänen}, and Greek carries \texttt{του}, \texttt{τον} and \texttt{τους}. Each also carries a conjunction, and each carries either a copula or a negation or both. Finnish is the sharpest case: it has no articles to occupy those positions, and a proper name, \texttt{juhani}, takes one of them.

The corresponding claim about the rare half is that it carries the actions, which needs counting. No part of speech tagger is available here. English verb inflection stands in for one, and the proxy is weak in one direction: a gerund used as a noun and an adjective built from a participle both carry these endings without being verbs, so every figure is an upper bound (\texttt{examples/\allowbreak{}language/\allowbreak{}4\_\allowbreak{}measure/\allowbreak{}head\_\allowbreak{}tail\_\allowbreak{}parts.\allowbreak{}py:31-\allowbreak{}36}).

\tiny
\setlength{\tabcolsep}{2pt}
\begin{longtable}{@{}>{\setlength{\rightskip}{0pt}\setlength{\parfillskip}{0pt}\arraybackslash}p{\dimexpr\textwidth/3-2\tabcolsep\relax}>{\setlength{\rightskip}{0pt}\setlength{\parfillskip}{0pt}\arraybackslash}p{\dimexpr\textwidth/3-2\tabcolsep\relax}>{\setlength{\rightskip}{0pt}\setlength{\parfillskip}{0pt}\arraybackslash}p{\dimexpr\textwidth/3-2\tabcolsep\relax}@{}}
\toprule
corpus & head, percent ending in \texttt{-\allowbreak{}ed} or \texttt{-\allowbreak{}ing} & rare half \\
\midrule
\texttt{english\_\allowbreak{}1667\_\allowbreak{}milton\_\allowbreak{}epic} & 10.1 & 22.3 \\
\texttt{english\_\allowbreak{}1623\_\allowbreak{}shakespeare} & 9.7 & 18.5 \\
\texttt{english\_\allowbreak{}1813\_\allowbreak{}austen} & 16.9 & 23.7 \\
\texttt{english\_\allowbreak{}1611\_\allowbreak{}kjv} & 12.9 & 15.4 \\
\texttt{monkey\_\allowbreak{}a26\_\allowbreak{}d18\_\allowbreak{}english} & 0.0 & 0.0 \\
\bottomrule
\end{longtable}
\normalsize

The rare half carries more of it in all four, by 1.19 to 2.21 times, and the memoryless control produces none, so the enrichment is not an artifact of counting endings. The King James is the weakest row and its own morphology explains that: it conjugates as \texttt{saith}, \texttt{cometh} and \texttt{spake}, which this proxy does not count at all. Its verbs are undercounted by an amount not estimated here.

So the frequent half holds the machinery for referring to people and predicating anything of them, and the rare half holds more of what is being done. That describes the two halves. The description reaches nothing about what any text means, and the statistics in Section 7.4 that separate a corpus from a shuffle of itself do not know that either half exists.

Scoring individual words for clustering against the same permutation null (\texttt{examples/\allowbreak{}language/\allowbreak{}4\_\allowbreak{}measure/\allowbreak{}word\_\allowbreak{}burstiness.\allowbreak{}py:44-\allowbreak{}72}) answers two questions, and only the second of them comes out.

The first was whether words for things people do constantly, instead of things a passage is about, spread more evenly than the rest. Section 7.4.4 predicts they should. A supplied English set of them gives 0.843 against a corpus average of 0.828 for Austen, 0.521 against 0.509 for the King James, and 0.809 against 0.796 for Milton. The direction is the predicted one every time and the size is not worth reporting, between 7 and 20 words cleared the occurrence floor, and the prediction is not established.

The corpus averages themselves separate, which the first question was not asking about.

\tiny
\setlength{\tabcolsep}{2pt}
\begin{longtable}{@{}>{\setlength{\rightskip}{0pt}\setlength{\parfillskip}{0pt}\arraybackslash}p{\dimexpr\textwidth/3-2\tabcolsep\relax}>{\setlength{\rightskip}{0pt}\setlength{\parfillskip}{0pt}\arraybackslash}p{\dimexpr\textwidth/3-2\tabcolsep\relax}>{\setlength{\rightskip}{0pt}\setlength{\parfillskip}{0pt}\arraybackslash}p{\dimexpr\textwidth/3-2\tabcolsep\relax}@{}}
\toprule
corpus & word burstiness & what the corpus is \\
\midrule
\texttt{monkey\_\allowbreak{}a26\_\allowbreak{}d18\_\allowbreak{}geom90} & 1.010 & a memoryless process \\
\texttt{greek\_\allowbreak{}iliad} & 0.833 & one poem \\
\texttt{english\_\allowbreak{}1813\_\allowbreak{}austen} & 0.828 & one novel \\
\texttt{english\_\allowbreak{}1667\_\allowbreak{}milton\_\allowbreak{}epic} & 0.796 & one poem \\
\texttt{spanish\_\allowbreak{}1605\_\allowbreak{}cervantes} & 0.780 & one novel \\
\texttt{english\_\allowbreak{}1720\_\allowbreak{}pope\_\allowbreak{}iliad\_\allowbreak{}epic} & 0.780 & one poem \\
\texttt{finnish\_\allowbreak{}1870\_\allowbreak{}kivi\_\allowbreak{}prose} & 0.770 & one novel \\
\texttt{french\_\allowbreak{}1862\_\allowbreak{}hugo} & 0.758 & one novel \\
\texttt{english\_\allowbreak{}1623\_\allowbreak{}shakespeare} & 0.667 & 37 plays \\
C source & 0.533 & 40 files \\
\texttt{english\_\allowbreak{}1611\_\allowbreak{}kjv} & 0.509 & 66 books \\
\bottomrule
\end{longtable}
\normalsize

Single works occupy 0.758 to 0.833 over four languages and three centuries, collections occupy 0.509 to 0.667, and the memoryless control sits at 1.010. The bands do not overlap. The C sources fall with the collections, and forty files each concerning a different module should give that.

Hugo is the lowest of the single works and that supports the reading instead of straining it. \_Les Misérables\_ runs to five volumes and turns aside into Waterloo, the sewers of Paris and the history of a convent. It is the most topically varied single work measured here and it sits nearest the collection band.

Every corpus in that table arrived already being one thing or the other. The reading rests on labels applied by hand and could be tracking whatever those labels correlate with. Two manipulations settle it (\texttt{examples/\allowbreak{}any\_\allowbreak{}corpus/\allowbreak{}3\_\allowbreak{}reference/\allowbreak{}heterogeneity\_\allowbreak{}control.\allowbreak{}py}). Joining English single works into one corpus, with the total length held near 900000 characters so the walk cannot be a length effect, gives 0.828 for one work, 0.685 for two and 0.626 for three. Cutting the 66 book collection into equal pieces walks the other way:

\tiny
\setlength{\tabcolsep}{2pt}
\begin{longtable}{@{}>{\setlength{\rightskip}{0pt}\setlength{\parfillskip}{0pt}\arraybackslash}p{\dimexpr\textwidth/4-2\tabcolsep\relax}>{\setlength{\rightskip}{0pt}\setlength{\parfillskip}{0pt}\arraybackslash}p{\dimexpr\textwidth/4-2\tabcolsep\relax}>{\setlength{\rightskip}{0pt}\setlength{\parfillskip}{0pt}\arraybackslash}p{\dimexpr\textwidth/4-2\tabcolsep\relax}>{\setlength{\rightskip}{0pt}\setlength{\parfillskip}{0pt}\arraybackslash}p{\dimexpr\textwidth/4-2\tabcolsep\relax}@{}}
\toprule
pieces & books per piece & mean & highest piece \\
\midrule
1 & 66 & 0.509 &  \\
8 & 8.2 & 0.612 & 0.700 \\
16 & 4.1 & 0.648 & 0.740 \\
32 & 2.1 & 0.681 & 0.835 \\
64 & 1.0 & 0.708 & 0.847 \\
\bottomrule
\end{longtable}
\normalsize

The climb continues at every cut, and at one book per piece the highest pieces reach 0.847, which is inside the band the single works occupy. The mean stays at 0.708 for two reasons that are properties of the cutting: the pieces are equal in length so they do not align with book boundaries and most of them straddle two books, and a short piece has fewer words clearing the occurrence floor.

So the quantity responds to subjects being added and removed while length is held fixed, which is a stronger statement than the table of labeled corpora supports on its own.

The measure is not a count on its own. Two English works joined give 0.685, close to the 0.667 of a corpus of 37 plays, because a novel and an epic poem stand further apart than two plays of one period do. What the measure reads is how far the vocabulary spreads across subjects, with the number of them and the distance between them both contributing and neither separated here.

That is a property of what a text is about, recovered without reading any of it, and it is the only measurement in this document that orders corpora on a scale with resolved bands. It still names no subject and identifies no meaning. It counts how far the vocabulary gathers.

\paragraph{4.13.1 Re-encoding Invariance}

Every row above is alphabetic, and a regularity common to all of them could belong to that way of writing instead of to language. Re-encoding one of them settles it, because a re-encoding holds the meaning fixed and changes nothing else. Morse carries the same words over two marks and two silences, and \texttt{examples/\allowbreak{}language/\allowbreak{}1\_\allowbreak{}represent/\allowbreak{}encode\_\allowbreak{}percussive.\allowbreak{}py} produces it.

\tiny
\setlength{\tabcolsep}{2pt}
\begin{longtable}{@{}>{\setlength{\rightskip}{0pt}\setlength{\parfillskip}{0pt}\arraybackslash}p{\dimexpr\textwidth/6-2\tabcolsep\relax}>{\setlength{\rightskip}{0pt}\setlength{\parfillskip}{0pt}\arraybackslash}p{\dimexpr\textwidth/6-2\tabcolsep\relax}>{\setlength{\rightskip}{0pt}\setlength{\parfillskip}{0pt}\arraybackslash}p{\dimexpr\textwidth/6-2\tabcolsep\relax}>{\setlength{\rightskip}{0pt}\setlength{\parfillskip}{0pt}\arraybackslash}p{\dimexpr\textwidth/6-2\tabcolsep\relax}>{\setlength{\rightskip}{0pt}\setlength{\parfillskip}{0pt}\arraybackslash}p{\dimexpr\textwidth/6-2\tabcolsep\relax}>{\setlength{\rightskip}{0pt}\setlength{\parfillskip}{0pt}\arraybackslash}p{\dimexpr\textwidth/6-2\tabcolsep\relax}@{}}
\toprule
 & boundary & gap & dispersion & Zipf & brevity \\
\midrule
Austen 1813, as written & \texttt{0x20} & 4.94 & 0.374 & -0.875 & -0.301 \\
Austen 1813, in Morse & \texttt{0x20} & 4.59 & 0.217 & -1.185 & -0.652 \\
\bottomrule
\end{longtable}
\normalsize

Given four distinct marks and no letters, the detector finds a boundary and both regularities are present. Brevity is more than twice as strong, which looked at first like the one result here that was designed and not discovered, since Morse assigns its shortest codes to the most frequent letters. Section 4.13.2 tests that and it is mostly not so.

\textbf{The two rows are not at the same level and should not be read as a pair.} English splits at words. Morse splits at the letter gap, because within a word that gap is more regular than the word gap is. The unit counts say so: 1187 distinct units is far more than the 36 codes Morse has, and the encoder writes the word slash with no space beside it. The last code of one word, the slash, and the first code of the next form one unit. Thirty six codes plus up to 1296 such pairs brackets the 1187 seen. The transcription invented a vocabulary of boundary spanning pairs that is in neither language.

What the row does establish is narrow and is the part worth having. Neither regularity is a property of alphabetic writing. Both survive into a four symbol percussive representation of the same text.

\paragraph{4.13.2 Code-Table Permutation}

Morse had regional variants that gave different codes to the same letters, and a code table is a choice and not part of what is being said. Permuting it holds the message, the language and the set of code lengths all fixed and changes only which letter received which code. That separates a property of the message from a property of the encoding, and the two regularities should come apart under it: how often a word recurs cannot depend on how its letters are spelled, and how long a code is was somebody's decision.

\tiny
\setlength{\tabcolsep}{2pt}
\begin{longtable}{@{}>{\setlength{\rightskip}{0pt}\setlength{\parfillskip}{0pt}\arraybackslash}p{\dimexpr\textwidth/4-2\tabcolsep\relax}>{\setlength{\rightskip}{0pt}\setlength{\parfillskip}{0pt}\arraybackslash}p{\dimexpr\textwidth/4-2\tabcolsep\relax}>{\setlength{\rightskip}{0pt}\setlength{\parfillskip}{0pt}\arraybackslash}p{\dimexpr\textwidth/4-2\tabcolsep\relax}>{\setlength{\rightskip}{0pt}\setlength{\parfillskip}{0pt}\arraybackslash}p{\dimexpr\textwidth/4-2\tabcolsep\relax}@{}}
\toprule
 & Zipf & brevity & encoded bytes \\
\midrule
Morse, the real table & -1.185 & -0.652 & 2,001,933 \\
Morse, the same codes permuted & -1.183 & -0.550 & 2,580,891 \\
\bottomrule
\end{longtable}
\normalsize

\textbf{Zipf is unchanged to two parts in a thousand}. A property of the message has to do that under a relabeling of its symbols.

\textbf{Brevity is nearly unchanged, and that refutes the reading offered in Section 4.13.1.} A permuted table keeps a correlation of -0.550. At most 0.10 of the -0.652 can be credited to Morse assigning short codes to frequent letters. The rest is structural: short strings recur more often because there are fewer short strings, and that holds for any table whatever. The stronger brevity in the Morse rows than in the English row is mostly an artifact of the unit length distribution and not evidence of design.

\textbf{What the design does buy is the length.} The real table encodes the same text in 22.4\% fewer bytes than a permutation of its own codes. That is Morse's engineering, and it appears in transmission cost and not in the correlation it was attributed to.

So the two are not the same kind of claim. Zipf is invariant under re-encoding and is about the message. The brevity correlation is also largely invariant, which makes it mostly a fact about how many short strings exist and only slightly a fact about anyone's choices.

\paragraph{4.13.3 Within-Language Variance}

Every result above compares one text in one language against another text in another, and a difference could belong to the language or to the text. Separating those needs replication inside a language, and the sections above have twelve English texts and one or two of everything else.

Four texts were taken in each of seven alphabetic languages from Project Gutenberg's own index for each, which is not a hand chosen selection, and four Chinese novels beside them (\texttt{maint/\allowbreak{}data/\allowbreak{}fetch/\allowbreak{}fetch\_\allowbreak{}by\_\allowbreak{}language.\allowbreak{}py}). Everything is measured at character width and not at byte width (\texttt{examples/\allowbreak{}language/\allowbreak{}4\_\allowbreak{}measure/\allowbreak{}language\_\allowbreak{}constant.\allowbreak{}py:44-\allowbreak{}70}). That is required for the Chinese to be comparable at all, since one of those novels carries 3164 distinct characters where no byte seating reaches, and it is also the width at which a Chinese symbol is a morpheme instead of a part of one.

The question is then a variance question. If a quantity belongs to a language, its spread within one language is small against its spread between languages, and an analysis of variance reports that (\texttt{src/\allowbreak{}engine/\allowbreak{}r/\allowbreak{}hypotheses/\allowbreak{}language\_\allowbreak{}variance.\allowbreak{}R}).

\tiny
\setlength{\tabcolsep}{2pt}
\begin{longtable}{@{}>{\setlength{\rightskip}{0pt}\setlength{\parfillskip}{0pt}\arraybackslash}p{\dimexpr\textwidth/5-2\tabcolsep\relax}>{\setlength{\rightskip}{0pt}\setlength{\parfillskip}{0pt}\arraybackslash}p{\dimexpr\textwidth/5-2\tabcolsep\relax}>{\setlength{\rightskip}{0pt}\setlength{\parfillskip}{0pt}\arraybackslash}p{\dimexpr\textwidth/5-2\tabcolsep\relax}>{\setlength{\rightskip}{0pt}\setlength{\parfillskip}{0pt}\arraybackslash}p{\dimexpr\textwidth/5-2\tabcolsep\relax}>{\setlength{\rightskip}{0pt}\setlength{\parfillskip}{0pt}\arraybackslash}p{\dimexpr\textwidth/5-2\tabcolsep\relax}@{}}
\toprule
quantity & spread within & spread between & \(F\) & \(p\) \\
\midrule
mean distance between boundaries & 0.3327 & 0.6980 & 13.21 & \(3.5 \times 10^{-6}\) \\
collision entropy & 0.0430 & 0.0746 & 9.02 & \(6.2 \times 10^{-5}\) \\
rare half against a permutation null & 0.0731 & 0.0342 & 0.66 & 0.68 \\
\bottomrule
\end{longtable}
\normalsize

The first two belong to a language. How long its words are and what its character inventory looks like separate the seven at \(p\) below \(10^{-4}\). Those are constants of a language and the earlier readings of them are readings of the language.

The third does not. Its spread between languages, 0.0342, is smaller than its spread within one, 0.0731. Two novels in one language differ on it by more than two languages do. That quantity carries no information about which language produced the text.

Chinese decides it. Measured in within language standard deviations from the mean of the seven, it stands 66.0 away on collision entropy and 22.1 away on the mean distance between boundaries, and 0.5 away on the rare half. A writing system that encodes morphemes instead of sounds, with thirty times the symbol inventory, moves every structural quantity by tens of deviations and leaves the arrangement measure where it was.

The results in Sections 4.13.05 and 4.13.1 hold over seven language families and two scripts for that reason. The measure was never reading which language it was given, and a change of language had nothing in it to disturb. It also bounds what those results can mean: a quantity that cannot tell Finnish from Chinese cannot be evidence about any particular language, and Section 7.4 records separately that it is not evidence about language at all.

\textbf{This conclusion is not new and the credit is not this document's.} Montemurro and Zanette, \_Universal Entropy of Word Ordering Across Linguistic Families\_, PLoS ONE 6(5):e19875 (2011), measure the entropy a text loses when its words are shuffled, the same null used here, over 7077 texts in eight corpora spanning Indo-European, Finno-Ugric, Austronesian, Afroasiatic and Sino-Tibetan families and the Sumerian isolate. They report the quantity bounded near 3.3 bits per word across all of them, with a relative variability of 0.07 against 0.23 for the entropy itself, and they name it a statistical linguistic universal. Their sample is two hundred times the size of this one and reaches Old Egyptian and Sumerian, which this one does not.

The statistic here is not theirs, since they measure a Shannon entropy reduction in bits per word and this measures a dispersion of gaps over the rare half of the symbols, and the agreement of two unlike statistics is worth something. The finding is theirs.

What this document adds has to be stated carefully, because two different instruments appear in it and they have been run together in earlier drafts. The dispersion measure used here has been applied to text, to source code and to recorded sound, which are unlike media, and it separates an arranged domain from an unarranged one in each. The dispersion measure has not been applied to the pictures in Section 4.11 or to the crystal lattices, which were measured by the shift detector instead, and that is a separate instrument with a separate null.

Even inside the media it does cover, the readings are not directly comparable. A waveform carries a smoothness that text does not, the symbol width differs between them, and Section 4.10 establishes that the width governs what is visible. So the supportable claim is that one measure runs on unlike media and separates in each of them, and not that it returns comparable values across them.

Four texts in a language is a small sample and all of them come from one publisher. The figures for any single language here carry the selection of that source with them. What the variance decomposition needs is replication, which it has, and not breadth.

\subsubsection{4.14 Contextual Diversity}

Frequency cannot separate two kinds of common unit. In an English bible \texttt{the} and \texttt{God} are both frequent, and one is frequent because it attaches to anything while the other is frequent because the book is about it. The company separates them: a unit doing grammatical work has nearly as many distinct followers as occurrences, and a unit belonging to a narrow subject has far fewer (\texttt{test/\allowbreak{}bench/\allowbreak{}bench\_\allowbreak{}ancorae\_\allowbreak{}ab.\allowbreak{}c:2007}).

The lowest variety among units seen at least 24 times, over the first 25000 units:

\tiny
\setlength{\tabcolsep}{2pt}
\begin{longtable}{@{}>{\setlength{\rightskip}{0pt}\setlength{\parfillskip}{0pt}\arraybackslash}p{\dimexpr\textwidth/8-2\tabcolsep\relax}>{\setlength{\rightskip}{0pt}\setlength{\parfillskip}{0pt}\arraybackslash}p{\dimexpr\textwidth/8-2\tabcolsep\relax}>{\setlength{\rightskip}{0pt}\setlength{\parfillskip}{0pt}\arraybackslash}p{\dimexpr\textwidth/8-2\tabcolsep\relax}>{\setlength{\rightskip}{0pt}\setlength{\parfillskip}{0pt}\arraybackslash}p{\dimexpr\textwidth/8-2\tabcolsep\relax}>{\setlength{\rightskip}{0pt}\setlength{\parfillskip}{0pt}\arraybackslash}p{\dimexpr\textwidth/8-2\tabcolsep\relax}>{\setlength{\rightskip}{0pt}\setlength{\parfillskip}{0pt}\arraybackslash}p{\dimexpr\textwidth/8-2\tabcolsep\relax}>{\setlength{\rightskip}{0pt}\setlength{\parfillskip}{0pt}\arraybackslash}p{\dimexpr\textwidth/8-2\tabcolsep\relax}>{\setlength{\rightskip}{0pt}\setlength{\parfillskip}{0pt}\arraybackslash}p{\dimexpr\textwidth/8-2\tabcolsep\relax}@{}}
\toprule
King James Bible & seen & followers & variety & Austen & seen & followers & variety \\
\midrule
\texttt{years,} & 32 & 1 & 0.031 & \texttt{Heading} & 29 & 1 & 0.034 \\
\texttt{king} & 25 & 1 & 0.040 & \texttt{Mrs.\allowbreak{}} & 73 & 17 & 0.233 \\
\texttt{said} & 133 & 11 & 0.083 & \texttt{Mr.\allowbreak{}} & 178 & 46 & 0.258 \\
\texttt{sons} & 53 & 7 & 0.132 & \texttt{Miss} & 119 & 32 & 0.269 \\
\texttt{And} & 788 & 120 & 0.152 & \texttt{I} & 354 & 107 & 0.302 \\
\bottomrule
\end{longtable}
\normalsize

Austen's three lowest, after one artifact, are the three honorifics. A title in a Regency novel precedes a small closed set of surnames, which makes the social system the book runs on visible as a statistic.

The bible's are its register and its formulas. \texttt{And} occurring 788 times with a variety of 0.152 is very low for a conjunction and is the paratactic style the translation is known for. \texttt{years,} and \texttt{king} each have exactly one follower across every occurrence. Those are the genealogy and the regnal formula.

Nothing here knows English, knows what a name or a title is, or has a dictionary of any kind.

\textbf{Two limits.} \texttt{Heading} is Project Gutenberg's markup and not Austen, the same class of artifact as the line wrapping in Section 4.13. And narrow company marks formulaic context in general, which returns stock phrases such as \texttt{said unto} alongside culturally weighted words and does not distinguish the two. Cultural weight is one of the things that produces this signature and not the only one.

\paragraph{4.14.1 Burstiness and Grammatical Connectives}

Narrow company catches connectives along with subjects, and a connective is the wrong answer twice over: it is needed everywhere, so it says nothing about what a text is about, and its form varies between languages while its job does not. Where a unit falls separates them. A unit doing grammatical work is spread evenly. The gaps between its occurrences are near geometric and their dispersion sits near one. A unit belonging to a subject arrives in bursts where that subject is discussed, and its dispersion runs far above one. That is Section 4.12's statistic read at its other end (\texttt{test/\allowbreak{}bench/\allowbreak{}bench\_\allowbreak{}ancorae\_\allowbreak{}ab.\allowbreak{}c:2187}).

\tiny
\setlength{\tabcolsep}{2pt}
\begin{longtable}{@{}>{\setlength{\rightskip}{0pt}\setlength{\parfillskip}{0pt}\arraybackslash}p{\dimexpr\textwidth/9-2\tabcolsep\relax}>{\setlength{\rightskip}{0pt}\setlength{\parfillskip}{0pt}\arraybackslash}p{\dimexpr\textwidth/9-2\tabcolsep\relax}>{\setlength{\rightskip}{0pt}\setlength{\parfillskip}{0pt}\arraybackslash}p{\dimexpr\textwidth/9-2\tabcolsep\relax}>{\setlength{\rightskip}{0pt}\setlength{\parfillskip}{0pt}\arraybackslash}p{\dimexpr\textwidth/9-2\tabcolsep\relax}>{\setlength{\rightskip}{0pt}\setlength{\parfillskip}{0pt}\arraybackslash}p{\dimexpr\textwidth/9-2\tabcolsep\relax}>{\setlength{\rightskip}{0pt}\setlength{\parfillskip}{0pt}\arraybackslash}p{\dimexpr\textwidth/9-2\tabcolsep\relax}>{\setlength{\rightskip}{0pt}\setlength{\parfillskip}{0pt}\arraybackslash}p{\dimexpr\textwidth/9-2\tabcolsep\relax}>{\setlength{\rightskip}{0pt}\setlength{\parfillskip}{0pt}\arraybackslash}p{\dimexpr\textwidth/9-2\tabcolsep\relax}>{\setlength{\rightskip}{0pt}\setlength{\parfillskip}{0pt}\arraybackslash}p{\dimexpr\textwidth/9-2\tabcolsep\relax}@{}}
\toprule
King James Bible & seen & burst & Cervantes & seen & burst & Austen & seen & burst \\
\midrule
\texttt{years,} & 32 & 20.8 & \texttt{Donde} & 48 & 18.5 & \texttt{you} & 205 & 4.3 \\
\texttt{begat} & 67 & 20.0 & \texttt{Sancho} & 52 & 9.7 & \texttt{I} & 354 & 3.6 \\
\texttt{lived} & 34 & 13.6 & \texttt{aventura} & 37 & 8.0 & \texttt{Bingley,} & 26 & 3.4 \\
\texttt{king} & 25 & 9.5 & \texttt{trata} & 25 & 7.8 & \texttt{had} & 179 & 3.0 \\
\texttt{waters} & 28 & 9.4 & \texttt{historia} & 25 & 4.0 & \texttt{Miss} & 119 & 2.8 \\
\texttt{Esau} & 28 & 8.3 & \texttt{dijo} & 51 & 3.9 & \texttt{he} & 195 & 2.5 \\
\bottomrule
\end{longtable}
\normalsize

\textbf{The connectives are gone.} Neither \texttt{the}, \texttt{and}, \texttt{of} nor \texttt{And} appears in any of the three lists, though each is among the most frequent units in its text. Being spread evenly puts them at the floor of this statistic. They are discarded without a stopword list, a grammar or a dictionary.

What comes back is the begetting and living of Genesis, Sancho and his adventures, and the pronouns of a novel carried by dialogue. Three languages and three centuries through one statistic.

\textbf{The size of the numbers carries something as well.} The bible bursts between 6.6 and 20.8 and the novel only between 2.5 and 4.3. A compilation of books written centuries apart has vocabulary that clusters hard by section, and a single novel is homogeneous enough that little clusters at all. The scale of the statistic reports how uniform a text is.

\subsection{5. Discussion}

\subsubsection{5.1 Algebraic Structure of the Filter}

Strip the cost away and what remains is small. The state is one bit for each alignment, alive or dead. A read of one cell selects a mask over the pattern's positions and clears the alignments whose position holds something else. Survivors are the complement. The method is that and no more, and it contains no comparison that steers control flow, no position computed from a value just read, no ordering on positions, no metric, and no dimension.

Proposition 1 is why the pieces commute: a refutation depends only on the cell that produced it, and a set of reads contributes a union and a union does not remember the order it was built in. The consequences are all mechanical from there. Reads can be issued together because none waits on another. Accumulation is a bitwise OR. The surviving set is a complement. The mask is \(m\) bits wide whatever the alphabet holds and however many dimensions the index set has. Neither quantity appears in the state and neither needs a bound.

Entropy is nowhere in that description. It enters one line later, when the question becomes how many reads are needed, and it never returns.

\subsubsection{5.2 Collision Entropy as Cost Parameter}

Section 4.0 lists eight quantities that are each a function of \(2^{-H_2}\) and the two sizes. They were derived at different times for different reasons and several were measured before anyone noticed they were the same parameter. The candidate rate, the refutation distance, the estimator's bias correction, the fraction an in-order walk collects, the offset where advance saturates, the number of anchors before the space is spent, the free order probe stride, and the candidate count in \(d\) dimensions.

The last of those is the one that makes the collection meaningful instead of coincidental. The same expression holds from a line to an eight dimensional hypercube and over an alphabet of complex numbers with irrational parts, which is where the parameter demonstrates it is not standing in for the geometry or for what a symbol happens to be.

A single scalar fixing a cost surface across every dimension and every alphabet would be a strong claim if it were unqualified. The claim is not unqualified, and the next three sections are the qualifications.

\subsubsection{5.3 Dependence on Symbolization}

There is no collision entropy of a corpus. There is a collision entropy of a corpus read \(w\) bits at a time, and Section 4.10 measures \(H_2(w)/w\) falling from 0.984 to 0.600 on English between one bit and sixteen. Every figure in this document uses \(w = 8\) because a byte is the unit the code happened to be written in.

The clearest evidence that the choice is doing work is the corpus this document calls flat. It is \texttt{0x41} repeated, its bit pattern has period eight, and a window of six bits or more therefore takes exactly eight values and carries \(\log_2 8 = 3\) bits. Five separate instruments report zero on it, in Sections 4.3, 4.4, 4.4.1 and 4.5.1, and their agreement looked like independent confirmation. All five read bytes. The object carries three bits one bit narrower, and the five agreements are one assumption held five times.

\subsubsection{5.4 Higher-Order Distributional Information}

Two results in this document depend on more of the distribution than its collision probability, and neither is predicted by it.

Section 4.3.1 gives what a measure can be worth as the uninformed rate over the expected minimum of \(m\) size biased draws, and that ratio is one exactly when the probabilities are equal on the support. It is a property of the order statistics. Entropy orders the corpora wrongly: \texttt{structured} has the lowest collision entropy of the four and the highest ceiling. Section 4.8 then prices a mismatched reference at 25 to 30 times, the largest effect measured anywhere here and one no marginal predicts.

\subsubsection{5.5 Dependence Violations}

Every expression in Section 5.2 treats anchors as independent. Four measurements say they are not.

The product rule overpromises by up to four orders of magnitude at six anchors on structured data, always in the direction that under-provisions. Correlation survives to separation seven, and a stride equal to a record period is the strongest effect in this study at \(z = 12.3\). The stack exhaustion count predicts 2.98 anchors and measures 6 at \(m = 16\), where sixteen offsets cannot be spread far enough apart. And the distance field, which wins 13.2\% elsewhere, is beaten on periodic data alone, because crediting every offset from one read assumes the symbol distribution does not depend on position.

Those are four faces of one thing. The parameter is a marginal and each failure is a statement about a joint, so the boundary is not an accident of these corpora. It is where a marginal stops being able to say anything.

\subsubsection{5.6 Transferable Results}

The product rule failure is the portable one. Any structure that stacks cheap filters and sizes itself by multiplying their rates inherits it on structured data, in the same direction, and under-provisions.

The estimator is a subsampled U statistic and loses to a histogram estimator wherever one can be built. Its interest is that it is computable when one cannot: no bins, no codebook, constant state, and free inside a search that was happening anyway.

And the shape of the argument transfers whether or not the algorithm does. Correctness was separated from cost at Proposition 1 and never rejoined it. Every claim in this document that survived is one where that separation held, and every claim that fell is one where a cost measurement was allowed to stand in for a property of the thing.

\subsection{6. Limitations and Validity}

Everything below is finite, and the point of this section is to say what finite means as a number.

\subsubsection{6.1 Finite Sample Domains}

Corpora are 1408 to 2048 bytes in the sift bench, 3375 to 4089 positions in the lattice bench, and 256 to 16384 symbols in the estimator bench. Needle lengths reach 2048 and anchor counts reach six. Proposition 1 quantifies over unbounded sets and the measurements do not. The invariant results are evidence that the implementation agrees with the proposition over the ranges named and are not evidence for the proposition itself, which needs none.

\subsubsection{6.2 Finite Carriers}

Symbols are drawn from an alphabet of at most 256 values in every bench, and the complex domain uses two. A claim that the construction tolerates an unenumerable alphabet is supported here by the core never reading a value (\texttt{test/\allowbreak{}bench/\allowbreak{}bench\_\allowbreak{}ancorae\_\allowbreak{}lattice.\allowbreak{}c:80}, \texttt{193-\allowbreak{}239}) and not by a measurement on an unenumerable alphabet, which cannot be performed.

\subsubsection{6.3 Finite Precision}

The complex domain's symbols are square roots of primes held in \texttt{double}, which makes the values stored rational approximations of irrationals. Equality is decided over storage, so the approximation affects which symbols are distinct and never whether the comparison is sound.

\subsubsection{6.4 Sensitivity to Entropy Estimation}

Sections 4.6.1, 4.9.2 and 4.9.4 derive constants that take \(H_2\) as an input. All three were computed here from a full histogram over the corpus, which presumes the alphabet can be enumerated. Section 2.1 refuses that assumption, and Section 4.5 measures what the alternative costs: estimating \(H_2\) from candidate counts carries a root mean square error between 0.02 and 0.60 bits depending on the probe budget.

Two of the three constants are exponential in \(H_2\), so that error does not stay small:

\tiny
\setlength{\tabcolsep}{2pt}
\begin{longtable}{@{}>{\setlength{\rightskip}{0pt}\setlength{\parfillskip}{0pt}\arraybackslash}p{\dimexpr\textwidth/4-2\tabcolsep\relax}>{\setlength{\rightskip}{0pt}\setlength{\parfillskip}{0pt}\arraybackslash}p{\dimexpr\textwidth/4-2\tabcolsep\relax}>{\setlength{\rightskip}{0pt}\setlength{\parfillskip}{0pt}\arraybackslash}p{\dimexpr\textwidth/4-2\tabcolsep\relax}>{\setlength{\rightskip}{0pt}\setlength{\parfillskip}{0pt}\arraybackslash}p{\dimexpr\textwidth/4-2\tabcolsep\relax}@{}}
\toprule
constant & sensitivity & at 0.25 bits & at 0.60 bits \\
\midrule
saturation radius \(3\cdot 2^{H_2}\) & \(\mathrm{d}r/r = \ln 2 \cdot \mathrm{d}H_2\) & 17\% & 51\% \\
harvest, through \(m\,2^{-H_2}\) & the same exponential & 17\% & 51\% \\
exhaustion \(\log_2 N / H_2\) & \(\mathrm{d}n/n = -\mathrm{d}H_2/H_2\) & 6.8\% & 16\% \\
\bottomrule
\end{longtable}
\normalsize

At a 16 probe budget the saturation radius is known to about half its value. For \texttt{periodic16} that moves the crossing from 73 to somewhere between 36 and 110, a range spanning two of the needle lengths tested. The cold and hot classification of a given search is not always decidable from a cheap estimate. The exhaustion count is the mildest of the three and stays inside 16\%.

None of this reaches Sections 2.2, 2.3 or 4.1 to 4.2. Soundness, the absence of completeness, and every invariant row are indifferent to whether the distribution is known at all. That property made them deductive in the first place.

\subsubsection{6.5 Corrections to the Analysis}

Recorded because each was invisible until something contradicted it.

A predicted rate that included the needle finding itself, inflating every ratio. A self similar corpus that produced a clean looking trend of correlation growing with needle length, which reversed once the corpora stopped repeating. A 1.24 control read as a finding when its \(z\) said it was noise. The corpus matched table prediction, refuted by Section 4.8. A uniform control summarized as a ratio when its counts were too small to carry one. A control generator changed from a 64 bit xorshift to SHA-256 in counter mode while its own comment still described the xorshift, which moved the uniform skip figures from 261.1, 262.7 and 261.1 to 271.9, 273.6 and 281.3 (\texttt{test/\allowbreak{}bench/\allowbreak{}bench\_\allowbreak{}ancorae\_\allowbreak{}sift.\allowbreak{}c:252-\allowbreak{}262}). A dimension sweep whose predicted column again omitted the guaranteed self match, drifting to a ratio of 1.22 at the smallest domain until the term was restored. That is the first correction in this list made a third time. An interpretation of Section 4.3 that named entropy as the quantity governing what a measure is worth, which the corpus ordering contradicts and Section 4.3.1 replaces. And a claim about prior art asserted from search engine snippets before either source had been read.

The first entry recurring three times is the finding, not the anecdote. A pattern drawn from the data it is searched in always matches itself, and every model written here has had to be told so separately.

\subsubsection{6.6 Unmeasured Quantities}

Any vectorized cost model. Every figure in Section 4.9 counts sequential reads, and a shift is inherently serial because the next position cannot be computed until the current one has been read. Under a model that tests many positions per word operation, the advance advantage that decides Section 4.9.1 may not be available at all, and rejection would be the only lever left. That measurement would decide whether the arrangement in Section 4.9.3 belongs in the library, and it has not been made.

The estimator's behavior against published estimators other than the two in Section 4.5. The construction on a domain whose alphabet is genuinely unenumerable. A spatially correlated domain of more than one dimension, on which a centroid based anchor rule would need to be tested, since on an independent domain every point of a pattern is interchangeable and such a rule correctly shows nothing. Any timing figure: what a SWAR word costs is a per part question and no host number answers it.

\subsection{7. Related Work}

Two sources were read in full. Both were reached because a search suggested they were close, and a search result is not a reading.

The \texttt{memchr} crate's \texttt{arch/\allowbreak{}all/\allowbreak{}packedpair} module selects two byte positions by a frequency rank and uses the pair as a prefilter. The only constraint that module places on where the two sit is that their indices differ, and nothing in it treats their separation as a variable or their rates as independent. That statement covers the one module read. It does not describe the crate, whose substring search, rare byte prefilter and vector backends are separate files that were not read, and no claim is made about them.

The arXiv preprint at \texttt{2601.\allowbreak{}03271} selects a single anchor, \(a = \arg\min_i \mathrm{freq}(P[i])\), and reports character comparisons on the \_Divina Commedia\_. Having one anchor, it contains no separation, independence or joint probability treatment, and uses no periodic or fixed width data.

Neither document contains the separation analysis here. That is a claim about two documents and not about the literature, and it is the strongest form the evidence supports.

\subsubsection{7.1 Order-Free Filtering and Deterministic Sampling}

Section 2.2 makes each refutation depend only on the cell that produced it. The set of alignments ruled out by a group of reads is a union and does not depend on the order they were read in. That observation is the dueling paradigm, introduced by Vishkin, and its consequences were worked out at the time. Read from Neuburger's survey, sections 1.3 and 1.4.

A duel eliminates candidates pairwise from single text symbols with no traversal order, and the survey states the setting directly: "This alphabet independent method is used for fast pattern matching in a parallel setting." Alphabet independence, which Section 2.1 assumes, is a stated property of the 1990 method.

Deterministic sampling then chooses probe positions from the pattern before any text is read. "The deterministic sample is a set A of positions of the pattern and a number \(x\), such that if the positions match a candidate \(i\) in the text, then all other candidates in the interval \([i-x+1, i+\tfrac{m}{2}-x]\) are eliminated." Its sample size is logarithmic in the pattern length, which is a stronger result than any fixed stride over the text.

Rare symbol selection appears inside that construction. The pattern is stacked \(\tfrac{m}{2}\) times, a witness column is chosen, and the rule is to "select the less frequent of the two symbols, which can occur in at most half of the copies," repeated at most \(\log \tfrac{m}{2}\) times. Choosing by rarity to halve a candidate set inside an order free elimination is the 1991 method.

Periodicity is a stated precondition there, with a published repair. The guarantee holds "for a non-periodic pattern," and a periodic pattern is reduced by taking its prefix of size \(2\pi-1\), which is not periodic, then linking neighboring occurrences. The failure recorded in Section 4.9.3, peaking at needle lengths equal to the record period and twice it, is that precondition being violated, and the repair for it predates this document.

\textbf{What that leaves, stated more carefully than a first reading of it allowed.} The prior art above is combinatorial and worst case. It selects a rarer symbol without pricing what rarity is worth, and it carries no distribution over the alphabet.

The prior art also carries a precondition that the construction in Section 2 does not. The deterministic sample is built from the pattern before any text is read: the pattern is stacked, a witness column is chosen, the less frequent of two symbols is taken, and the procedure repeats over the pattern's own self overlap structure. That is a preprocessing pass requiring the pattern in hand and analyzable in advance, and it buys a worst case guarantee at \(\log m\) probes.

An adaptive search declines that pass. It has no probe set until answers start arriving and no guarantee either, and Section 4.9.6 measures what that costs: probing without confirming is wrong on a third to almost all of the searches at any stride above one. What it has instead is that the error is one directional. A miss is impossible under Proposition 1. More probing can only reduce the error and never introduce one, and accuracy is monotone in effort with no preprocessing and no bound on how far it can be taken.

Section 4.9.6 also measures why the two cannot be traded freely. The safe stride is a property of the pattern, and a per pattern preprocessing pass would determine it. A stride calibrated on one pattern is wrong on 0 to 62 of 63 others. A fixed probe set needs the per pattern analysis, and skipping the analysis while keeping a fixed probe set is the arrangement that fails.

\subsubsection{7.3 Average-Case String Matching}

Read from Tsai, \_Average Case Analysis of the Boyer-Moore Algorithm\_, Random Structures and Algorithms 28, 481 (2006), pages 1 to 4 and 14 to 16 of 18. Sections 4 and 5, which carry the Markov chain derivation, were not read.

The quantity that literature computes is \(C_n\), the number of character comparisons, and the constant is \(\mu = \lim_{n\to\infty} \mathbb{E}(C_n/n)\). Baeza-Yates and colleagues "applied an analytic approach to the average-case analysis of the BMH algorithm. Under the assumption that the text is independent and identically distributed (iid), they derived an exact expression for the linearity constant". Mahmoud, Smythe and Régnier established asymptotic normality of \(C_n\), Smythe extended the text model to Markov, and Tsai adds a Berry-Esseen bound.

\textbf{The parameter is alphabet cardinality.} The text is written over "a set of \(q\) characters", every numeric result is computed at \(q = 2\), 4 and 8, the error term is \((m-K-1)/q^K\), and the experiment uses text "built randomly". Under a uniform alphabet \(\sum_\sigma p_\sigma^2 = 1/q\) exactly, so \(2^{H_2} = 1/\sum p^2\) is the effective alphabet size and the collision probability used throughout this document is that \(q\) carried to the non-uniform case. The transition probabilities in their framework are written for a general distribution, so the generalization is available there and is not taken. This document does not claim a new quantity, it claims the classical one at a weaker assumption.

\textbf{A comparison count cannot see a declined refutation.} Section 4.9.4 prices what an in-order shift leaves uncollected, between 3\% and 99\% depending on the needle length. That quantity is invisible to \(C_n\), because a comparison count records what the algorithm did and a forfeited refutation is something it never did. The gap is not a defect in thirty years of analysis. It is outside what the metric measures, and it becomes visible only once the traversal order is treated as a free variable, which Section 7.1 credits to Vishkin.

\subsubsection{7.2 Thermodynamic Information Bounds}

Sections 4.3 and 4.8 reach for a physical analogy: a reference measure is worth the divergence of the domain from it, and at equilibrium there is no gradient and nothing to extract. Read from Parrondo, Horowitz and Sagawa, \_Thermodynamics of information\_, Nature Physics 11, 131 (2015).

The results are sharper than the analogy. For a system in a statistical state \(\rho\) with Hamiltonian \(\mathcal{H}_0\), the non-equilibrium free energy is \(\mathcal{F}(\rho;\mathcal{H}_0) \equiv \langle \mathcal{H}_0\rangle_\rho - TS(\rho)\), and Box 1 states that "\(-W_{\text{diss}}\) is the maximum work that can be extracted from the non-equilibrium state \(\rho\)", with \(\mathcal{F}(\rho;\mathcal{H}_0) \ge \mathcal{F}(\rho_0)\) for any \(\rho\). The second law for isothermal processes between non-equilibrium states is \(T\Delta S_{\text{tot}} = W_{\text{diss}} \equiv W - \Delta\mathcal{F} \ge 0\).

For measurement the bound is on mutual information, not on a divergence from an arbitrary reference: \(\Delta \mathcal{F}_{\text{meas}} = kTI(X;M)\), and for a cyclic feedback process \(W \ge -kTI(X(t_{\text{ms}});M)\), so the extractable work is "proportional to the information acquired in the measurement". The Szilárd engine saturates it because its outcome is left or right with equal probability, "giving \(H(M) = \ln 2\)".

\textbf{Where this document departs from the analogy.} The equilibrium reading holds: a domain with no structure offers nothing, and Sections 4.3 and 4.8 measure exactly that in the rows that refuse to separate. The reference measure reading does not survive Section 4.3.1. What a measure can be worth here is the uninformed rate over the expected minimum of \(m\) size biased draws, which equals one when the distribution is uniform on its support. That quantity is not a relative entropy between the domain and the table, and the corpus ordering shows entropy does not govern it. The physics is invoked as background and no result in this document rests on it.

\subsubsection{7.4 Statistical Identification of Language}

Section 4.13 measures a Zipf slope, a brevity correlation and a boundary regularity on ten corpora and reads their agreement as a property of how people produce language. Exactly that inference has been argued over in print. Four items were read in full: Sproat, \_Ancient Symbols, Computational Linguistics, and the Reviewing Practices of the General Science Journals\_, Computational Linguistics 36(3):585, and in 36(4) the replies from Rao and colleagues at 795, from Lee and colleagues at 791, and Sproat's answer to both at 807.

Sproat's case is that these statistics do not separate language from anything else, and he demonstrates it instead of asserting it. Applying the published classifier of Lee and colleagues to Mesopotamian deity symbols from kudurru boundary stones, a system known to be non-linguistic, returns \(C_r = 8.0\) and \(U_r = 1.55\) and classifies it as writing. Applying it to 75 "texts" produced by successive tosses of seven six-sided dice returns \(C_r = 12.64\) and \(U_r = 1.18\), which the same decision tree reports as a syllabic writing system. For the conditional entropy measure he notes that a memoryless process with a Zipfian and non-equiprobable unigram distribution, carrying conditional independence between symbols, reproduces the published curves. His conclusion is that neither result distinguishes "structure that derives from linguistic constraints from structure that derives from some other kind of constraints."

Rao and colleagues answer that sufficiency was never claimed. Their framework is inductive, estimating a posterior over a linguistic hypothesis from several properties at once, of which entropy is one, and they extend the measurement to block entropies up to \(N = 6\) where they report the artificial counterexamples failing to track. They also supply non-linguistic controls that are not artificial, namely DNA, protein, Fortran and music, which occupy entropic ranges away from the linguistic ones. Their own description of the Zipf-Mandelbrot property is that it is "often considered a necessary (though not sufficient) condition for language", the asymmetry between Proposition 1 and Proposition 2 reached from a different direction and by different people.

\textbf{What this costs Section 4.13, measured and not conceded in the abstract.} Sproat's counterexample class was built and run through the same code that produced every other row (\texttt{maint/\allowbreak{}data/\allowbreak{}fetch/\allowbreak{}gen\_\allowbreak{}monkey\_\allowbreak{}corpus.\allowbreak{}py}). Characters are drawn independently and a delimiter is dropped in at a fixed rate.

A first version of that control had two arms, one uniform and one weighted by English letter frequencies, with the alphabet fixed at 26 and the delimiter rate at 0.18. Three of those four numbers came from English, and a result resembling English established nothing. The sweep below varies the alphabet from 8 to 64 symbols, the delimiter rate from 0.06 to 0.35, and the letter distribution across uniform, two geometric falls and the English weights, with only the last row carrying any value taken from a language.

\tiny
\setlength{\tabcolsep}{2pt}
\begin{longtable}{@{}>{\setlength{\rightskip}{0pt}\setlength{\parfillskip}{0pt}\arraybackslash}p{\dimexpr\textwidth/4-2\tabcolsep\relax}>{\setlength{\rightskip}{0pt}\setlength{\parfillskip}{0pt}\arraybackslash}p{\dimexpr\textwidth/4-2\tabcolsep\relax}>{\setlength{\rightskip}{0pt}\setlength{\parfillskip}{0pt}\arraybackslash}p{\dimexpr\textwidth/4-2\tabcolsep\relax}>{\setlength{\rightskip}{0pt}\setlength{\parfillskip}{0pt}\arraybackslash}p{\dimexpr\textwidth/4-2\tabcolsep\relax}@{}}
\toprule
measure & ten memoryless arms & ten natural corpora & separates \\
\midrule
Zipf slope & -0.978 to -1.332 & -0.723 to -1.185 & no \\
brevity correlation & -0.124 to -0.401 & -0.247 to -0.364 & no \\
types at 25k & 5824 to 22160 & 4171 to 11854 & no \\
dispersion against a permutation null & 0.99 to 1.01 & 1.79 to 4.82 & yes \\
\bottomrule
\end{longtable}
\normalsize

Three of the four measures fail. The Zipf slope from a geometric distribution carrying no language at all is -0.996, against -0.988 for the English weighted arm. The seeding was never what produced a natural looking value and non-uniformity alone is enough. This is the distinction between random and random-and-equiprobable that Sproat says the original work confused, and the uniform arms do sit outside the natural range at -1.224 and -1.332. The brevity correlation does something worse than fail to separate: an arm over 8 symbols with a delimiter rate of 0.35 reaches -0.401, stronger than any natural text measured here. The type count overlaps at three arms.

Section 4.13 therefore may not read a Zipf slope, a brevity correlation or a type count as evidence of anything about production, and those claims are withdrawn. What remains of that section is the boundary regularity.

That one separates completely, and across the whole sweep. Every memoryless arm scores between 0.99 and 1.01, which is a corpus indistinguishable from its own shuffle and is what a process with no memory has to give, whatever its alphabet size or delimiter rate. Every natural corpus scores at or above 1.79. This is the measure Sproat grants, writing of the randomization method that it "seems to make sense as a way to spot memoryless non-equiprobable processes masquerading as structured systems", and his stated limit on it applies here unchanged: a non-linguistic system that has structure passes it as well. So the surviving measure establishes that a corpus is not memoryless. It does not establish that a corpus is a language, and this document does not claim it does.

The surviving measure is also the one Section 7.2 can price. A corpus equal to its own shuffle carries no gradient and offers nothing to extract. The equilibrium reading holds exactly here, and the ratio above one is a distance from that equilibrium.

\textbf{What a cipher does to the surviving measure.} Sproat's stated limit is that a structured non-linguistic system passes this test as well. The question is what it takes to make a text stop passing it. One English corpus was put through transforms that keep every symbol in the same seat range, so the same measurement path reads all of them (\texttt{data/\allowbreak{}transform\_\allowbreak{}corpus.\allowbreak{}py}).

\tiny
\setlength{\tabcolsep}{2pt}
\begin{longtable}{@{}>{\setlength{\rightskip}{0pt}\setlength{\parfillskip}{0pt}\arraybackslash}p{\dimexpr\textwidth/4-2\tabcolsep\relax}>{\setlength{\rightskip}{0pt}\setlength{\parfillskip}{0pt}\arraybackslash}p{\dimexpr\textwidth/4-2\tabcolsep\relax}>{\setlength{\rightskip}{0pt}\setlength{\parfillskip}{0pt}\arraybackslash}p{\dimexpr\textwidth/4-2\tabcolsep\relax}>{\setlength{\rightskip}{0pt}\setlength{\parfillskip}{0pt}\arraybackslash}p{\dimexpr\textwidth/4-2\tabcolsep\relax}@{}}
\toprule
transform & key length & ratio & mean gap \\
\midrule
none &  & 2.91 & 5.29 \\
one fixed permutation of the seats & 1 permutation & 2.91 & 5.29 \\
repeating key & 1 & 2.91 & 5.29 \\
repeating key & 2 & 1.51 & 10.47 \\
repeating key & 3 & 1.17 & 15.10 \\
repeating key & 4 & 1.06 & 25.16 \\
full length pseudorandom addend & 728751 & 1.01 & 29.66 \\
\bottomrule
\end{longtable}
\normalsize

A permutation of the symbols reproduces the measurement to four decimal places, and the boundary is found at a different seat carrying an identical dispersion of 0.2815. Relabeling every symbol changes nothing the measure reads. A bijection has to do that.

A repeating key of length \(k\) sends one plaintext symbol to \(k\) ciphertext symbols by position, and the mean gap grows with \(k\) as that predicts. The measurement is being divided and not destroyed, and by \(k = 4\) the split gaps exceed what a statistic over gaps can read. The pattern is still recoverable there by separating the positions that share a key offset, and a ratio near one in those rows is a limit of this measure and is not evidence that anything was erased. An attempt to recover it that way is reported in the next paragraph and failed for an unrelated reason.

The last row is different in kind. When the key is as long as the message the ratio reaches 1.01, the floor the memoryless controls occupy. That is the one case where the structure is absent from the text instead of hidden in it, because it now resides in the key, and it is the condition for perfect secrecy. A bijection cannot remove what this measure reads unless it spends key material equal to the message.

One check on that was attempted and does not support anything. Taking every eighth symbol of the \(k = 8\) cipher isolates one key offset, which should leave a pure substitution, and it returns 0.95. The word boundary in this corpus sits at a period near 5.3 symbols and sampling at a stride of 8 is above that period. The subsampling destroys the boundary regularity on its own, with or without a cipher. The row measures the sampling and not the transform.

\paragraph{7.4.1 Frequency-Stratified Information}

Every result above comes from a detector that returns one symbol, the one whose gaps are most regular, and that rejects any candidate occurring less often than once in 64 symbols. Two consequences follow from that construction and neither was intended. The occurrence floor admits only frequent symbols, so the detector can only ever report on the head of the distribution. And ranking by regularity treats variability as failure, and a symbol whose occurrences cluster is scored as a poor candidate.

Under a Zipf distribution the head carries the token count and the tail carries the information, because the surprisal of a symbol is \(-\log p\) and the many rare symbols each contribute more of it. Scoring every symbol against the same permutation null instead of one (\texttt{data/\allowbreak{}measure\_\allowbreak{}spectrum.\allowbreak{}py:34-\allowbreak{}51}) gives the mean ratio over each half of the symbols by frequency:

\tiny
\setlength{\tabcolsep}{2pt}
\begin{longtable}{@{}>{\setlength{\rightskip}{0pt}\setlength{\parfillskip}{0pt}\arraybackslash}p{\dimexpr\textwidth/3-2\tabcolsep\relax}>{\setlength{\rightskip}{0pt}\setlength{\parfillskip}{0pt}\arraybackslash}p{\dimexpr\textwidth/3-2\tabcolsep\relax}>{\setlength{\rightskip}{0pt}\setlength{\parfillskip}{0pt}\arraybackslash}p{\dimexpr\textwidth/3-2\tabcolsep\relax}@{}}
\toprule
corpus & head half & tail half \\
\midrule
\texttt{english\_\allowbreak{}1611\_\allowbreak{}kjv} & 0.96 & 0.59 \\
\texttt{finnish\_\allowbreak{}1870\_\allowbreak{}kivi\_\allowbreak{}prose} & 0.97 & 0.62 \\
\texttt{french\_\allowbreak{}1862\_\allowbreak{}hugo} & 0.96 & 0.69 \\
\texttt{english\_\allowbreak{}1813\_\allowbreak{}austen} & 0.95 & 0.73 \\
\texttt{greek\_\allowbreak{}iliad} & 1.06 & 0.75 \\
\texttt{monkey\_\allowbreak{}a26\_\allowbreak{}d18\_\allowbreak{}geom90} & 1.00 & 1.00 \\
\bottomrule
\end{longtable}
\normalsize

The tail departs from the null by 0.25 to 0.41 where the head departs by 0.03 to 0.06, and the memoryless control sits at 1.00 in both halves. The departure is a property of the corpora and not of the measure. The direction is below one: the real corpus carries gaps more variable than its shuffle. That is clustering: a rare word appears several times in one passage and nowhere else. Clustering is structure, and the ranking used everywhere above scores it as the opposite.

The averages also separate better than the winner does. A head average of 0.95 to 1.06 against a control at 1.00 does not separate at all, and every separation reported earlier in this section came from the single best head symbol, which has to be found. The tail average separates from the control across all five corpora without selecting any symbol.

So the measure that survived Section 7.4 is the weaker of two, and the stronger one was excluded by an occurrence floor added in Section 4.13.0 to reject decoration. Nothing above is withdrawn, since the boundary results stand as measured. The withdrawal covers only the implication that the boundary is where the structure sits.

\paragraph{7.4.2 Invariance and Periodic Obfuscation}

The transforms in Section 7.4 leave the tail measure exactly where the head measure was left, and running them through it characterizes the measure completely. Cosets are scored with \texttt{data/\allowbreak{}supersample\_\allowbreak{}cosets.\allowbreak{}py}, which splits a corpus at a stride, scores every coset and averages, reading all of the text and spending none of the length.

A relabeling of the symbols changes nothing, because the measure reads gaps between occurrences and a permutation moves which seat carries which gaps. The single symbol result is identical to four decimals at 0.2815, and over the whole spectrum a substitution is likewise invisible.

A repeating key of length 8 divides the structure without removing it. Read whole, the ciphertext gives a tail figure of 0.99 against 0.73 for the plaintext, which is a corpus that looks memoryless. Split at a stride of 8 and averaged over all eight cosets, the ciphertext gives 0.896 and the plaintext gives 0.896. The agreement is exact because each coset was enciphered by one substitution, and a substitution is what the previous paragraph shows the measure cannot see.

The period does not have to be known. Scanning the stride over two ciphertexts of the same plaintext, one with a key length of 3 and one with 8, against the plaintext scanned the same way:

\tiny
\setlength{\tabcolsep}{2pt}
\begin{longtable}{@{}>{\setlength{\rightskip}{0pt}\setlength{\parfillskip}{0pt}\arraybackslash}p{\dimexpr\textwidth/11-2\tabcolsep\relax}>{\setlength{\rightskip}{0pt}\setlength{\parfillskip}{0pt}\arraybackslash}p{\dimexpr\textwidth/11-2\tabcolsep\relax}>{\setlength{\rightskip}{0pt}\setlength{\parfillskip}{0pt}\arraybackslash}p{\dimexpr\textwidth/11-2\tabcolsep\relax}>{\setlength{\rightskip}{0pt}\setlength{\parfillskip}{0pt}\arraybackslash}p{\dimexpr\textwidth/11-2\tabcolsep\relax}>{\setlength{\rightskip}{0pt}\setlength{\parfillskip}{0pt}\arraybackslash}p{\dimexpr\textwidth/11-2\tabcolsep\relax}>{\setlength{\rightskip}{0pt}\setlength{\parfillskip}{0pt}\arraybackslash}p{\dimexpr\textwidth/11-2\tabcolsep\relax}>{\setlength{\rightskip}{0pt}\setlength{\parfillskip}{0pt}\arraybackslash}p{\dimexpr\textwidth/11-2\tabcolsep\relax}>{\setlength{\rightskip}{0pt}\setlength{\parfillskip}{0pt}\arraybackslash}p{\dimexpr\textwidth/11-2\tabcolsep\relax}>{\setlength{\rightskip}{0pt}\setlength{\parfillskip}{0pt}\arraybackslash}p{\dimexpr\textwidth/11-2\tabcolsep\relax}>{\setlength{\rightskip}{0pt}\setlength{\parfillskip}{0pt}\arraybackslash}p{\dimexpr\textwidth/11-2\tabcolsep\relax}>{\setlength{\rightskip}{0pt}\setlength{\parfillskip}{0pt}\arraybackslash}p{\dimexpr\textwidth/11-2\tabcolsep\relax}@{}}
\toprule
stride & 2 & 3 & 4 & 5 & 6 & 7 & 8 & 9 & 12 & 16 \\
\midrule
plaintext & 0.800 & 0.830 & 0.876 & 0.889 & 0.895 & 0.895 & 0.896 & 0.897 & 0.896 & 0.941 \\
key length 3 & 0.909 & 0.830 & 0.935 & 0.952 & 0.895 & 0.975 & 0.950 & 0.897 & 0.896 & 0.984 \\
key length 8 & 0.963 & 1.002 & 0.913 & 1.004 & 0.977 & 1.010 & 0.896 & 1.007 & 0.958 & 0.941 \\
\bottomrule
\end{longtable}
\normalsize

Each ciphertext equals the plaintext exactly at every stride that is a multiple of its key length, at 3, 6, 9 and 12 for the first and at 8 and 16 for the second, and lies above the plaintext everywhere else. The cipher is not producing a dip at the period. It is withholding the plaintext's own structure at every stride except the ones that align the cosets with the key, where it withholds nothing at all.

The plaintext row makes that readable, and taking a scan without it invites two errors. The baseline is not flat, climbing from 0.800 at a stride of 2 to 0.941 at 16 as the subsampling removes structure on its own, and a dip has to be judged against that curve and not against 1.0. And the value coprime strides take is a property of the key length: with a key of 8 they sit near 1.0 because the structure is divided eight ways, while with a key of 3 they sit between 0.909 and 0.975 because a three way division hides less.

The baseline is not an artifact of the shortening. Each coset holds one stride's worth of the corpus, which makes a scan compare estimates from different amounts of text and from whatever symbols cleared the occurrence floor at that length. Truncating every coset to 45000 symbols removes both, and the curve still climbs, from 0.769 at a stride of 2 to 0.936 at 16.

Its shape is \(1 - c/\sqrt{s}\) in the stride \(s\). Fitting the capped curve gives \(c\) between 0.237 and 0.277 at eight of the ten strides measured, with strides 2 and 12 at 0.327 and 0.343. The exponent follows from what the tail measure reads. Its signal is rare symbols clustering into passages, and at a fixed coset length a larger stride spans proportionally more of the original text. Each cluster contributes \(1/s\) as many points to the coset. Deciding that a cluster is present is a counting problem and counting significance grows as the square root of the count, so the signal falls as \(1/\sqrt{s}\).

That sets what a scan can reach. Structure is not lost at a stride, it is attenuated by a known factor, and a longer corpus buys back any stride at quadratic cost.

\paragraph{7.4.3 Structural Ranking of Formal and Natural Languages}

Sproat's stated limit is that a system with structure and no language passes this test, and every control in Section 7.4 so far has been memoryless, which tests only the other side. The C sources in this repository were measured as a corpus, 422887 bytes over 94 distinct symbols, since a programming language is produced under constraints that share nothing with a vocal tract or a listener decoding in real time.

\tiny
\setlength{\tabcolsep}{2pt}
\begin{longtable}{@{}>{\setlength{\rightskip}{0pt}\setlength{\parfillskip}{0pt}\arraybackslash}p{\dimexpr\textwidth/3-2\tabcolsep\relax}>{\setlength{\rightskip}{0pt}\setlength{\parfillskip}{0pt}\arraybackslash}p{\dimexpr\textwidth/3-2\tabcolsep\relax}>{\setlength{\rightskip}{0pt}\setlength{\parfillskip}{0pt}\arraybackslash}p{\dimexpr\textwidth/3-2\tabcolsep\relax}@{}}
\toprule
corpus & head half & tail half \\
\midrule
C source & 0.59 & 0.50 \\
\texttt{english\_\allowbreak{}1813\_\allowbreak{}austen} & 0.95 & 0.73 \\
\texttt{monkey\_\allowbreak{}a26\_\allowbreak{}d18\_\allowbreak{}geom90} & 1.00 & 1.00 \\
\bottomrule
\end{longtable}
\normalsize

The C source departs from the null further than the English does on both halves. So the measure orders these three by how much structure they carry, and a formal language carries more of it than prose does, and fixed syntax, repeated identifiers and aligned indentation produce that. The measure is a structure meter. Nothing in it identifies a natural language, and the earlier statement that it establishes only that a corpus is not memoryless is now measured on the instrument itself instead of taken from the literature.

One reading survives this and is not tested by it. Both corpora that score away from the null were written by people, which leaves the result consistent with the measure detecting production by a person while failing to distinguish which kind. Separating those needs a system carrying structure that no person produced, of the sort Rao and colleagues reach for with genomes and protein sequences. No such control is measured here, so no claim of that kind is made.

\paragraph{7.4.4 Scale Localization of Structure}

The null used above shuffles single symbols, which destroys every arrangement at once and cannot say which span an arrangement occupied. Cutting the corpus into blocks of \(B\) symbols and shuffling the blocks keeps everything shorter than \(B\) and destroys everything longer (\texttt{data/\allowbreak{}block\_\allowbreak{}shuffle\_\allowbreak{}null.\allowbreak{}py:52-\allowbreak{}59}), so sweeping \(B\) locates the signal.

\tiny
\setlength{\tabcolsep}{2pt}
\begin{longtable}{@{}>{\setlength{\rightskip}{0pt}\setlength{\parfillskip}{0pt}\arraybackslash}p{\dimexpr\textwidth/8-2\tabcolsep\relax}>{\setlength{\rightskip}{0pt}\setlength{\parfillskip}{0pt}\arraybackslash}p{\dimexpr\textwidth/8-2\tabcolsep\relax}>{\setlength{\rightskip}{0pt}\setlength{\parfillskip}{0pt}\arraybackslash}p{\dimexpr\textwidth/8-2\tabcolsep\relax}>{\setlength{\rightskip}{0pt}\setlength{\parfillskip}{0pt}\arraybackslash}p{\dimexpr\textwidth/8-2\tabcolsep\relax}>{\setlength{\rightskip}{0pt}\setlength{\parfillskip}{0pt}\arraybackslash}p{\dimexpr\textwidth/8-2\tabcolsep\relax}>{\setlength{\rightskip}{0pt}\setlength{\parfillskip}{0pt}\arraybackslash}p{\dimexpr\textwidth/8-2\tabcolsep\relax}>{\setlength{\rightskip}{0pt}\setlength{\parfillskip}{0pt}\arraybackslash}p{\dimexpr\textwidth/8-2\tabcolsep\relax}>{\setlength{\rightskip}{0pt}\setlength{\parfillskip}{0pt}\arraybackslash}p{\dimexpr\textwidth/8-2\tabcolsep\relax}@{}}
\toprule
block span & 1 & 8 & 32 & 128 & 512 & 2048 & 8192 \\
\midrule
\texttt{english\_\allowbreak{}1813\_\allowbreak{}austen} tail & 0.725 & 0.727 & 0.729 & 0.811 & 0.846 & 0.934 & 0.975 \\
C source tail & 0.500 & 0.524 & 0.550 & 0.623 & 0.713 & 0.825 & 0.915 \\
\texttt{monkey\_\allowbreak{}a26\_\allowbreak{}d18\_\allowbreak{}geom90} tail & 0.997 & 1.001 & 1.001 & 0.998 & 0.999 & 1.001 & 1.000 \\
\bottomrule
\end{longtable}
\normalsize

The memoryless control reads 1.0 at every span. The sweep introduces nothing by itself.

The English text is flat from a span of 1 to a span of 32 and then climbs. Blocks of 32 symbols, around six words, destroy its tail signal as completely as shuffling single symbols does, and 1.5\% of the departure from the null has returned by that point. Its structure begins returning at 128 and is still returning at 8192, which is a passage.

The C source returns 10\% of its departure by a span of 32, close to seven times as much, and climbs from a span of 8 upward without a flat region. So the two carry this signal at different spans: the source has arrangements at every scale, including an identifier reused a few lines later, while the prose has almost none of it below a sentence and all of it above one.

That is a statement about the clustering measure and not about the text. Section 4.6 measures a cascade of three anchors inside a needle of 24 symbols surviving 20992 times more often than the product of the anchor rates predicts, which is correlation between positions closer together than 24 and far larger than anything in this table. Both hold. The block sweep says the clustering signal has no component below a span of 32, and the cascade says a different arrangement lives below 24 that the clustering measure cannot see.

This is the first measurement in this document that orders a natural language differently from a formal one, and it does so by where the structure sits and not by how much of it there is. A reading of it is available and is not established here: a rare word recurs across a passage because the passage concerns what the word names, which is an arrangement no shorter than the subject being discussed. Nothing in this measurement reaches what a text is about, and the span is the only quantity being reported.

The corresponding figures for the frequent half move non-monotonically over the same sweep, from 0.945 at a span of 1 to 0.913 at 32 and back to 0.993 at 8192. That half is a mean over many symbols carrying little structure individually, which leaves it too diluted to support a claim and none is made from it.

The minimum over the scan is the key length in both cases, at 0.830 for the first and 0.896 for the second. Following a doubling ladder finds the second and misses the first, since 3 is not reached by doubling, and the harmonic cluster carries traps in both directions: a stride of 12 reads 0.958 against the key of 8 because it shares a factor of 4 with it, which is deeper than the correct divisor 2 reads at 0.963. Taking the first strong dip while scanning upward returns 12 for a key of 8. The scan has to be read as a whole.

One transform defeats all of it. A pseudorandom addend as long as the message gives 1.004 read whole and 1.004 under every coset scan, because there is no period to split on and no coset in which a fixed substitution was applied. The memoryless controls give 1.005 by the same procedure.

So the measure is invariant under relabeling, divided but not destroyed by any periodic key, and recoverable from that division without knowing the period. It is defeated only where key material equals message length. That is a complete statement of what it detects, and it is the same statement in both directions: the structure this document keys on cannot be removed by any transform that a shorter description could undo.

\textbf{The same finding points the other way for the method.} Section 4.13.08 argues that these regularities are regenerated by the conditions that produce a message and are not carried between instances. Sproat's position is the stronger form of that: the structure arises from generators with no communicative intent at all, dice included. A method keying on a property that nearly every generator produces has a wider domain than one keying on a property special to language, and Section 2.1 never required the source to be a language. It requires that something regular produced the bytes.

\subsection{8. Conclusion}

The safety of an anchor filter and the saving from one are different kinds of claim and were measured as such. The first needs no data and was confirmed to hold in the implementation over 9,396,207 true occurrences on byte strings and 213,840 more across thirteen other geometries, with none refused. The second is a property of a domain, moves by up to a factor of 6.4 between measures, and has a ceiling that equals one exactly when the domain's distribution is uniform on its support.

One quantity turned out to do three jobs. The collision probability \(2^{-H_2}\) sets the rate at which an uninformed anchor admits a candidate, sets the fraction of alignments a single read settles, and sets the size of the bias correction separating the two estimator forms. None of the three was derived from the others and all three were measured independently.

Three results carry outside this construction. The candidate count is an estimator of collision entropy that needs no histogram and no codebook, at an accuracy cost of one to twenty times against published estimators depending on the probe budget. Sizing a cascade by multiplying anchor rates underestimates the survivors by up to four orders of magnitude on structured data, always in the direction that under-provisions. And a filter of this shape settles \(m(1-2^{-H_2})\) alignments per read instead of one, which is a property of the necessary condition and not of any particular shift table.

Where an order on positions exists, that order is most of the available structure and a shift takes it. The candidate rate this document spends its length measuring is not what governs a search. Section 4.9 records that as a loss of 1.2 to 25 times and then recovers it, first by treating rejection and advance as one product and then by measuring the advance directly. The arrangement that wins holds no model of the alphabet at all, the only form of it Section 2.1's assumptions permit.

\subsection{9. Reproducibility}

Every figure comes from four self contained programs and one cost profile. Build and run them with:

\begin{lstlisting}[breaklines=true,breakatwhitespace=false,basicstyle=\ttfamily\small]
gcc -O2 -std=c11 -Itest/bench -Isrc -Iinclude -Ideps/embedded_types/include -Itest/support \
    test/bench/bench_ancorae_sift.c test/support/mmgr_sha256.c \
    src/impensa_ancorae_acus/impensa_ancorae_acus_english.c -lm -o sift
\end{lstlisting}

The same line builds \texttt{bench\_\allowbreak{}ancorae\_\allowbreak{}lattice.\allowbreak{}c}, \texttt{bench\_\allowbreak{}ancorae\_\allowbreak{}entropy.\allowbreak{}c} and \texttt{bench\_\allowbreak{}ancorae\_\allowbreak{}ab.\allowbreak{}c}. Every sift row carries a fingerprint of the 256 linked costs, \texttt{69c2e2df} for the english profile, because the cost table is a link time singleton and a binary reports on one profile only. The estimator bench takes \texttt{-\allowbreak{}DPROBE\_\allowbreak{}SAMPLES=k} to reproduce the budget sweep in Section 4.5.

\subsection{Appendix A: Digest Oracle}

\texttt{mmgr\_\allowbreak{}sha256} is test support and nothing in \texttt{src/\allowbreak{}} reaches it. A digest computed by the code under test agrees with that code whatever the code did, and testing precept 1 in this tree is that the library is never its own oracle.

\texttt{mmgr\_\allowbreak{}sha256} is held to five published vectors by \texttt{mmgr\_\allowbreak{}sha256\_\allowbreak{}self\_\allowbreak{}test} (\texttt{test/\allowbreak{}support/\allowbreak{}mmgr\_\allowbreak{}sha256.\allowbreak{}c:398}). RFC 6234 section 8.5 supplies four that sit on the padding boundaries: \texttt{"abc"}, the 56 octet message whose padding overflows into a second block, 64 octets fed ten times so padding forms a whole block alone, and one million \texttt{'a'} which pushes the length field past \(2^{23}\) bits over 15,625 compressions. The empty message is the fifth, which RFC 8448 section 3 prints as the TLS 1.3 \texttt{Transcript-\allowbreak{}Hash("")}.

\texttt{harness.\allowbreak{}py vectors} runs the normative suites offline against bytes vendored under \texttt{test/\allowbreak{}vectors}, each one's source and digest recorded in \texttt{MANIFEST.\allowbreak{}json}, with every file's hash checked before a vector is read from it.

\tiny
\setlength{\tabcolsep}{2pt}
\begin{longtable}{@{}>{\setlength{\rightskip}{0pt}\setlength{\parfillskip}{0pt}\arraybackslash}p{\dimexpr\textwidth/4-2\tabcolsep\relax}>{\setlength{\rightskip}{0pt}\setlength{\parfillskip}{0pt}\arraybackslash}p{\dimexpr\textwidth/4-2\tabcolsep\relax}>{\setlength{\rightskip}{0pt}\setlength{\parfillskip}{0pt}\arraybackslash}p{\dimexpr\textwidth/4-2\tabcolsep\relax}>{\setlength{\rightskip}{0pt}\setlength{\parfillskip}{0pt}\arraybackslash}p{\dimexpr\textwidth/4-2\tabcolsep\relax}@{}}
\toprule
suite & source & count & what it reaches that nothing else does \\
\midrule
CAVP ShortMsg & NIST & 65 & one shot, 0 to 64 bytes \\
CAVP LongMsg & NIST & 64 & 163 to 6400 bytes \\
CAVP Monte & NIST & 100 × 1000 & chained state: each digest feeds the next \\
CAVP bit ShortMsg & NIST & 513 & 448 do not end on a byte \\
CAVP bit LongMsg & NIST & 512 & 448 do not end on a byte \\
CAVP HMAC & NIST & 225 & 150 carry truncated tags \\
Wycheproof HMAC & C2SP \texttt{b61843a9} & 46 & 20 modified tags that must not reproduce \\
streaming splits & invariant & 3321 cuts & every split of one message must reach one digest \\
differential & second implementation & 424 & padding edges plus random lengths \\
\bottomrule
\end{longtable}
\normalsize

1554 published vectors from two bodies, none of them written here. Monte earns its place twice over: a hash whose state carries wrongly between blocks passes every one shot table and fails on Monte's first checkpoint. Every vector in every published suite hashes its message in a single call. None reaches \texttt{mmgr\_\allowbreak{}sha256\_\allowbreak{}take}, and the split and differential rows are the only thing testing the streaming path.

This is two bodies and not the published corpus of cryptographic oracles. Widening it is open work.

\texttt{mmgr\_\allowbreak{}sha256\_\allowbreak{}bits} is the general entry and takes a \texttt{uint64\_\allowbreak{}t} count, because the standard's length field is 64 bits. Trailing bits sit in the high end of the final byte:

\begin{lstlisting}[breaklines=true,breakatwhitespace=false,basicstyle=\ttfamily\small]
const uint8_t kept   = (uint8_t)(bytes[whole] & (uint8_t)(0xFFu << (8u - spare)));
const uint8_t marker = (uint8_t)(kept | (uint8_t)(0x80u >> spare));
\end{lstlisting}

There is one padding implementation, \texttt{mmgr\_\allowbreak{}sha256\_\allowbreak{}finish\_\allowbreak{}with(running, marker, total\_\allowbreak{}bits, digest)}, and both public entries are that function with its two arguments filled in differently. The byte case passes \texttt{0x80} and \texttt{length * 8}, so no second copy exists to drift from the first.

The interface is \texttt{begin}, \texttt{take}, \texttt{finish} with the one shot built on top (\texttt{test/\allowbreak{}support/\allowbreak{}mmgr\_\allowbreak{}sha256.\allowbreak{}h:59-\allowbreak{}99}). It is streaming because a one shot form cannot run the standard's own vectors: TEST3 is a megabyte and an ATSAMD51 has 192 KB of RAM.

---

\textbf{Author:} dstroy0 (Douglas Quigg) <dquigg123@gmail.com> \textbf{Date:} 2026-09-01 \textbf{Copyright:} (c) 2026 Douglas Quigg (dstroy0). All rights reserved. \textbf{License:} AGPL-3.0-or-later OR LicenseRef-Commercial OR LicenseRef-Educational. Every use falls under AGPL-3.0-or-later unless you hold explicit permission, which is either a negotiated commercial licensing contract or an educator's license issued to you personally.

